\documentclass[11pt]{article}
% Shared preamble for all solution files.
\usepackage[margin=1.1in]{geometry}
\usepackage{enumitem}
\usepackage{tikz}
\usepackage[hidelinks]{hyperref}
\usepackage{amsmath,amssymb,amsthm}
\newcommand{\fa}{\mathfrak a}
\newcommand{\fb}{\mathfrak b}
\newcommand{\fc}{\mathfrak c}
\newcommand{\fp}{\mathfrak p}
\newcommand{\fq}{\mathfrak q}
\newcommand{\fm}{\mathfrak m}
\newcommand{\fn}{\mathfrak n}
\newcommand{\fr}{\mathfrak r}
\newcommand{\Nil}{\mathfrak N}
\newcommand{\Jac}{\mathfrak R}
\newcommand{\rad}{r}
\newcommand{\Spec}{\operatorname{Spec}}
\newcommand{\Max}{\operatorname{Max}}
\newcommand{\Ann}{\operatorname{Ann}}
\newcommand{\Ker}{\operatorname{Ker}}
\newcommand{\Z}{\mathbf Z}
\newcommand{\Q}{\mathbf Q}
\newcommand{\R}{\mathbf R}
\newcommand{\C}{\mathbf C}
\newcommand{\ci}{\mathrm{c}}
\newcommand{\ex}{\mathrm{e}}

\setlength{\parindent}{0pt}
\setlength{\parskip}{0.55em}
% \exercise{label}{statement}  -- followed by the solution text
\newcommand{\exercise}[2]{\par\bigskip\noindent{\large\textbf{Exercise #1}}\par\smallskip
  \begin{quote}\small\textbf{Problem.}\ #2\end{quote}\par}
\newcommand{\pf}{\noindent\textbf{Solution.}\ }

\title{Liu, Chapter 4: Some local properties\\[2pt]\large Solutions to the exercises}
\author{}
\date{}
\begin{document}
\maketitle
\noindent\textbf{Conventions.} Rings are commutative with $1$. Results of the text are cited by their numbers in Liu (``Prop.~3.25'', ``Thm.~3.12'', \dots), and results of earlier chapters as ``Prop.~2.3.12'', ``Ex.~3.2.4'' etc. Exercise labels are those of the book, numbered within each section of the chapter: for instance Exercise~3.12 is the twelfth exercise at the end of Section 4.3. Each exercise is restated before its solution; sub-items that the book displays as nested lists are written as running prose.

Honest notes. Many printed statements need a correction or an interpretation, and each solution says what is proved: Ex.~1.5 (integrality over $A[T]$), Ex.~1.9(c) (the formula for the restriction to $V$), Ex.~1.15, 1.17, 1.18, 1.19(b) (empty open sets, an integrality convention, and counterexamples in characteristic $p$), Ex.~2.1 (the printed claim fails for $\mathbb A^1_k\to\operatorname{Spec}k$; replacement statements are proved), Ex.~3.1 (needs a one-dimensional local ring), Ex.~3.2, 3.8(b), 3.11 (counterexamples and corrected versions), Ex.~3.19(c) and 3.22(a). Some solutions cite standard facts which are named where used.

\section*{4.1 Normal schemes}

Throughout, rings are commutative with $1$, and (as in Section 4.1.1 of the book) a \emph{normal ring} means a normal integral domain. Results of the text are cited by their numbers in Liu (``Prop.~1.5'', ``Thm.~1.14'', \dots), and results of earlier chapters as ``Prop.~2.4.18'' etc.

Honest notes. In Ex.~1.5 the book says ``$f$ integral over $A$''; this must be ``integral over $A[T]$'' (this is what (b) and (c) use), and that is what is proved. In Ex.~1.9(c) the printed formula for $\varphi|_V$ (``$A[T_1/T_0]\to\mathcal O_X(V)$, $T_1/T_0\mapsto f$'', although $A[T_1/T_0]$ is the coordinate ring of $D_+(T_0)$, not of $D_+(T_1)$) has to be read as $T_0/T_1\mapsto f^{-1}$, the coordinate of $D_+(T_1)$; this is the only reading for which the two morphisms agree on $U\cap V$. In Ex.~1.4 the phrase ``this property is verified at the closed points'' is read as: the local rings at the closed points are reduced (resp.\ integral domains, resp.\ normal domains). In Ex.~1.10 I use, for the quotient $X/G$, only the description recalled at the beginning of that solution. Exercises 1.2, 1.5(e), 1.8 use earlier exercises, and the needed arguments are given.

\exercise{1.1}{Let $A$ be a ring and $B$ an $A$-algebra. We suppose that there exists a faithful $B$-module $M$ (i.e., $bM=0$ for $b\in B$ implies $b=0$) which is finitely generated over $A$. Using the Cayley--Hamilton identity on the matrices $\operatorname{Hom}_A(M,M)$ (i.e., any matrix cancels its characteristic polynomial), show that $B$ is integral over $A$ (this generalizes Lemma~1.11).}
\pf
Let $\varphi:A\to B$ be the structure map; $M$ is an $A$-module through $\varphi$. ``Integral'' means that every $b\in B$ is a root of a monic polynomial whose coefficients lie in $\varphi(A)$. Let $m_1,\dots,m_n$ generate $M$ over $A$. If $n=0$ then $M=0$, hence $B=0$ since $M$ is faithful, and there is nothing to prove.

Fix $b\in B$ and let $u:M\to M$ be multiplication by $b$; it is $A$-linear. Since $bm_i\in M$, we can write $bm_i=\sum_j\lambda_{ij}m_j$ with $\lambda_{ij}\in A$. Let $\varepsilon:A^n\to M$, $e_i\mapsto m_i$, which is surjective, and let $L:A^n\to A^n$ be the $A$-linear map $e_i\mapsto\sum_j\lambda_{ij}e_j$. Then $u\circ\varepsilon=\varepsilon\circ L$, hence $P(u)\circ\varepsilon=\varepsilon\circ P(L)$ for every $P\in A[X]$ (where $A[X]$ acts on $M$ through $\varphi$ and $u$). Let $\chi(X)=X^n+c_{n-1}X^{n-1}+\dots+c_0\in A[X]$ be the characteristic polynomial of $L$. By Cayley--Hamilton $\chi(L)=0$, so $\chi(u)\circ\varepsilon=0$, and since $\varepsilon$ is onto, $\chi(u)=0$. This says
\[ \bigl(b^n+\varphi(c_{n-1})b^{n-1}+\dots+\varphi(c_0)\bigr)\,m=0\quad\text{for all }m\in M. \]
As $M$ is faithful, $b^n+\varphi(c_{n-1})b^{n-1}+\dots+\varphi(c_0)=0$. So $b$ is integral over $A$. \hfill$\square$

\exercise{1.2}{Let $A$ be an integral domain, and let $I$ be a finitely generated ideal of $A$. Show that for any $\phi\in\operatorname{Hom}_A(I,I)$ and for any non-zero $a\in I$, $a^{-1}\phi(a)\in\operatorname{Frac}(A)$ is integral over $A$ (use the Cayley--Hamilton identity).}
\pf
Let $K=\operatorname{Frac}(A)$. Put $c=a^{-1}\phi(a)\in K$. For nonzero $x\in I$ we have, since $x,a\in A$ and $\phi$ is $A$-linear,
\[ a\,\phi(x)=\phi(ax)=x\,\phi(a),\qquad\text{so}\qquad \phi(x)=cx. \]
This also holds for $x=0$. Hence $cI=\phi(I)\subseteq I$. Let $B=A[c]\subseteq K$. Then $I$ is a $B$-module (because $cI\subseteq I$), finitely generated over $A$ by hypothesis, and it is faithful: if $b\in B\subseteq K$ and $bI=0$, then $bx=0$ for the nonzero $x=a\in I$, so $b=0$ in the field $K$. By Exercise~1.1 (with $M=I$), $B$ is integral over $A$. In particular $c$ is integral over $A$. \hfill$\square$

\exercise{1.3}{Let $X$ be a normal Noetherian local scheme of dimension $2$, with closed point $s$. Show that $X\setminus\{s\}$ is a non-affine Dedekind scheme.}
\pf
Write $X=\operatorname{Spec}A$ with $(A,\fm)$ a Noetherian local ring; $s=\fm$, $\dim A=2$. Since $X$ is normal, $A=\mathcal O_{X,s}$ is a normal domain, and $X$ is irreducible. Put $U=X\setminus\{s\}$.

\emph{$U$ is a Dedekind scheme.} $U$ is open in the Noetherian scheme $X$, hence Noetherian, and it is non-empty and irreducible (it contains the generic point). Its local rings are $A_\fp$ with $\fp\neq\fm$, which are localizations of the normal domain $A$, hence normal (Lemma~1.4(b)). So $U$ is a normal locally Noetherian scheme. Let $\fp\in U$, i.e.\ $\fp\neq\fm$, and let $\fp_0\subsetneq\dots\subsetneq\fp_r=\fp$ be a chain of primes. Appending $\fm$ gives a chain of length $r+1\le\dim A=2$, so $\operatorname{ht}\fp\le1$. Thus every chain of irreducible closed subsets of $U$ has length $\le1$, i.e.\ $\dim U\le1$. Moreover, as $\dim A=2$ there is a chain $0\subsetneq\fp_1\subsetneq\fm$, and $0\subsetneq\fp_1$ is a chain in $U$; so $\dim U=1$. By Definition~1.2, $U$ is a Dedekind scheme.

\emph{$U$ is not affine.} Since $\operatorname{codim}(\{s\},X)=\dim\mathcal O_{X,s}=\dim A=2$, Thm.~1.14 shows that the restriction $\mathcal O_X(X)=A\to\mathcal O_X(U)$ is an isomorphism. Suppose $U=\operatorname{Spec}B$ were affine. The open immersion $j:U\to X$ corresponds to a ring homomorphism $A\to B$, and this is precisely the restriction map $\mathcal O_X(X)\to\mathcal O_X(U)$, i.e.\ an isomorphism. A morphism of affine schemes whose ring homomorphism is an isomorphism is an isomorphism; so $j$ would be surjective. But $s\notin j(U)$, a contradiction. \hfill$\square$

\exercise{1.4}{Let $X$ be a connected algebraic variety. Show that $X$ is reduced (resp.\ integral; resp.\ normal) if and only if this property is verified at the closed points of $X$.}
\pf
An algebraic variety is of finite type over a field, hence Noetherian and quasi-compact. ``At a point $x$'' means: $\mathcal O_{X,x}$ is reduced (resp.\ an integral domain, resp.\ a normal domain). A closed point means a point closed in $X$.

\emph{Step 0.} If $x\in X$ is any point, there is a closed point $y\in\overline{\{x\}}$ (Exercise~2.4.8, as used in the proof of Prop.~1.5). Every affine open neighbourhood $V=\operatorname{Spec}R$ of $y$ contains $x$ (an open set containing $y\in\overline{\{x\}}$ contains $x$), and $x$ corresponds to a prime contained in the prime of $y$; hence $\mathcal O_{X,x}$ is a localization of $\mathcal O_{X,y}$. A localization of a reduced ring is reduced, of a domain is a domain, and of a normal domain is a normal domain (Lemma~1.4(b)). So if the property holds at all closed points, it holds at all points.

\emph{Reduced.} By Step 0 and the definition of reducedness (all local rings reduced), $X$ is reduced iff it is reduced at its closed points. (Connectedness is not needed.)

\emph{Integral.} If $X$ is integral, all local rings are domains. Conversely suppose $\mathcal O_{X,y}$ is a domain for every closed $y$. By Step 0 all local rings are domains, in particular reduced, so $X$ is reduced; it remains to see that $X$ is irreducible (reduced and irreducible means integral). Let $Z_1,\dots,Z_r$ be the irreducible components of $X$ (finitely many since $X$ is Noetherian). Suppose $Z_i\cap Z_j\neq\emptyset$ for some $i\neq j$. This is a nonempty closed subset, so it contains a closed point $y$ of $X$ (apply Exercise~2.4.8 to a point of it). Take an affine open $V=\operatorname{Spec}R\ni y$, $y\leftrightarrow\fp$. The irreducible components of $V$ passing through $y$ correspond to the minimal primes of $R$ contained in $\fp$, i.e.\ to the minimal primes of $R_\fp=\mathcal O_{X,y}$, which is a domain; so exactly one irreducible component $W$ of $V$ passes through $y$. Now $Z_i\cap V$ and $Z_j\cap V$ are irreducible closed subsets of $V$ containing $y$, so both lie in $W$. Hence $Z_i=\overline{Z_i\cap V}\subseteq\overline W$ and likewise $Z_j\subseteq\overline W$, with $\overline W$ irreducible; by maximality of components $Z_i=\overline W=Z_j$, a contradiction. Therefore the $Z_i$ are pairwise disjoint closed sets covering $X$; each $Z_i$ is the complement of the union of the others, hence also open. As $X$ is connected, $r=1$ and $X$ is irreducible, hence integral.

\emph{Normal.} If $X$ is normal, it is normal at every point. Conversely, suppose $\mathcal O_{X,y}$ is a normal domain for all closed $y$. By the integral case (here connectedness is used), $X$ is integral, in particular irreducible; by Step 0 it is normal at all points. So $X$ is normal. \hfill$\square$

\exercise{1.5}{Let $A$ be a normal ring. We want to show that $A[T]$ is also normal. Let $f\in\operatorname{Frac}(A[T])$ be integral over $A[T]$ (the book writes ``over $A$''; see the remark at the beginning).
\begin{enumerate}[label=(\alph*),leftmargin=*]
\item Show that $f\in\operatorname{Frac}(A)[T]$.
\item Show that there exists a subring $A_0$ of $A$ such that $A_0$ is finitely generated over $\mathbb Z$, $f\in\operatorname{Frac}(A_0)[T]$, and $f$ is integral over $A_0[T]$.
\item Show that there exists an $a\in A_0\setminus\{0\}$ such that $af^n\in A_0[T]$ for every $n\ge0$.
\item Let us write $f=\sum_{i\le d}\alpha_iT^i$. Show that $a\alpha_0^n\in A_0$ for every $n\ge0$, and that $a^{i+1}\alpha_i^n\in A_0$ for every $i\le d$ and for every $n\ge0$.
\item Show that $A_0[\alpha_i]\subseteq a^{-i-1}A_0$, and that $f\in A[T]$.
\item (A second method when $A$ is Noetherian.) Show that $A[T]$ is normal if $A$ is a discrete valuation ring; then prove the general case using Lemma~1.13.
\end{enumerate}}
\pf
Let $K=\operatorname{Frac}(A)$, so $\operatorname{Frac}(A[T])=K(T)$. Since $A[T]$ is a domain, showing that it is normal means showing that every $f\in K(T)$ integral over $A[T]$ lies in $A[T]$. We may assume $f\neq0$. Fix an equation
\[ f^N+g_{N-1}f^{N-1}+\dots+g_0=0,\qquad g_j\in A[T]. \]

\textbf{(a)} $K[T]$ is a principal ideal domain, hence factorial, hence normal by Lemma~1.4(a), and $\operatorname{Frac}(K[T])=K(T)$. The element $f$ is integral over $A[T]\subseteq K[T]$, hence over $K[T]$, so $f\in K[T]$.

\textbf{(b)} Write $f=\sum_{i\le d}\alpha_iT^i$ with $\alpha_i\in K$, and write each $\alpha_i=b_i/c_i$ with $b_i,c_i\in A$, $c_i\neq0$. Let $A_0$ be the subring of $A$ generated over $\mathbb Z$ by all $b_i,c_i$ and by all the (finitely many) coefficients of $g_0,\dots,g_{N-1}$. Then $A_0$ is a domain finitely generated over $\mathbb Z$, $\alpha_i\in K_0:=\operatorname{Frac}(A_0)\subseteq K$, so $f\in K_0[T]$, and the displayed equation has $g_j\in A_0[T]$, so $f$ is integral over $A_0[T]$. Note that $A_0$ is Noetherian (Hilbert's basis theorem).

\textbf{(c)} By the integral equation, $A_0[T][f]=\sum_{k<N}A_0[T]f^k$ (every $f^n$ with $n\ge N$ is an $A_0[T]$-combination of $1,f,\dots,f^{N-1}$, by induction on $n$). Each $f^k$ ($k<N$) lies in $K_0[T]$, so it has finitely many coefficients in $K_0$. Let $a\in A_0\setminus\{0\}$ be a common denominator of all these coefficients, so $af^k\in A_0[T]$ for $k<N$. For $n\ge N$, $f^n=\sum_{k<N}h_kf^k$ with $h_k\in A_0[T]$, so $af^n\in A_0[T]$ as well.

\textbf{(d)} Define $f_0=f$ and recursively $\alpha_i:=f_i(0)\in K_0$ and $f_{i+1}:=(f_i-\alpha_i)/T\in K_0[T]$. By induction $f_i=\sum_{k\ge0}\alpha_{i+k}T^k$ with $\alpha_j$ the coefficients of $f$ (and $\alpha_j=0$ for $j>d$); the constant term of $f_i^n$ is $\alpha_i^n$. Let $B_0=A_0$ and $B_i=A_0[\alpha_0,\dots,\alpha_{i-1}]$. We prove by induction on $i$ that for all $n\ge0$:
\[ (1_i)\ \ af_i^n\in B_i[T];\qquad (2_i)\ \ a\alpha_i^n\in B_i;\qquad (3_i)\ \ B_i\subseteq a^{-i}A_0. \]
For $i=0$, $(1_0)$ is (c), $(2_0)$ is the constant term of $(1_0)$, and $(3_0)$ is trivial. Suppose $(1_i),(2_i),(3_i)$ hold. Note $(2_i)$ follows from $(1_i)$ by taking constant terms. Since $a\alpha_i^n\in B_i$ for all $n$, we get $B_{i+1}=B_i[\alpha_i]=\sum_nB_i\alpha_i^n\subseteq a^{-1}B_i\subseteq a^{-i-1}A_0$, which is $(3_{i+1})$. Next, $Tf_{i+1}=f_i-\alpha_i$, so
\[ T^nf_{i+1}^n=(f_i-\alpha_i)^n=\sum_{k=0}^n\binom nk(-\alpha_i)^{n-k}f_i^k. \]
Multiply by $a$: by $(1_i)$, $af_i^k\in B_i[T]\subseteq B_{i+1}[T]$, and $\alpha_i\in B_{i+1}$; hence $aT^nf_{i+1}^n\in B_{i+1}[T]$. The coefficients of $af_{i+1}^n$ are those of $aT^nf_{i+1}^n$ shifted by $n$ places, so $af_{i+1}^n\in B_{i+1}[T]$: this is $(1_{i+1})$, and $(2_{i+1})$ follows. This completes the induction.

Now $(2_0)$ gives $a\alpha_0^n\in A_0$, and $(2_i),(3_i)$ give $a\alpha_i^n\in B_i\subseteq a^{-i}A_0$, i.e.\ $a^{i+1}\alpha_i^n\in A_0$ for all $i,n$.

\textbf{(e)} By (d), $\alpha_i^n\in a^{-i-1}A_0$ for all $n\ge0$, so $A_0[\alpha_i]=\sum_nA_0\alpha_i^n\subseteq a^{-i-1}A_0$. Now $A_0$ is Noetherian and $a^{-i-1}A_0\simeq A_0$ as $A_0$-modules, so $A_0[\alpha_i]$ is a finitely generated $A_0$-module. It is a faithful module over the ring $A_0[\alpha_i]$ (it contains $1$), so by Exercise~1.1 (with $B=M=A_0[\alpha_i]$) the ring $A_0[\alpha_i]$ is integral over $A_0$; in particular $\alpha_i$ is integral over $A_0\subseteq A$. As $\alpha_i\in K$ and $A$ is normal, $\alpha_i\in A$. This holds for each $i$, so $f\in A[T]$. Hence $A[T]$ is normal.

\textbf{(f)} \emph{Case $A$ a discrete valuation ring}, with uniformizer $t$ and fraction field $K$. Let $f\in K(T)$ be integral over $A[T]$; by (a) $f\in K[T]$. Let $r\ge0$ be minimal with $g:=t^rf\in A[T]$. Suppose $r>0$; then $g\notin tA[T]$ (otherwise $t^{r-1}f\in A[T]$). Multiplying the integral equation by $t^{rN}$ gives $g^N+\sum_{j<N}t^{r(N-j)}g_jg^j=0$, so $g^N\in tA[T]$. But $A[T]/tA[T]\simeq(A/tA)[T]$ is a domain (as $A/tA$ is a field), so $g\in tA[T]$, a contradiction. Hence $r=0$ and $f\in A[T]$.

\emph{General Noetherian case.} Let $A$ be a normal Noetherian domain. If $\dim A=0$, $A$ is a field and $A[T]$ is a principal ideal domain, hence normal. Let $\dim A\ge1$ and let $f\in K(T)$ be integral over $A[T]$. For each prime $\fp$ of height $1$, $A_\fp$ is a Noetherian normal local domain of dimension $1$ (Lemma~1.4(b)); it is a principal ideal domain by Prop.~1.12, and not a field, so it is a discrete valuation ring. Since $A[T]\subseteq A_\fp[T]$ and $\operatorname{Frac}(A_\fp[T])=K(T)$, the first case shows $f\in A_\fp[T]$. Thus every coefficient of $f$ lies in $\bigcap_{\operatorname{ht}\fp=1}A_\fp=A$ (Lemma~1.13), i.e.\ $f\in A[T]$. \hfill$\square$

\exercise{1.6}{Let $A\to B$ be a faithfully flat ring homomorphism, and let us suppose $B$ is normal.
\begin{enumerate}[label=(\alph*),leftmargin=*]
\item Let $I,J$ be ideals of $A$ such that $IB\subseteq JB$. Show that $I\subseteq J$.
\item Let $f=ab^{-1}\in\operatorname{Frac}(A)$ be integral over $A$. Show that $aB\subseteq bB$. Deduce from this that $f\in A$ and that $A$ is normal.
\end{enumerate}}
\pf
Recall that faithfully flat means: $B$ is flat over $A$, and $M\otimes_AB=0$ implies $M=0$ for every $A$-module $M$.

\textbf{(a)} Let $M=(I+J)/J$. Tensoring the exact sequence $0\to J\to I+J\to M\to0$ with the flat module $B$, and using that for an ideal $\fa$ of $A$ the map $\fa\otimes_AB\to B$ is injective with image $\fa B$, we get
\[ M\otimes_AB\simeq\frac{(I+J)B}{JB}=\frac{IB+JB}{JB}=0 \]
because $IB\subseteq JB$. By faithful flatness $M=0$, i.e.\ $I\subseteq J$.

(In particular, taking $I=\operatorname{Ker}(A\to B)$ and $J=0$, we get $IB=0$, so $I=0$: the map $A\to B$ is injective. Thus $A\subseteq B$ is a domain.)

\textbf{(b)} Let $f=ab^{-1}$ with $a,b\in A$, $b\neq0$, integral over $A$. Since $A\subseteq B$, $f\in\operatorname{Frac}(A)\subseteq\operatorname{Frac}(B)$ is integral over $B$, so $f\in B$ because $B$ is normal. Hence $a=fb\in bB$, i.e.\ $aB\subseteq bB$. Applying (a) to the ideals $aA$ and $bA$ (note $(aA)B=aB$ and $(bA)B=bB$) gives $aA\subseteq bA$, so $a=bc$ with $c\in A$ and $f=c\in A$. Since $A$ is a domain in which every element of $\operatorname{Frac}(A)$ integral over $A$ lies in $A$, $A$ is normal. \hfill$\square$

\exercise{1.7}{Let $k$ be a field. Determine the normalization of $\operatorname{Proj}k[x,y,z]/(x^3-y^2z)$.}
\pf
Let $C=\operatorname{Proj}k[x,y,z]/(x^3-y^2z)$ (any characteristic). \textbf{Answer:} the normalization is
\[ \varphi:\mathbb P^1_k\to C,\qquad [s:u]\mapsto[s^2u:s^3:u^3], \]
an isomorphism away from the cusp $[0:0:1]$, with $\varphi^{-1}([0:0:1])=\{[0:1]\}$.

\emph{The charts of $C$.} Let $U_y=D_+(y)$ and $U_z=D_+(z)$. If $y=z=0$ at a point then $x^3=0$, so $U_y\cup U_z=C$. We have
\[ \mathcal O(U_y)=k[X,Z]/(X^3-Z)\simeq k[X]\ \ (X=x/y,\ Z=z/y),\qquad \mathcal O(U_z)=k[a,b]/(a^3-b^2)\ \ (a=x/z,\ b=y/z). \]
The ring $k[a,b]/(a^3-b^2)$ is isomorphic to $k[v^2,v^3]\subset k[v]$ via $a\mapsto v^2$, $b\mapsto v^3$: this map is onto $k[v^2,v^3]$, and it is injective because every class has a representative $p_0(b)+p_1(b)a+p_2(b)a^2$ (reduce $a^3=b^2$), whose image $p_0(v^3)+p_1(v^3)v^2+p_2(v^3)v^4$ has the three summands supported on exponents $\equiv0,2,1\pmod 3$ respectively, so it vanishes only if all $p_i=0$. Thus $U_y\simeq\mathbb A^1_k$ and $U_z\simeq\operatorname{Spec}k[v^2,v^3]$ are integral, and $U_y\cap U_z=D(X)\neq\emptyset$; two irreducible opens that meet generate an irreducible space, so $C$ is irreducible and reduced, i.e.\ an integral scheme.

\emph{The morphism $\varphi$.} Let $V_1=\{s\ne0\}=\operatorname{Spec}k[w]$ with $w=u/s$, and $V_2=\{u\neq0\}=\operatorname{Spec}k[v]$ with $v=s/u$. Define $\varphi_1:V_1\to U_y$ by $X\mapsto w$ (then $Z=X^3\mapsto w^3$, i.e.\ $[x:y:z]=[w:1:w^3]$), and $\varphi_2:V_2\to U_z$ by $a\mapsto v^2$, $b\mapsto v^3$. On $V_1\cap V_2=\operatorname{Spec}k[w^{\pm1}]$, $\varphi_1$ lands in $U_y\cap U_z=D(X)$, where $a=X/Z=X^{-2}$ and $b=1/Z=X^{-3}$; so $\varphi_1$ sends $a\mapsto w^{-2}=v^2$ and $b\mapsto w^{-3}=v^3$, which is $\varphi_2$. So $\varphi_1,\varphi_2$ glue to a morphism $\varphi:\mathbb P^1_k\to C$, and it is given by the stated formula. Moreover $\varphi^{-1}(U_y)=V_1$ (a point of $V_2$ goes to $U_y$ iff $b=v^3\neq0$, i.e.\ it lies in $V_1\cap V_2$) and $\varphi^{-1}(U_z)=V_2$ (a point of $V_1$ goes to $U_z$ iff $z/y=w^3\neq0$, i.e.\ it lies in $V_2$).

\emph{$\varphi$ is the normalization.} $\mathbb P^1_k$ is normal (Example~1.6) and $\varphi$ is dominant (it is an isomorphism $V_1\to U_y$, since $k[X]\to k[w]$, $X\mapsto w$, is an isomorphism). Let $\pi:C'\to C$ be the normalization (Prop.~1.22). By Def.~1.19, $\varphi=\pi\circ g$ for a morphism $g:\mathbb P^1_k\to C'$. Over $U_y$: $\mathcal O(U_y)=k[X]$ is normal, so by Lemma~1.21 the normalization of $U_y$ is the identity, and $\varphi^{-1}(U_y)=V_1\to U_y$ is an isomorphism. Over $U_z$: $A=k[v^2,v^3]$ has fraction field $k(v)$ (as $v=v^3/v^2$), and its integral closure in $k(v)$ is $k[v]$ (the ring $k[v]$ is factorial, hence normal, and it is integral over $A$ because $v$ is a root of the monic polynomial $T^2-v^2\in A[T]$). By Lemma~1.21, $\operatorname{Spec}k[v]=V_2\to U_z$ is the normalization of $U_z$. Since the restriction of a normalization to an open subset of the target is a normalization (remark after Def.~1.19), $\pi^{-1}(U_y)\to U_y$ and $\pi^{-1}(U_z)\to U_z$ are the normalizations of $U_y,U_z$, and $g$ restricts to morphisms $\varphi^{-1}(U_i)\to\pi^{-1}(U_i)$ between two normalizations of $U_i$ over $U_i$, which are isomorphisms by uniqueness of normalization. As $U_y,U_z$ cover $C$, $g$ is an isomorphism, so $\varphi$ is the normalization. \hfill$\square$

\exercise{1.8}{Let $\mathcal O_K$ be a Dedekind domain, and let $a\in\mathcal O_K$ be non-zero. Show that the scheme $\operatorname{Proj}\mathcal O_K[x,y,z]/(xy-az^2)$ is normal.}
\pf
Write $A=\mathcal O_K$, $F=\operatorname{Frac}(A)$, and $X=\operatorname{Proj}A[x,y,z]/(xy-az^2)$. We use that $A[T]$ is normal for $A$ a normal Noetherian domain: this was proved in Exercise~1.5(f) (for a DVR by the $t$-adic argument, in general through $A=\bigcap_{\operatorname{ht}\fp=1}A_\fp$ of Lemma~1.13; if $A$ is a field it is clear). Then $A[T]$ and its localizations (Lemma~1.4(b)) are normal.

\emph{The three charts.} $X=D_+(x)\cup D_+(y)\cup D_+(z)$. We have
\[ D_+(x)=\operatorname{Spec}A[Y,W]/(Y-aW^2)\simeq\operatorname{Spec}A[W],\quad D_+(y)=\operatorname{Spec}A[X',W']/(X'-aW'^2)\simeq\operatorname{Spec}A[W'], \]
with $Y=y/x$, $W=z/x$, $X'=x/y$, $W'=z/y$, and
\[ D_+(z)=\operatorname{Spec}B,\qquad B=A[\xi,\eta]/(\xi\eta-a),\quad \xi=x/z,\ \eta=y/z. \]
Two of the charts are polynomial rings over $A$, hence normal domains. It remains to study $B$.

\emph{Structure of $B$.} Since $\xi\eta=a\in A$, $B$ is spanned over $A$ by $\xi^i$ ($i\ge0$) and $\eta^j$ ($j\ge1$). Consider $\theta:B\to F[\xi,\xi^{-1}]$ (Laurent polynomial ring in a new variable $\xi$) with $\xi\mapsto\xi$, $\eta\mapsto a\xi^{-1}$; it is a well-defined ring homomorphism. The images $\xi^i$ ($i\ge0$) and $a^j\xi^{-j}$ ($j\ge1$) of the spanning elements are linearly independent over $A$ (as $a\ne0$), so these elements of $B$ form an $A$-basis, and $\theta$ is injective with image
\[ B_0=\Bigl\{\sum_ic_i\xi^i\ :\ c_i\in A\ (i\ge0),\ c_{-j}\in a^jA\ (j\ge1)\Bigr\}. \]
Hence $B$ is a domain with $\operatorname{Frac}(B)=F(\xi)$. Under $\theta$:
\begin{itemize}
\item $B[1/\xi]$ becomes $A[\xi,\xi^{-1}]=\{\sum c_i\xi^i:\ c_i\in A\ \forall i\}$ (because $\eta=a\xi^{-1}$ lies in it);
\item $B[1/\eta]=A[\eta,\eta^{-1}]$ (because $\xi=a\eta^{-1}$) becomes $\{\sum c_i\xi^i:\ c_{-j}\in a^jA\ (j\ge0),\ c_i\in a^{-i}A\ (i>0)\}$, since $\eta^j=a^j\xi^{-j}$ and, for $j>0$, $\eta^{-j}=a^{-j}\xi^j$.
\end{itemize}
The intersection of these two rings inside $F[\xi,\xi^{-1}]$ is the set of $\sum c_i\xi^i$ with $c_i\in A\cap a^{-i}A=A$ for $i>0$ and $c_{-j}\in A\cap a^jA=a^jA$ for $j\ge0$, that is, exactly $B_0$. So
\[ B=B[1/\xi]\cap B[1/\eta]. \]
Both $B[1/\xi]\simeq A[\xi^{\pm1}]$ and $B[1/\eta]\simeq A[\eta^{\pm1}]$ are normal (localizations of polynomial rings over $A$). \emph{$B$ is normal}: let $g\in\operatorname{Frac}(B)$ be integral over $B$. It is integral over the normal domains $B[1/\xi]$ and $B[1/\eta]$ (both have fraction field $\operatorname{Frac}(B)$), hence lies in both, hence in $B$.

\emph{$X$ is normal.} Every point of $X$ lies in one of the three affine charts, each the spectrum of a normal domain, so its local ring is a localization of a normal domain, hence normal. It remains to see that $X$ is irreducible. The three charts are integral. The open set $D_+(z)\cap D_+(x)=\operatorname{Spec}B[1/\xi]$ is nonempty, and so is $D_+(z)\cap D_+(y)$. A nonempty open subset of an irreducible space contains its generic point, which is also the generic point of the open subset; so the generic points of $D_+(x)$, $D_+(z)$ coincide, and likewise for $D_+(y)$ and $D_+(z)$. Call it $\zeta$. Then each chart is contained in $\overline{\{\zeta\}}$, so $X=\overline{\{\zeta\}}$ is irreducible. Thus $X$ is normal. \hfill$\square$

\exercise{1.9}{Let $X=\operatorname{Spec}A$ be an affine Dedekind scheme. We are going to show that every open subscheme $U$ of $X$ is affine. We may assume that $U$ is the complement of a closed point $x_0$.
\begin{enumerate}[label=(\alph*),leftmargin=*]
\item Let $\fm_0$ be the maximal ideal of $A$ corresponding to $x_0$. Let $t\in A$ be such that $t\mathcal O_{X,x_0}=\fm_0\mathcal O_{X,x_0}$ and let $V(t)=\{x_0,x_1,\dots,x_n\}$. Show that for every $1\le i\le n$, there exists $t_i\in A$ such that $t_i\notin\fm_0$ and that $t_i$ generates the maximal ideal of $\mathcal O_{X,x_i}$.
\item Show that there exists an $f\in\mathcal O_X(U)\setminus A$ (take $f=t^{-1}\prod_{1\le i\le n}t_i^{d_i}$ with large $d_i$).
\item Let $Z=V(f)\subset U$ and $V=X\setminus Z$. Show that there exists a (unique) morphism $\varphi:X\to\mathbb P^1_X=\operatorname{Proj}A[T_0,T_1]$ such that $\varphi|_U$ is the morphism of $A$-schemes $U\to D_+(T_0)$ corresponding to $A[T_1/T_0]\to\mathcal O_X(U)$, $T_1/T_0\mapsto f$; and $\varphi|_V$ is the morphism of $A$-schemes $V\to D_+(T_1)$ corresponding to $A[T_0/T_1]\to\mathcal O_X(V)$, $T_0/T_1\mapsto f^{-1}$ (the printed formula for $V$ is corrected, see the beginning).
\item Show that $\varphi$ is a section of the canonical morphism $\mathbb P^1_X\to X$, and that $U=\varphi^{-1}(D_+(T_0))$. Deduce from this that $U$ is affine, and that $\mathcal O_X(U)=A[f]$.
\end{enumerate}}
\pf
\emph{Preliminaries.} $A$ is a Noetherian normal domain of dimension $\le1$; let $F=\operatorname{Frac}(A)$. If $\dim A=0$, $X$ is a point and every open subset ($\emptyset$ or $X$) is affine; so let $\dim A=1$. Every nonzero prime is then maximal. For a maximal ideal $\fm$ the ring $A_\fm$ is a local principal ideal domain (Prop.~1.12) which is not a field, i.e.\ a discrete valuation ring; let $v_\fm$ be its valuation on $F^*$, so $A_\fm=\{g\in F:v_\fm(g)\ge0\}$.

\emph{Reduction to the complement of one closed point.} Let $U\subseteq X$ be open and nonempty; it contains the generic point. $X\setminus U=V(I)$ with $I\neq0$, and $A/I$ is a Noetherian ring of dimension $0$, so $X\setminus U=\{x_1',\dots,x_r'\}$ is a finite set of closed points. Put $U_0=X$ and $U_k=U_{k-1}\setminus\{x_k'\}$. If $U_{k-1}$ is affine, it is an affine Dedekind scheme (an open subscheme of a Dedekind scheme is Dedekind, Rem.~1.3), and $U_k$ is the complement of a closed point in it. So by induction it suffices to treat $U=X\setminus\{x_0\}$ with $X$ an affine Dedekind scheme and $x_0$ a closed point, as in the statement. Also recall that for the integral scheme $X$ one has $\mathcal O_X(W)=\bigcap_{x\in W}\mathcal O_{X,x}\subseteq F$ for every open $W$ (compare Prop.~2.4.18 as used in the proof of Prop.~1.5).

\textbf{(a)} The existence of $t$: $\mathcal O_{X,x_0}=A_{\fm_0}$ is a discrete valuation ring, so $\fm_0A_{\fm_0}=(\tau)$ is principal, with $\tau=t/s$, $t\in\fm_0$, $s\in A\setminus\fm_0$; as $s$ is a unit in $A_{\fm_0}$, $tA_{\fm_0}=\fm_0A_{\fm_0}$. Also $t\neq0$, so $V(t)$ consists of closed points (it misses the generic point) and is finite ($A/tA$ is Noetherian of dimension $0$). Let $\fm_i$ correspond to $x_i$ ($1\le i\le n$); $\fm_i\neq\fm_0$.

Fix $i$. As $A_{\fm_i}$ is a discrete valuation ring, $\fm_iA_{\fm_i}=cA_{\fm_i}$ for some $c\in A$ (as above); then $c\in\fm_iA_{\fm_i}\cap A=\fm_i$. The ideals $\fm_0$ and $\fm_i^2$ are comaximal, so by the Chinese remainder theorem there is $t_i\in A$ with $t_i\equiv1\pmod{\fm_0}$ and $t_i\equiv c\pmod{\fm_i^2}$. Then $t_i\notin\fm_0$. Write $t_i=c+w$ with $w\in\fm_i^2$; in $A_{\fm_i}$ we have $w\in c^2A_{\fm_i}$, so $t_i=c(1+c\,w')$ with $w'\in A_{\fm_i}$, and $1+cw'$ is a unit of the local ring $A_{\fm_i}$ (as $c\in\fm_iA_{\fm_i}$). Hence $t_iA_{\fm_i}=cA_{\fm_i}=\fm_iA_{\fm_i}$: $t_i$ generates the maximal ideal of $\mathcal O_{X,x_i}$.

\textbf{(b)} Let $e_i=v_{\fm_i}(t)\ge1$ (since $t\in\fm_i$), choose integers $d_i\ge e_i$, and set $f=t^{-1}\prod_{i=1}^nt_i^{d_i}\in F$. We check membership in each local ring $\mathcal O_{X,x}$, $x\in U$.
\begin{itemize}
\item $x$ the generic point: $\mathcal O_{X,x}=F$. Fine.
\item $x=\fm$ a maximal ideal with $t\notin\fm$, i.e.\ $x\notin V(t)$: $t$ is a unit of $A_\fm$ and $t_i\in A$, so $f\in A_\fm$.
\item $x=x_i$, $1\le i\le n$: $v_{\fm_i}(t_i)=1$ and $v_{\fm_i}(t_j)\ge0$ for all $j$ (as $t_j\in A$), so $v_{\fm_i}(f)\ge-e_i+d_i\ge0$.
\end{itemize}
Hence $f\in\bigcap_{x\in U}\mathcal O_{X,x}=\mathcal O_X(U)$. But at $x_0$ we have $v_{\fm_0}(t)=1$ and $v_{\fm_0}(t_i)=0$ (as $t_i\notin\fm_0$), so $v_{\fm_0}(f)=-1<0$, i.e.\ $f\notin A_{\fm_0}\supseteq A$. So $f\in\mathcal O_X(U)\setminus A$. The same computation shows $f^{-1}\in\fm_0A_{\fm_0}$, which we use next.

\textbf{(c)} $Z=V(f)=\{x\in U:f\in\fm_xA_x\}$ is closed in $U$ and does not contain $x_0$; $V=X\setminus Z=\{x_0\}\cup\{x\in U:f\ \text{is a unit in}\ \mathcal O_{X,x}\}$ is open in $X$, and $U\cup V=X$ since $x_0\in V$. On $U\cap V=U\setminus Z$ the function $f$ is invertible. Moreover $f^{-1}\in\mathcal O_X(V)$: for $x\in U\cap V$, $f^{-1}\in\mathcal O_{X,x}$; for $x=x_0$, $v_{\fm_0}(f^{-1})=1\ge0$. We have $D_+(T_0)=\operatorname{Spec}A[S]$ with $S=T_1/T_0$ and $D_+(T_1)=\operatorname{Spec}A[S']$ with $S'=T_0/T_1=S^{-1}$, and a morphism from a scheme $W$ to an affine scheme $\operatorname{Spec}R$ corresponds to a ring homomorphism $R\to\mathcal O(W)$. Let $\varphi_U:U\to D_+(T_0)$ correspond to $S\mapsto f\in\mathcal O_X(U)$ and $\varphi_V:V\to D_+(T_1)$ to $S'\mapsto f^{-1}\in\mathcal O_X(V)$. Both are morphisms of $A$-schemes. On $U\cap V$ both land in $D_+(T_0)\cap D_+(T_1)=\operatorname{Spec}A[S,S^{-1}]$ (since $f$ and $f^{-1}$ are units there) and both correspond to the ring homomorphism $A[S,S^{-1}]\to\mathcal O_X(U\cap V)$, $S\mapsto f$. Hence they agree on $U\cap V$ and glue to a morphism $\varphi:X\to\mathbb P^1_X$ with the required restrictions. It is unique because $U\cup V=X$ and its restrictions to $U$ and $V$ are prescribed.

\textbf{(d)} Let $\pi:\mathbb P^1_X\to X$ be the structure morphism. As $\varphi_U,\varphi_V$ are $A$-morphisms, $\pi\circ\varphi$ restricts to the inclusions on $U$ and on $V$; since $U\cup V=X$, $\pi\circ\varphi=\mathrm{id}_X$: $\varphi$ is a section.

Clearly $U\subseteq\varphi^{-1}(D_+(T_0))$. The only point of $X$ not in $U$ is $x_0$, and $\varphi(x_0)\in D_+(T_1)$ is the point where the coordinate $S'$ vanishes: indeed $S'$ maps to $f^{-1}\in\fm_0A_{\fm_0}$, whose value in $k(x_0)$ is $0$. A point of $D_+(T_0)\cap D_+(T_1)=D(S')\subset D_+(T_1)$ has $S'\neq0$ in its residue field, so $\varphi(x_0)\notin D_+(T_0)$. Hence $\varphi^{-1}(D_+(T_0))=U$.

\emph{$U$ is affine.} A section of a separated morphism is a closed immersion (it is the base change of the diagonal $\mathbb P^1_X\to\mathbb P^1_X\times_X\mathbb P^1_X$ along $(\mathrm{id},\varphi\circ\pi)$), and $\pi$ is separated. So $\varphi$ is a closed immersion, and therefore so is its restriction $\varphi_U:U=\varphi^{-1}(D_+(T_0))\to D_+(T_0)=\operatorname{Spec}A[S]$. A closed subscheme of an affine scheme is affine, so $U$ is affine; the closed immersion $\operatorname{Spec}\mathcal O_X(U)\to\operatorname{Spec}A[S]$ corresponds to a \emph{surjective} homomorphism $A[S]\to\mathcal O_X(U)$, which is $S\mapsto f$. Hence $\mathcal O_X(U)=A[f]$. \hfill$\square$

\exercise{1.10}{Let $X$ be a quasi-projective scheme over a locally Noetherian scheme $S$. Let $G$ be a finite group acting on $X$, so that the quotient scheme $X/G$ exists (Exercise~3.3.23). Let us moreover suppose that $X$ is normal.
\begin{enumerate}[label=(\alph*),leftmargin=*]
\item Show that $Y:=X/G$ is normal.
\item Show that the first projection $X\times_YX\to X$ induces an isomorphism between any irreducible component $Z$ of $X\times_YX$, endowed with the reduced subscheme structure, and the scheme $X$. (Show that $Z\to X$ is finite and birational.)
\end{enumerate}}
\pf
\emph{What is used about $X/G$.} Let $\pi:X\to Y=X/G$ be the quotient morphism. By Exercise~3.3.23 (and its proof, which goes through Exercise~2.3.21(c)), $X$ is covered by $G$-stable affine open sets $W_i=\operatorname{Spec}B_i$; by Exercise~2.3.21(a),(b), $Y$ is obtained by gluing the affine schemes $\operatorname{Spec}(B_i^{G})$, which are open in $Y$ and cover it, and on each of them $\mathcal O_Y=(\pi_*\mathcal O_X)^G$ (this is also the quotient as a ringed space, Exercise~2.2.14). Since $V\mapsto\mathcal O_X(\pi^{-1}(V))^G$ is a sheaf on $Y$, we get $\mathcal O_Y(V)=\mathcal O_X(\pi^{-1}(V))^G$ for every open $V\subseteq Y$. Moreover $\pi\circ\sigma=\pi$ for $\sigma\in G$, $\pi$ is surjective (on each $W_i$ it is $\operatorname{Spec}B_i\to\operatorname{Spec}B_i^G$, surjective because $B_i$ is integral over $B_i^G$) and finite (Exercise~3.3.23, as $S$ is locally Noetherian). Consequently, for every affine open $V\subseteq Y$ the open set $\pi^{-1}(V)$ is affine (finiteness) and $G$-stable, and with $C=\mathcal O_Y(V)$, $B=\mathcal O_X(\pi^{-1}(V))$ we have $\pi^{-1}(V)=\operatorname{Spec}B$, $V=\operatorname{Spec}C$ and $C=B^\Gamma$, where $\Gamma\subseteq\operatorname{Aut}(B)$ is the finite group of ring automorphisms $\sigma^\sharp$ induced by $\sigma\in G$.

\textbf{(a)} $X$ is irreducible, and $\pi$ is surjective and continuous, so $Y$ is irreducible. Let $V=\operatorname{Spec}C\subseteq Y$ be an affine open, $\pi^{-1}(V)=\operatorname{Spec}B$, $C=B^\Gamma$. By Prop.~1.5 ((i)$\Rightarrow$(ii)), $B$ is a normal domain, hence $C$, a subring of $B$, is a domain. Let $\alpha\in\operatorname{Frac}(C)\subseteq\operatorname{Frac}(B)$ be integral over $C$. Then $\alpha$ is integral over $B$, so $\alpha\in B$ since $B$ is normal. The automorphisms in $\Gamma$ extend to $\operatorname{Frac}(B)$ and fix $C$, hence fix $\operatorname{Frac}(C)\ni\alpha$; so $\alpha\in B^\Gamma=C$. Thus $C$ is normal. Every local ring $\mathcal O_{Y,y}$ is a localization of such a $C$, hence normal (Lemma~1.4(b)). So $Y$ is irreducible and normal at all its points: $Y$ is normal.

\textbf{(b)} We first prove an algebraic lemma.

\emph{Lemma.} Let $B$ be a ring, $\Gamma\subseteq\operatorname{Aut}(B)$ a finite group, $C=B^\Gamma$, $\Omega$ a field and $j_1,j_2:B\to\Omega$ ring homomorphisms with $j_1|_C=j_2|_C$. Then there is $\tau\in\Gamma$ with $j_2=j_1\circ\tau$.

\emph{Proof.} Let $\Gamma$ act on $B[X]$ coefficientwise, so $B[X]^\Gamma=C[X]$. Extend $j_1,j_2$ to $\tilde j_1,\tilde j_2:B[X]\to\Omega[X]$ by $X\mapsto X$; they agree on $C[X]$. For $\beta\in B[X]$, the polynomial $P_\beta(T)=\prod_{\tau\in\Gamma}(T-\tau\beta)\in B[X][T]$ has $\Gamma$-invariant coefficients, hence coefficients in $C[X]$. Applying $\tilde j_1$ or $\tilde j_2$ to its coefficients gives the same polynomial $\prod_\tau(T-\tilde j_1(\tau\beta))=\prod_\tau(T-\tilde j_2(\tau\beta))$ in $\Omega[X][T]$. Its right-hand form has the root $\tilde j_2(\beta)$ (the factor $\tau=1$), so in the domain $\Omega[X]$ we get $\tilde j_2(\beta)=\tilde j_1(\tau\beta)$ for some $\tau=\tau(\beta)\in\Gamma$. Let $E_\tau=\{b\in B:j_2(b)=j_1(\tau b)\}$. Suppose no $E_\tau$ equals $B$: pick $b_\tau\in B\setminus E_\tau$ for each $\tau$, pairwise distinct integers $n_\tau\ge0$, and $\beta=\sum_\tau b_\tau X^{n_\tau}$. Then for some $\tau$, $\tilde j_2(\beta)=\tilde j_1(\tau\beta)$; comparing the coefficients of $X^{n_\tau}$ gives $j_2(b_\tau)=j_1(\tau b_\tau)$, i.e.\ $b_\tau\in E_\tau$, a contradiction. Hence $E_\tau=B$ for some $\tau$. $\square$

\emph{Local description of $X\times_YX$.} Let $V=\operatorname{Spec}C$ be an affine open of $Y$, $\pi^{-1}(V)=\operatorname{Spec}B$ as above. Then $W_V:=(\pi\circ p_1)^{-1}(V)=\pi^{-1}(V)\times_V\pi^{-1}(V)=\operatorname{Spec}(B\otimes_CB)$, and these $W_V$ cover $X\times_YX$. For $\sigma\in G$, let $\gamma_\sigma=(\mathrm{id}_X,\sigma):X\to X\times_YX$ (a morphism since $\pi\sigma=\pi$); then $p_1\gamma_\sigma=\mathrm{id}_X$ and $\gamma_\sigma^{-1}(W_V)=\pi^{-1}(V)$. Over $V$, $\gamma_\sigma$ corresponds to $\psi_\sigma:B\otimes_CB\to B$, $b\otimes b'\mapsto b\,\sigma^\sharp(b')$, which is surjective ($b\otimes1\mapsto b$). So $\gamma_\sigma$ is a closed immersion over each $W_V$, hence a closed immersion, and $\Gamma_\sigma:=\gamma_\sigma(X)$ is a closed irreducible subset of $X\times_YX$, with $\gamma_\sigma\!:X\to\Gamma_\sigma$ a homeomorphism.

\emph{Claim: $X\times_YX=\bigcup_{\sigma\in G}\Gamma_\sigma$ as sets.} It suffices to check this inside each $\operatorname{Spec}(B\otimes_CB)$. Let $Q$ be a prime of $B\otimes_CB$, $\Omega=\operatorname{Frac}((B\otimes_CB)/Q)$, and let $j_1(b)=\overline{b\otimes1}$, $j_2(b')=\overline{1\otimes b'}$ in $\Omega$. They agree on $C$ since $c\otimes1=1\otimes c$ in $B\otimes_CB$. By the Lemma, $j_2=j_1\circ\tau$ with $\tau=\sigma^\sharp$ for some $\sigma\in G$, i.e.\ $1\otimes b'\equiv\sigma^\sharp(b')\otimes1\pmod Q$ for all $b'$. Then for $x=\sum_ib_i\otimes b_i'\in\operatorname{Ker}\psi_\sigma$, $x\equiv(\sum_ib_i\sigma^\sharp(b_i'))\otimes1=\psi_\sigma(x)\otimes 1=0\pmod Q$. So $Q\supseteq\operatorname{Ker}\psi_\sigma$, i.e.\ $Q\in\Gamma_\sigma\cap W_V$. This proves the claim.

\emph{Conclusion.} The finitely many irreducible closed sets $\Gamma_\sigma$ cover $X\times_YX$. An irreducible closed subset lying in a finite union of closed sets lies in one of them, so the irreducible components of $X\times_YX$ are the maximal ones among the $\Gamma_\sigma$. Let $Z$ be such a component, $Z=\Gamma_\sigma$, with its reduced structure. Since $X$ is reduced (it is integral, being normal) and $\gamma_\sigma:X\to X\times_YX$ is a closed immersion with image $Z$, $\gamma_\sigma$ factors through an isomorphism $X\xrightarrow{\ \sim\ }Z$ (the closed subscheme of $X\times_YX$ defined by $\gamma_\sigma$ is reduced with underlying set $Z$, hence equals $Z$ with the reduced structure). Since $p_1\circ\gamma_\sigma=\mathrm{id}_X$, the restriction $p_1|_Z:Z\to X$ is the inverse of this isomorphism, hence an isomorphism. In particular $Z\to X$ is finite and birational. \hfill$\square$

Remark on (b): for the reduced structure only the reducedness of $X$ was used; normality of $X$ enters in (a) only. The route via ``finite and birational onto a normal scheme'' suggested in the book is replaced by the direct identification of the components with the graphs $\Gamma_\sigma$.

\exercise{1.11}{Let $X$ be an integral algebraic variety over a field $k$. Let $k'=K(X)\cap\bar k$.
\begin{enumerate}[label=(\alph*),leftmargin=*]
\item Let $X'$ be the normalization of $X$. Show that $X'$ is an algebraic variety over $k'$.
\item Show that for any point $x\in X'$, $k'$ is a subfield of $k(x)$. Deduce from this that $k'$ is finite over $k$.
\end{enumerate}}
\pf
\textbf{Reading of $k'$.} Here $\bar k$ is an algebraic closure of $k$, and $K(X)\cap\bar k$ is the set of elements of $K(X)$ that are algebraic over $k$ (any common overfield of $K(X)$ and $\bar k$ gives the same set). So $k'$ is the algebraic closure of $k$ in $K(X)$; it is a subfield of $K(X)$ containing $k$. An \emph{algebraic variety} is a separated scheme of finite type over a field.

\textbf{Step 1: $k'$ lies in every local ring of $X'$.} Let $\pi:X'\to X$ be the normalization. It is birational (Prop.~1.22), so $K(X')=K(X)$, and it is finite (Cor.~1.30). Let $x\in X'$. The ring $\mathcal O_{X',x}$ is a normal domain with fraction field $K(X)$ and contains $k$. An element $\alpha\in k'$ is algebraic over $k$, hence integral over $k$ (divide its minimal polynomial by the leading coefficient; $k$ is a field), hence integral over $\mathcal O_{X',x}$. It lies in the fraction field, so $\alpha\in\mathcal O_{X',x}$. Thus $k'\subseteq\mathcal O_{X',x}$ for all $x\in X'$. For a non-empty open $V\subseteq X'$ we have $\mathcal O_{X'}(V)=\bigcap_{x\in V}\mathcal O_{X',x}$ inside $K(X')$, so $k'\subseteq\mathcal O_{X'}(V)$ compatibly with restrictions.

\textbf{(a).} By Step 1, $\mathcal O_{X'}$ is a sheaf of $k'$-algebras, so $X'$ is a $k'$-scheme (Cor.~2.3.26), and the $k$-structure is the one induced by $k\subseteq k'$. The scheme $X'$ is of finite type over $k$ (finite over $X$), so it is covered by finitely many affine opens $\operatorname{Spec}B_j$ with $B_j$ finitely generated $k$-algebras. They are finitely generated $k'$-algebras (add $k'$ to the generators, using $k\subseteq k'$). Hence $X'\to\operatorname{Spec}k'$ is of finite type. The morphism $X'\to X$ is finite, hence affine, hence separated, and $X$ is separated over $k$, so $X'$ is separated over $k$. If $g\circ f$ is separated then so is $f$; applying this to $X'\to\operatorname{Spec}k'\to\operatorname{Spec}k$ shows that $X'\to\operatorname{Spec}k'$ is separated. So $X'$ is an algebraic variety over $k'$ (it is even integral and normal).

\textbf{(b).} Let $x\in X'$. By Step 1 we have a ring homomorphism $k'\to\mathcal O_{X',x}\to k(x)$. It is injective because $k'$ is a field, and it is $k$-linear. So $k'$ is a subfield of $k(x)$, containing $k$.

Now choose a non-empty affine open $\operatorname{Spec}B\subseteq X'$ ($B$ a finitely generated $k$-algebra, $B\neq0$) and a maximal ideal $\mathfrak m$ of $B$, with corresponding point $x$. Then $k(x)=B/\mathfrak m$ is a finite extension of $k$ (Cor.~2.1.12). Hence $k\subseteq k'\subseteq k(x)$ and $[k':k]\le[k(x):k]<\infty$. \hfill$\square$

\exercise{1.12}{Let $X$ be a normal and proper algebraic variety over a field $k$. Show that $\mathcal O_X(X)=K(X)\cap\bar k$.}
\pf
Since $X$ is normal, it is irreducible with normal (in particular reduced) local rings, hence integral. As in Ex.~1.11, $K(X)\cap\bar k$ means the set $k'$ of elements of $K(X)$ algebraic over $k$.

\emph{$k'\subseteq\mathcal O_X(X)$.} Let $\alpha\in k'$. It is integral over $k$. For each $x\in X$ the ring $\mathcal O_{X,x}$ is a normal domain with fraction field $K(X)$ containing $k$, so $\alpha\in\mathcal O_{X,x}$. As $X$ is integral, $\mathcal O_X(X)=\bigcap_x\mathcal O_{X,x}$ in $K(X)$ (a family of germs in $K(X)$ that are all regular glues, since restriction maps to $K(X)$ are injective). Hence $\alpha\in\mathcal O_X(X)$.

\emph{$\mathcal O_X(X)\subseteq k'$.} Since $X$ is proper over $k$, $\mathcal O_X(X)$ is integral over $k$ (Prop.~3.3.18). Thus every element of $\mathcal O_X(X)\subseteq K(X)$ is integral, hence algebraic, over $k$, so it lies in $k'$. \hfill$\square$

\exercise{1.13}{Let $X$ be a reduced algebraic variety. Show that $X$ is normal at all of the generic points of $X$. Show that $X$ contains an everywhere dense normal open subvariety.}
\pf
\textbf{Conventions.} ``Normal at $x$'' means that $\mathcal O_{X,x}$ is a normal domain (Def.~1.1). Since $X$ need not be irreducible, ``normal open subvariety'' is understood as: an open subscheme all of whose local rings are normal domains. It is a disjoint union of normal schemes in the sense of Def.~1.1, and it is normal in that sense if $X$ is irreducible.

\textbf{Generic points.} The scheme $X$ is Noetherian, so it has finitely many irreducible components $X_1,\dots,X_n$ with generic points $\eta_1,\dots,\eta_n$. The point $\eta_i$ lies in no $X_j$ with $j\ne i$ (otherwise $X_i=\overline{\{\eta_i\}}\subseteq X_j$, forcing $X_i=X_j$). The primes of $\mathcal O_{X,\eta_i}$ correspond to the points of $X$ that generalize $\eta_i$, and $\eta_i$ is the only one. So $\mathcal O_{X,\eta_i}$ is a local ring with a unique prime ideal, and it is reduced, so its maximal ideal (the nilradical) is $0$. Hence $\mathcal O_{X,\eta_i}$ is a field, in particular a normal domain.

\textbf{Dense normal open.} Put $U_i=X_i\setminus\bigcup_{j\ne i}X_j$. It is open in $X$ (a finite union of closed sets is closed) and contains $\eta_i$. It is irreducible (open in $X_i$) and reduced, hence integral, and it is an algebraic variety (open in $X$). The $U_i$ are pairwise disjoint. By Cor.~1.30, the set $N_i$ of normal points of $U_i$ is open in $U_i$, hence in $X$, and it contains $\eta_i$ by the first part. Put $N=\bigcup_iN_i$. It is open, the $N_i$ are pairwise disjoint, and $N$ contains all $\eta_i$, so $\overline N\supseteq\bigcup_i\overline{\{\eta_i\}}=X$, i.e.\ $N$ is dense. For $x\in N_i$ we have $\mathcal O_{N,x}=\mathcal O_{X,x}=\mathcal O_{U_i,x}$, which is normal. Thus $N$ is a dense open subvariety of $X$ all of whose local rings are normal domains. \hfill$\square$

\exercise{1.14}{Let $X$ be an integral scheme of dimension $1$, and let $f:X\to Y$ be a separated birational morphism of finite type from $X$ to a Dedekind scheme. Show that $f$ is an open immersion.}
\pf
\textbf{Conventions.} $Y$ is normal, hence integral, locally Noetherian of dimension $\le1$. Let $\xi,\eta$ be the generic points of $X,Y$. \emph{Birational} (Ex.~3.2.6) means $f(\xi)=\eta$ and $f^\#_\xi:K(Y)\to K(X)$ is an isomorphism; we identify $K:=K(Y)=K(X)$. The hypothesis $\dim X=1$ will not be used.

\textbf{Lemma A.} Let $S$ be a scheme, $Z$ an integral $S$-scheme with function field $K(Z)$ and $j:\operatorname{Spec}K(Z)\to Z$ the canonical morphism, and $T$ a separated $S$-scheme. If two $S$-morphisms $g_1,g_2:Z\to T$ satisfy $g_1\circ j=g_2\circ j$, then $g_1=g_2$.

\emph{Proof.} Both send the generic point $\zeta$ of $Z$ to the same point $p$. Take an affine open $W=\operatorname{Spec}R\ni p$ of $T$ and a non-empty affine open $V_0=\operatorname{Spec}C\subseteq g_1^{-1}(W)\cap g_2^{-1}(W)$. Then $C$ is a domain with fraction field $K(Z)$. The restrictions $g_i|_{V_0}:V_0\to W$ correspond to ring maps $\gamma_i:R\to C$. The compositions $R\to C\subseteq K(Z)$ correspond to $g_i\circ j$, so they agree, and since $C\to K(Z)$ is injective, $\gamma_1=\gamma_2$. So $g_1,g_2$ agree on the dense open $V_0$. As $Z$ is reduced and $T$ separated over $S$, Prop.~3.3.11 gives $g_1=g_2$. $\square$

\textbf{Step 1: local rings.} Let $x\in X$, $y=f(x)$, $R=\mathcal O_{Y,y}$, $S=\mathcal O_{X,x}$. These are domains with fraction field $K$, and $R\subseteq K$, $S\subseteq K$ (localizations of domains). The local homomorphism $f^\#_x:R\to S$ is compatible with the embeddings into $K$ (because $f^\#_\xi$ is the identity of $K$), so $R\subseteq S\subseteq K$. We claim $R=S$.

If $x=\xi$, then $y=\eta$ and $R=S=K$. Let $x\ne\xi$. Then $S$ is not a field, since $\xi$ is a proper generization of $x$ and $\operatorname{Spec}S$ has at least two points. If $y=\eta$, then $K=R\subseteq S\subseteq K$ would make $S$ a field; so $y\ne\eta$. Then $R$ is a normal Noetherian local domain of dimension $\ge1$, hence of dimension $1$ (as $\dim Y\le1$), hence a principal ideal domain (Prop.~1.12) and a discrete valuation ring; let $\pi$ be a uniformizer. As $f^\#_x$ is local, $\pi\in\mathfrak m_x$, so $\pi$ is not a unit of $S$. If $s\in S\setminus R$, write $s=u\pi^n$ with $u\in R^*$ and $n<0$; then $\pi^{-1}=s\,u^{-1}\pi^{-n-1}\in S$, contradicting that $\pi$ is not a unit of $S$. Hence $S=R$.

In particular we showed: $x\ne\xi\Rightarrow f(x)\ne\eta$.

\textbf{Step 2: $f$ is injective.} Let $f(x_1)=f(x_2)=y$. If $y=\eta$, then $x_1=x_2=\xi$ by the remark at the end of Step 1. Otherwise $x_1,x_2\neq\xi$ and $R:=\mathcal O_{Y,y}=\mathcal O_{X,x_1}=\mathcal O_{X,x_2}$ as subrings of $K$. Let $g_i:\operatorname{Spec}R=\operatorname{Spec}\mathcal O_{X,x_i}\to X$ be the canonical morphisms. They are $Y$-morphisms (both composites with $f$ are the canonical $\operatorname{Spec}\mathcal O_{Y,y}\to Y$), and both restrict on $\operatorname{Spec}K$ to the canonical morphism $\operatorname{Spec}K(X)\to X$. As $X$ is separated over $Y$ and $\operatorname{Spec}R$ is integral, Lemma A gives $g_1=g_2$. Evaluating at the closed point of $\operatorname{Spec}R$ gives $x_1=x_2$.

\textbf{Step 3: $f$ is a local isomorphism.} Let $x\in X$, $y=f(x)$. Choose affine opens $W=\operatorname{Spec}A\ni y$ and $V=\operatorname{Spec}B\ni x$ with $f(V)\subseteq W$; $B=A[b_1,\dots,b_n]$ is a finitely generated $A$-algebra since $f$ is of finite type. Both are domains inside $K$, with $A\subseteq B$. Let $\mathfrak p,\mathfrak q$ be the primes of $y,x$. By Step 1, $B_{\mathfrak q}=A_{\mathfrak p}$. So each $b_i\in A_{\mathfrak p}$ and there is $s\in A\setminus\mathfrak p$ with all $b_i\in A_s$. Since $\mathfrak q\cap A=\mathfrak p$, $s\notin\mathfrak q$, so $x\in V':=D(s)\subseteq V$. Now $B_s=A_s[b_1,\dots,b_n]=A_s$ (the inclusion $A_s\subseteq B_s\subseteq A_s$). So $f$ maps $V'$ isomorphically onto the open subset $D_A(s)$ of $Y$.

\textbf{Conclusion.} By Step 3, $f$ is an open map and a local isomorphism; by Step 2 it is injective. Hence $f$ is a homeomorphism onto the open subset $f(X)$ and induces isomorphisms on all stalks, i.e.\ an isomorphism $X\to f(X)$ onto an open subscheme of $Y$. So $f$ is an open immersion. \hfill$\square$

\exercise{1.15}{Let $X$ be a locally Noetherian, separated, integral scheme, and $U$ an affine open subset of $X$ different from $X$. Let us suppose that the normalization $X'$ of $X$ is also locally Noetherian (e.g., if $X'\to X$ is finite). We want to show that the irreducible components of $X\setminus U$ are of codimension $1$ in $X$. Let $\xi$ be a generic point of $X\setminus U$.
\begin{enumerate}[label=(\alph*),leftmargin=*]
\item Show that there exists an affine open neighborhood $V$ of $\xi$ in $X$ such that $\xi$ is the unique generic point of $V\setminus(U\cap V)$.
\item Show that $V\cap U$ is an affine open subset of $V$. Using Theorem~1.14, show that $\dim\mathcal O_{X,\xi}=1$ if $X$ is normal.
\item Let $\pi:X'\to X$ be the normalization morphism. Show that any point of $\pi^{-1}(\xi)$ is a generic point of $X'\setminus U'$, where $U'=\pi^{-1}(U)$. Show that $\dim\mathcal O_{X,\xi}=\max_{\xi'\in\pi^{-1}(\xi)}\dim\mathcal O_{X',\xi'}=1$.
\item Show that the statement still holds if $X$ is not necessarily integral, but is such that the normalization of its (reduced) irreducible components are locally Noetherian.
\end{enumerate}}
\pf
\textbf{Hypothesis needed.} We assume $U\ne\emptyset$. (If $U=\emptyset$, which is an affine open subset different from $X$, then $X\setminus U=X$ has the single component $X$, of codimension $0$.) In (d) the analogous hypothesis is that $U$ meets every irreducible component of $X$, and without it (d) is false (see (d)). Recall $\operatorname{codim}(\overline{\{\xi\}},X)=\dim\mathcal O_{X,\xi}$. The statement to prove is therefore $\dim\mathcal O_{X,\xi}=1$.

\textbf{(a).} Let $V_0$ be an affine open neighborhood of $\xi$. It is a Noetherian space, so $V_0\cap(X\setminus U)$ has finitely many irreducible components $Z_0,Z_1,\dots,Z_r$, where $Z_0$ is the closure of $\{\xi\}$ in $V_0$ (the closure of $\{\xi\}$ is an irreducible component of $X\setminus U$, so $Z_0$ is a component of the trace). We have $\xi\notin Z_j$ for $j\ge1$ (otherwise $Z_0\subseteq Z_j$, hence $Z_0=Z_j$). Then $V_1=V_0\setminus\bigcup_{j\ge1}Z_j$ is an open neighborhood of $\xi$ and $V_1\cap(X\setminus U)=Z_0\setminus\bigcup_{j\ge1}Z_j$ is a non-empty open subset of the irreducible $Z_0$, hence irreducible with generic point $\xi$. Finally choose a principal open $V=D(g)\subseteq V_1$ containing $\xi$ (it is affine). Then $V\cap(X\setminus U)$ is a non-empty open subset of $V_1\cap(X\setminus U)$, hence irreducible with generic point $\xi$.

\textbf{(b).} Since $X$ is separated and $U,V$ are affine, $U\cap V$ is affine (it is the preimage of the closed diagonal in $U\times V$, a closed subscheme of an affine scheme). Now assume $X$ is normal. Write $V=\operatorname{Spec}A$. Then $A$ is a normal Noetherian domain (Prop.~1.5, and $X$ is locally Noetherian). Let $F=V\setminus(U\cap V)=V\cap(X\setminus U)$; by (a) it is irreducible with generic point $\xi$; let $\mathfrak p$ be the prime of $\xi$, so $F=V(\mathfrak p)$ and $\operatorname{codim}(F,V)=\operatorname{ht}\mathfrak p=\dim\mathcal O_{X,\xi}$.

Suppose $\operatorname{ht}\mathfrak p\ge2$. By Theorem~1.14 the restriction $A=\mathcal O(V)\to\mathcal O(V\setminus F)=\mathcal O(U\cap V)$ is an isomorphism. As $U\cap V$ is affine, the open immersion $U\cap V\to V$ is $\operatorname{Spec}$ of this restriction map, hence an isomorphism. So $F=\emptyset$, contradicting $\xi\in F$. Hence $\operatorname{ht}\mathfrak p\le1$. Also $\operatorname{ht}\mathfrak p\ge1$: the generic point of $X$ lies in $U\neq\emptyset$, so $\xi$ is not the generic point of $X$ and $\mathcal O_{X,\xi}$ is not a field. So $\dim\mathcal O_{X,\xi}=1$.

\textbf{(c).} The morphism $\pi$ is integral and birational (Prop.~1.22). By Def.~1.20, preimages of affine opens are affine, so $U'$ is affine and $\pi$ is affine, hence separated. So $X'$ is separated (over $\mathbb Z$, as $X$ is). The ring map of $\pi$ over each affine open is injective and integral, so $\pi$ is surjective (Prop.~2.5.10(b)). Thus $U'\ne X'$ (a point over $\xi$ lies in $X'\setminus U'$) and $U'\neq\emptyset$ (it contains the generic point). The scheme $X'$ is integral, normal, and locally Noetherian by hypothesis.

\emph{Points over $\xi$ are generic points of $X'\setminus U'$.} Take $V=\operatorname{Spec}A$ as in (a), so $\pi^{-1}(V)=\operatorname{Spec}A'$ with $A\subseteq A'$ integral (both inside $K(X)$). Let $\mathfrak p\subset A$ be the prime of $\xi$ and $F=V(\mathfrak p)$ as before. Then $\pi^{-1}(F)=V(\mathfrak pA')=\pi^{-1}(V)\cap(X'\setminus U')$. Let $\mathfrak q\subset A'$ be a prime with $\mathfrak q\cap A=\mathfrak p$ (these are the points of $\pi^{-1}(\xi)$). If $\mathfrak q'$ is a prime with $\mathfrak pA'\subseteq\mathfrak q'\subseteq\mathfrak q$, then $\mathfrak p\subseteq\mathfrak q'\cap A\subseteq\mathfrak q\cap A=\mathfrak p$, so $\mathfrak q'$ and $\mathfrak q$ have the same contraction, and by incomparability for integral extensions (proof of Prop.~2.5.10(a)) $\mathfrak q'=\mathfrak q$. So $\mathfrak q$ is a minimal prime over $\mathfrak pA'$, i.e.\ a generic point of an irreducible component of $\pi^{-1}(F)$. An irreducible component of the trace on an open subset of a closed set $C$ has closure an irreducible component of $C$, so $\mathfrak q$ is a generic point of an irreducible component of $X'\setminus U'$.

\emph{Equality of dimensions.} Applying (b) to $X'$, $U'$ and each $\xi'\in\pi^{-1}(\xi)$ gives $\dim\mathcal O_{X',\xi'}=1$. It remains to show $\dim\mathcal O_{X,\xi}=\max_{\xi'}\dim\mathcal O_{X',\xi'}$. Let $T=A\setminus\mathfrak p$. Then $A_{\mathfrak p}\subseteq T^{-1}A'=:A'_{\mathfrak p}$ is injective and integral, so $\dim A_{\mathfrak p}=\dim A'_{\mathfrak p}$ (Prop.~2.5.10(b)). The maximal ideals of $A'_{\mathfrak p}$ are the $\mathfrak qA'_{\mathfrak p}$ with $\mathfrak q\cap A=\mathfrak p$ (a prime of an integral extension is maximal iff its contraction is, by Prop.~2.5.10(b)), and $(A'_{\mathfrak p})_{\mathfrak qA'_{\mathfrak p}}=A'_{\mathfrak q}=\mathcal O_{X',\mathfrak q}$. By Prop.~2.5.8(c), $\dim A'_{\mathfrak p}=\max_{\mathfrak q}\dim A'_{\mathfrak q}$ (the set of such $\mathfrak q$ is non-empty by surjectivity). Therefore $\dim\mathcal O_{X,\xi}=\max_{\xi'\in\pi^{-1}(\xi)}\dim\mathcal O_{X',\xi'}=1$.

\textbf{(d).} Let $X$ be locally Noetherian and separated, with reduced irreducible components $X_i$ (integral closed subschemes) whose normalizations are locally Noetherian, and let $U\ne\emptyset$ be an affine open subset with $X\setminus U\neq\emptyset$.

\emph{Hypothesis needed.} We must assume that $U$ meets every irreducible component of $X$. Otherwise the statement is false: for $X=\mathbb A^1_k\sqcup\mathbb A^1_k$ (locally Noetherian, separated, with normal Noetherian components) and $U$ the first copy, $X\setminus U$ is the second copy, which is an irreducible component of $X$ of codimension $0$.

Let $\xi$ be a generic point of an irreducible component $D$ of $X\setminus U$ (so $\dim\mathcal O_{X,\xi}$ is what we must compute). The primes of $\mathcal O_{X,\xi}$ are the points generalizing $\xi$; since $\mathcal O_{X,\xi}$ is Noetherian it has finitely many minimal primes, corresponding to the (finitely many) irreducible components $X_1,\dots,X_m$ of $X$ through $\xi$, and $\mathcal O_{X,\xi}/\mathfrak p_i=\mathcal O_{X_i,\xi}$ for the reduced structure. Every chain of primes starts at a minimal prime, so $\dim\mathcal O_{X,\xi}=\max_{i\le m}\dim\mathcal O_{X_i,\xi}$.

Fix $i\le m$. The scheme $X_i$ is integral, locally Noetherian, separated, with locally Noetherian normalization. The set $U_i=U\cap X_i$ is non-empty (by the hypothesis) and affine (a closed subscheme of the affine $U$). It differs from $X_i$ because $\xi\in X_i\setminus U_i$. The point $\xi$ is a generic point of an irreducible component of $X_i\setminus U_i=X_i\cap(X\setminus U)$: the closure $D$ of $\{\xi\}$ is contained in $X_i\setminus U_i$ (as $\xi\in X_i$ and $X_i$ is closed), and any irreducible closed subset of $X_i\setminus U_i\subseteq X\setminus U$ containing $D$ equals $D$ by maximality of $D$. By (c) applied to $(X_i,U_i)$, $\dim\mathcal O_{X_i,\xi}=1$. Hence $\dim\mathcal O_{X,\xi}=1$, i.e.\ every irreducible component of $X\setminus U$ has codimension $1$ in $X$. \hfill$\square$

\exercise{1.16}{Let $C$ be a normal proper curve over a field $k$. We are going to show that $C$ is projective over $k$.
\begin{enumerate}[label=(\alph*),leftmargin=*]
\item Let $\{U_i\}_i$ be a finite covering of $C$ by affine open subsets. Each $U_i$ admits an open immersion into a projective variety $Y_i$ over $k$ (Exercise 2.3.18). Let $Y$ be the fibered product of the $Y_i$ over $k$. Show that the canonical morphism $\bigcap_iU_i\to Y$ extends uniquely to a proper morphism $h:C\to Y$.
\item Let $Z=h(C)$ be the closed subvariety of $Y$ endowed with the reduced structure. Show that $h$ induces a surjective morphism $f:C\to Z$ and that the projection $Y\to Y_i$ induces a dominant morphism $Z\to Y_i$.
\item Show that for any open subset $U_i$, the open immersion $U_i\to Y_i$ factors into $U_i\to C\to Z\to Y_i$. Deduce from this that $f^\#_x:\mathcal O_{Z,f(x)}\to\mathcal O_{C,x}$ is an isomorphism for every $x\in U_i$. Show that $f$ is an isomorphism and that $C$ is projective over $k$.
\end{enumerate}}
\pf
A curve is an algebraic variety of dimension $1$. As $C$ is normal, it is irreducible with normal local rings, hence integral. We may discard empty $U_i$ and take $r\ge1$ pieces $U_1,\dots,U_r$.

\textbf{(a).} \emph{Construction of $Y_i$.} The affine variety $U_i$ (integral) is a closed subscheme of some $\mathbb A^{n}_k\subseteq\mathbb P^n_k$. Let $Y_i$ be the closure of $U_i$ in $\mathbb P^n_k$ with its reduced structure. It is an integral projective $k$-scheme, and $U_i=Y_i\cap\mathbb A^n_k$ is a dense open subscheme of $Y_i$ (because $U_i$ is closed in $\mathbb A^n_k$). Let $j_i:U_i\to Y_i$ be the open immersion. The fibered product $Y=Y_1\times_k\dots\times_kY_r$ is projective over $k$ (Segre embedding); let $p_i:Y\to Y_i$ be the projections.

Let $W=\bigcap_iU_i$; it is non-empty and open since $C$ is irreducible. Let $j:W\to Y$ be the morphism with $p_i\circ j=j_i|_W$. Apply Cor.~1.17 to $S=\operatorname{Spec}k$: $C$ is normal of finite type over $k$ of dimension $1$, $Y\to S$ is proper, and $j$ is defined on the non-empty open set $W$. It extends uniquely to a morphism $h:C\to Y$. Since $C\to\operatorname{Spec}k$ is proper and $Y\to\operatorname{Spec}k$ is separated, $h$ is proper (Prop.~3.3.16(e)).

\textbf{(b).} A proper morphism is closed, so $Z=h(C)$ is closed in $Y$; it is irreducible (image of an irreducible space), so with the reduced structure it is an integral closed subscheme, proper over $k$. Since $C$ is reduced, $h$ factors through a morphism $f:C\to Z$, which is surjective by definition of $Z$. Let $q_i=p_i|_Z:Z\to Y_i$. It is proper (Prop.~3.3.16(e)), hence closed. Its image contains $p_i(h(W))=j_i(W)$, which is dense in $Y_i$ ($W$ is dense in $U_i$, and $U_i$ is dense in $Y_i$). So $q_i(Z)=Y_i$: the morphism $q_i$ is surjective, in particular dominant.

\textbf{(c).} \emph{Factorization.} The morphisms $p_i\circ h|_{U_i}$ and $j_i$ from $U_i$ to $Y_i$ are $k$-morphisms that agree on the dense open $W$, $U_i$ is reduced, and $Y_i$ is separated over $k$. By Prop.~3.3.11 they coincide: $j_i=q_i\circ f|_{U_i}$. So $U_i\to C\to Z\to Y_i$ is the open immersion $j_i$.

\emph{Stalks.} Let $x\in U_i$ and $z=f(x)$. Then $\mathcal O_{Y_i,j_i(x)}\xrightarrow{q_i^\#}\mathcal O_{Z,z}\xrightarrow{f^\#_x}\mathcal O_{C,x}$ has composite $j^\#_{i,x}$, an isomorphism. So $f^\#_x$ is surjective. It is also injective: $f$ is surjective and $C$ is irreducible, so $f$ maps the generic point of $C$ to the generic point of $Z$, and $f^\#_x$ is compatible with the embeddings $\mathcal O_{Z,z}\subseteq K(Z)\hookrightarrow K(C)\supseteq\mathcal O_{C,x}$ (all rings are domains, $Z$ and $C$ integral). So $f^\#_x$ is an isomorphism for every $x\in\bigcup U_i=C$. In particular $K(Z)\to K(C)$ is an isomorphism and we identify the two fields.

\emph{Injectivity of $f$.} Let $f(x_1)=f(x_2)=z$. By the above, $\mathcal O_{C,x_1}$ and $\mathcal O_{C,x_2}$ both equal $R:=\mathcal O_{Z,z}$ as subrings of $K(C)$. If $R=K(C)$ is a field, then $x_1=x_2$ is the generic point of $C$ (the local ring of an integral scheme at a point is a field only at the generic point). Otherwise $R$ is a normal Noetherian local domain of dimension $1$ (as $\dim C=1$ and $x_i$ is not generic), hence a discrete valuation ring. Let $g_i:\operatorname{Spec}R=\operatorname{Spec}\mathcal O_{C,x_i}\to C$ be the canonical morphisms. They are $k$-morphisms that agree on $\operatorname{Spec}K(C)$, and $C$ is separated over $k$, so $g_1=g_2$ by Lemma~A of Ex.~1.14 (applied with $S=\operatorname{Spec}k$, $Z=\operatorname{Spec}R$, $T=C$). Hence $x_1=g_1(\text{closed point})=g_2(\text{closed point})=x_2$.

\emph{Conclusion.} The morphism $f$ is bijective and proper (Prop.~3.3.16(e), $Z$ being separated), hence closed, so it is a homeomorphism. For a homeomorphism the stalk of $f_*\mathcal O_C$ at $z$ is $\mathcal O_{C,f^{-1}(z)}$, so $f^\#:\mathcal O_Z\to f_*\mathcal O_C$ is an isomorphism of sheaves (on stalks, by the above). Thus $f$ is an isomorphism of schemes. Therefore $C\simeq Z$ is a closed subscheme of the projective $k$-scheme $Y$, so $C$ is projective over $k$. \hfill$\square$

\exercise{1.17}{Let $U$ be an integral affine algebraic curve over a field $k$.
\begin{enumerate}[label=(\alph*),leftmargin=*]
\item Show that there exists a proper curve $\hat U$ over $k$ verifying the following properties: (1) there exists an open immersion $i:U\to\hat U$; (2) the points of $\hat U\setminus U$ are normal in $\hat U$.
\item Show that $\hat U$ is the `smallest proper completion' of $U$, in the sense that any morphism $U\to X$ into a proper variety over $k$ factors uniquely through the open immersion $U\to\hat U$. Deduce from this that $\hat U$ is unique up to isomorphism.
\item Show that there exists a finite morphism $U\to\mathbb A^1_k$, and that it extends to a finite morphism $f:\hat U\to\mathbb P^1_k$ such that $U=f^{-1}(\mathbb A^1_k)$.
\end{enumerate}}
\pf
\textbf{Conventions.} In (a) and (b) the curve $\hat U$ is required to be \emph{integral} (equivalently $i(U)$ is dense in $\hat U$); without this the uniqueness in (b) fails (one could add a disjoint proper curve). Morphisms in (b) are $k$-morphisms. We call such a pair $(\hat U,i)$ a \emph{completion} of $U$. We write $U=\operatorname{Spec}A$, $K=K(U)=\operatorname{Frac}(A)$. Lemma~A of Ex.~1.14 is used again: two $S$-morphisms from an integral scheme to a separated $S$-scheme that agree on $\operatorname{Spec}K(Z)$ are equal.

\textbf{Construction (giving (a) and the existence in (c)).} $A$ is a finitely generated $k$-domain of dimension $1$. By Noether normalization (Prop.~2.1.9) there is $t\in A$, transcendental over $k$, with $k[t]\subseteq A$ finite. So $\psi:U\to\mathbb A^1_k=\operatorname{Spec}k[t]$ is a finite morphism; this is the first claim of (c). Put $s=t^{-1}\in K$; the field $K$ is finite over $k(t)=k(s)$. Let $C_0$ be the integral closure of $k[s]$ in $K$; it is a normal domain, finite over $k[s]$ (Prop.~1.27 applied to $\operatorname{Spec}k[s]$ and $L=K$), of dimension $1$. Put
\[ C'=C_0\cap A_t,\quad\text{inside }K,\ \text{where }A_t=A[s]. \]

(i) $C'$ is a finite $k[s]$-module (a submodule of the finite module $C_0$ over the Noetherian ring $k[s]$) and $k[s]\subseteq C'\subseteq C_0$, with $C_0$ integral over $C'$.

(ii) $C'[1/s]=A_t$. The inclusion $\subseteq$ is clear. Let $a\in A_t$. As $A$ is finite over $k[t]$, $A_t$ is finite over $k[t,t^{-1}]=k[s,s^{-1}]$, so $a$ satisfies an equation $a^n+c_1a^{n-1}+\dots+c_n=0$ with $c_i\in k[s,s^{-1}]$. Choose $N$ with $s^Nc_i\in k[s]$ for all $i$. Then $s^Na$ satisfies the monic equation with coefficients $s^{iN}c_i=s^{(i-1)N}\cdot s^Nc_i\in k[s]$, so $s^Na\in C_0\cap A_t=C'$ and $a=s^{-N}(s^Na)\in C'[1/s]$.

(iii) \emph{For every prime $\mathfrak q'$ of $C'$ with $s\in\mathfrak q'$, the ring $C'_{\mathfrak q'}$ is normal.} Let $J=\{c\in C_0\mid cC_0\subseteq C'\}$ (the conductor); it is an ideal of $C_0$, contained in $C'$. It is non-zero: $C_0$ is a finitely generated $k[s]$-module inside $K=\operatorname{Frac}(C')$ (by (ii), $\operatorname{Frac}C'=\operatorname{Frac}A_t=K$), so a common denominator $v\in C'\setminus\{0\}$ of generators satisfies $vC_0\subseteq C'$. Moreover $J$ is saturated with respect to $s$: if $x\in C_0$ and $sx\in J$, then for $y\in C_0$ we get $sxy\in C'\subseteq A_t$, hence $xy\in A_t\cap C_0=C'$ (as $s$ is a unit of $A_t$), so $x\in J$.

Suppose $J\subseteq\mathfrak q'$. By Prop.~2.5.10(b) there is a prime $\mathfrak q$ of $C_0$ with $\mathfrak q\cap C'=\mathfrak q'$; then $J\subseteq\mathfrak q$ and $s\in\mathfrak q$. As $C_0$ is a Noetherian domain of dimension $1$ and $J\neq0$, the prime $\mathfrak q\neq0$ is maximal and minimal over $J$, hence an associated prime of $C_0/J$: $\mathfrak q=\operatorname{Ann}(\bar z)$ for some $z\in C_0$ with $z\notin J$. Then $sz\in J$, so $z\in J$ by saturation, a contradiction. Hence $J\not\subseteq\mathfrak q'$. Pick $u\in J\setminus\mathfrak q'$. For every $x\in C_0$ we have $ux\in C'$, so $x=ux/u\in C'_{\mathfrak q'}$. Thus $C'_{\mathfrak q'}=(C'\setminus\mathfrak q')^{-1}C_0$, a localization of the normal domain $C_0$, hence normal (Lemma~1.4(b)).

\emph{Gluing.} Let $E=\operatorname{Spec}C'$. By (ii), the open subschemes $D(t)=\operatorname{Spec}A_t\subseteq U$ and $D(s)=\operatorname{Spec}C'[1/s]\subseteq E$ are the same $\operatorname{Spec}A_t$. Let $\hat U=U\cup_{\operatorname{Spec}A_t}E$. Let $f:\hat U\to\mathbb P^1_k=\operatorname{Spec}k[t]\cup\operatorname{Spec}k[s]$ be given on $U$ by $\psi$ (from $k[t]\subseteq A$) and on $E$ by $k[s]\subseteq C'$. They agree on the overlap since both come from $k[s,s^{-1}]\subseteq A_t$. We have $f^{-1}(\mathbb A^1_k)=U$ (the points of $E$ outside $V(s)$ lie in $D(s)\subseteq U$) and $f^{-1}(\operatorname{Spec}k[s])=E$. Both are finite over their targets, so $f$ is finite.

Consequently $\hat U$ is separated (finite morphisms are affine, and $\mathbb P^1_k$ is separated), of finite type, and proper over $k$ (a finite morphism is proper, and $\mathbb P^1_k$ is proper; Prop.~3.3.16(b)). It is integral, being glued from the integral schemes $U,E$ along a non-empty open set, and has dimension $1$. $U$ is an open subscheme (so (1) holds), and $\hat U\setminus U=V(s)\subseteq E$. By (iii) its local rings are normal, so (2) holds. This proves (a) and the existence statement in (c).

\textbf{(b).} Let $(\hat U,i)$ be any completion; identify $U$ with an open dense subscheme. The set $T=\hat U\setminus U$ is closed and proper in the integral curve $\hat U$, so it is a finite set of closed points. Let $\varphi:U\to X$ be a $k$-morphism to a proper variety.

\emph{Uniqueness} of an extension $\Phi:\hat U\to X$ follows from Lemma~A ($\hat U$ integral, $X$ separated, and two extensions agree on $U\ni$ generic point).

\emph{Existence.} Let $x\in T$ and $R=\mathcal O_{\hat U,x}$. It is a normal Noetherian local domain of dimension $1$ (as $x$ is a closed point of a curve, normal by (2)), hence a discrete valuation ring (Prop.~1.12) with fraction field $K(\hat U)=K(U)=K$. The morphism $\operatorname{Spec}K\to U\xrightarrow{\varphi}X$ extends uniquely to a $k$-morphism $\operatorname{Spec}R\to X$ by the valuative criterion (Cor.~3.3.26, as $X\to\operatorname{Spec}k$ is proper), exactly as in the proof of Prop.~1.16. By Ex.~3.2.4 (as $X$ is of finite type over $\operatorname{Spec}k$) it extends to a morphism $g_x:U_x\to X$ on an open neighborhood $U_x$ of $x$ in $\hat U$; shrinking, $U_x\cap T=\{x\}$. The restrictions of $g_x$ and $\varphi$ to the integral scheme $U\cap U_x$ agree on $\operatorname{Spec}K$, hence agree by Lemma~A. Likewise $g_x$ and $g_{x'}$ agree on $U_x\cap U_{x'}$. So $\varphi$ and all $g_x$ glue to a morphism $\Phi$ on $U\cup\bigcup_{x\in T}U_x=\hat U$, which extends $\varphi$.

\emph{Uniqueness of $\hat U$.} Let $(\hat U,i)$ and $(\hat U',i')$ be two completions. Both $\hat U,\hat U'$ are proper varieties, so $i':U\to\hat U'$ extends to $\alpha:\hat U\to\hat U'$ with $\alpha\circ i=i'$, and $i$ extends to $\beta:\hat U'\to\hat U$ with $\beta\circ i'=i$. Then $\beta\circ\alpha$ and $\operatorname{id}_{\hat U}$ both extend $i:U\to\hat U$, so they are equal by uniqueness; similarly $\alpha\circ\beta=\operatorname{id}$. Hence $\hat U\simeq\hat U'$ compatibly with $i,i'$.

\textbf{(c), conclusion.} The finite morphism $\psi:U\to\mathbb A^1_k$ was constructed. For the completion built above, $f$ is a finite morphism extending $\psi$ (composed with $\mathbb A^1_k\subseteq\mathbb P^1_k$) with $f^{-1}(\mathbb A^1_k)=U$. By (b) any other completion is isomorphic to this one by an isomorphism compatible with $i$, and the extension of $\psi$ to it is unique, so the same holds for any completion $\hat U$ of $U$. \hfill$\square$

\exercise{1.18}{Let $k$ be a field, and let $P(x)\in k[x]$ be a monic polynomial.
\begin{enumerate}[label=(\alph*),leftmargin=*]
\item Show that $X=\operatorname{Spec}k[x,y]/(y^2-P(x))$ is integral if and only if $P(x)$ is not a square.
\item Let us suppose $\operatorname{char}(k)\neq2$. Show that $X$ is normal if and only if $P(x)$ is a polynomial without a square factor.
\item Let us suppose $\operatorname{char}(k)=2$ and $k$ is perfect. Then $P(x)$ can be written in a unique way as $P(x)=xQ_1(x)^2+Q_2(x)^2$. Show that $X$ is normal if and only if $P(x)\in k^*x+k[x^2]$. If $k$ is not perfect, show that this condition is sufficient.
\item Supposing $X$ integral, and $\operatorname{char}(k)\neq2$ or $k$ perfect, determine the normalization of $X$.
\end{enumerate}}
\pf
Let $B=k[x,y]/(y^2-P)$, $K_0=k(x)$. As a $k[x]$-module $B$ is free with basis $1,y$ (division by the monic polynomial $y^2-P$), so every element of $B$ is uniquely $a+by$ with $a,b\in k[x]$, and $B$ is integral over $k[x]$. By Prop.~1.5, $X=\operatorname{Spec}B$ is normal iff $B$ is a normal domain.

\textbf{(a).} If $P=Q^2$ with $Q\in k[x]$, then $(y-Q)(y+Q)=y^2-P=0$ in $B$, while $y\pm Q\neq0$ (the coefficient of $y$ is $1$). So $B$ is not a domain. Conversely, suppose $P$ is not a square in $k[x]$. If $y^2-P$ were reducible in the UFD $k[x][y]$, being monic of degree $2$ in $y$ it would be a product $(y-a)(y-b)$ with $a,b\in k[x]$ (a factor of $y$-degree $0$ would divide the monic polynomial, hence be a unit); then $a+b=0$ and $ab=-P$, so $P=a^2$, a contradiction. So $y^2-P$ is irreducible, hence prime in the UFD $k[x,y]$, and $B$ is a domain. (We also need $B\ne0$, which is clear.)

If $B$ is a domain, then $y\notin K_0$ and $L:=\operatorname{Frac}B=K_0(y)=K_0\oplus K_0y$ has degree $2$ over $K_0$. Indeed $y\in K_0$ would make $P=y^2$ a square in $K_0$, hence in $k[x]$ ($k[x]$ is normal, Lemma~1.4(a)).

\textbf{(b).} Assume $\operatorname{char}k\neq2$. We read ``without a square factor'' as: $P\neq1$ and there is no non-constant $Q$ with $Q^2\mid P$ (the constant $P=1$ is a square, and $X$ is then not integral; so for $P=1$ the printed statement needs this reading).

\emph{Normal $\Rightarrow$ squarefree and $P\neq1$.} If $X$ is normal then $B$ is a domain, so $P$ is not a square by (a); in particular $P\ne1$. Suppose $P=Q^2R$ with $Q$ non-constant and $R\in k[x]$. Let $z=y/Q\in L$. Then $z^2=R\in k[x]$, so $z$ is integral over $k[x]$, hence over $B$. But $z=0+Q^{-1}y$ and $Q^{-1}\notin k[x]$, so $z\notin B$ by uniqueness of the form $a+by$ ($1,y$ is a basis of $L$ over $K_0$). So $B$ is not normal, a contradiction.

\emph{Squarefree and $P\neq1$ $\Rightarrow$ normal.} Then $P$ is not a square, so $B$ is a domain and $L=K_0\oplus K_0y$. The automorphism $\sigma$ of $L$ over $K_0$ with $\sigma(y)=-y$ exists (the minimal polynomial $T^2-P$ has the distinct roots $\pm y$, as $2\ne0$). Let $z=a+by\in L$ ($a,b\in K_0$) be integral over $B$; it is integral over $k[x]$, and so is $\sigma(z)$ (apply $\sigma$ to an integral equation with coefficients in $k[x]$). Hence $z+\sigma(z)=2a$ and $z\sigma(z)=a^2-Pb^2$ are elements of $K_0$ integral over the normal ring $k[x]$, so they lie in $k[x]$. So $a\in k[x]$ (as $2$ is invertible) and $Pb^2\in k[x]$. Write $b=u/v$ with $u,v\in k[x]$ coprime, $v\ne0$. Then $v^2\mid Pu^2$, so $v^2\mid P$, hence $v$ is a constant since $P$ has no square factor. Thus $b\in k[x]$ and $z\in B$. So $B$ is integrally closed in $L$, i.e.\ normal.

\textbf{(c).} Let $\operatorname{char}k=2$.

\emph{Decomposition.} If $k$ is perfect, write $P=\sum a_ix^i$. The even part is $\sum a_{2j}x^{2j}=Q_2^2$ with $Q_2=\sum\sqrt{a_{2j}}x^j$, and the odd part is $x\cdot Q_1^2$ with $Q_1=\sum\sqrt{a_{2j+1}}x^j$. If $xQ_1^2+Q_2^2=xQ_1'^2+Q_2'^2$, then $x(Q_1-Q_1')^2=(Q_2-Q_2')^2$ (characteristic $2$); the left side has odd degree if $Q_1\ne Q_1'$, the right side even degree, so $Q_1=Q_1'$ and then $Q_2=Q_2'$. Note $k[x^2]\supseteq Q_2^2$, and over a perfect field $k[x^2]$ consists of the squares of polynomials. Hence $P\in k^*x+k[x^2]$ iff $Q_1\in k^*$.

\emph{Sufficiency (any field $k$ of characteristic $2$).} Suppose $P=cx+S(x^2)$ with $c\in k^*$ and $S\in k[T]$. Then $P'=c\ne0$ (formal derivative in $x$), whereas a square $Q^2$ has derivative $2QQ'=0$; so $P$ is not a square and $B$ is a domain with $L=K_0\oplus K_0y$ by (a). Let $z=a+by\in L$ be integral over $B$ (hence over $k[x]$). Then $z^2=a^2+b^2P$ (characteristic $2$) is an element of $K_0$ integral over $k[x]$, so $a^2+b^2P\in k[x]$. Write $a=\alpha/v$, $b=\beta/v$ with $\alpha,\beta,v\in k[x]$, $v\neq0$ and $\gcd(\alpha,\beta,v)=1$. Then $\alpha^2+\beta^2P=v^2g$ with $g\in k[x]$. Differentiating, $(\alpha^2)'=0$ and $(\beta^2P)'=\beta^2P'=c\beta^2$, $(v^2g)'=v^2g'$. So $c\beta^2=v^2g'$, hence $v^2\mid\beta^2$, so $v\mid\beta$. Then $v^2\mid\alpha^2=v^2g-\beta^2P$, so $v\mid\alpha$. Therefore $v\mid\gcd(\alpha,\beta,v)=1$, so $a,b\in k[x]$ and $z\in B$. Hence $B$ is normal.

\emph{Necessity ($k$ perfect).} Suppose $B$ is normal, so $P$ is not a square and $Q_1\neq0$ (otherwise $P=Q_2^2$). Suppose $Q_1$ is non-constant. Put $z=(y+Q_2)/Q_1\in L$. Then $z^2=(y^2+Q_2^2)/Q_1^2=(P+Q_2^2)/Q_1^2=x$ (characteristic $2$), so $z$ is integral over $k[x]\subseteq B$. But $z=Q_2/Q_1+Q_1^{-1}y\notin B$, as $Q_1^{-1}\notin k[x]$. This contradicts normality. Hence $Q_1\in k^*$, i.e.\ $P\in k^*x+k[x^2]$.

So for $k$ perfect: $X$ normal $\iff P\in k^*x+k[x^2]$. For $k$ not perfect the condition is still sufficient, as shown, but not necessary: for $k=\mathbb F_2(a)$ with $a\notin k^2$ and $P=x^2+a$ (not in $k^*x+k[x^2]$), put $w=x+y$; then $w^2=x^2+y^2=a$ in $B$, and $B=k[x,w]/(w^2-a)=k(\sqrt a)[x]$ is a polynomial ring over a field, hence normal.

\textbf{(d).} Assume $B$ is a domain ($P$ not a square) and $\operatorname{char}k\ne2$ or $k$ perfect. In each case we exhibit a normal domain $B'$ with $B\subseteq B'\subseteq L$, integral over $B$; then $B'$ is the integral closure of $B$ in $L$ and $\operatorname{Spec}B'\to X$ is the normalization (Lemma~1.21).

\emph{Case $\operatorname{char}k\neq2$.} Factor $P=Q^2R$ with $Q,R\in k[x]$ monic and $R$ squarefree (take $Q$ the product of $p_i^{\lfloor e_i/2\rfloor}$ and $R$ the product of $p_i^{e_i-2\lfloor e_i/2\rfloor}$ over the prime factorization $P=\prod p_i^{e_i}$). $R\ne1$ since $P$ is not a square. Put $w=y/Q\in L$, so $w^2=R$, and $B'=k[x,w]=k[x][w]\subseteq L$; it is $\simeq k[x,w]/(w^2-R)$ (the polynomial $T^2-R$ is the minimal polynomial of $w$ over $K_0$ and $k[x]$ is normal), which is normal by (b) applied to $R$. Moreover $B\subseteq B'$ ($y=Qw$), $B'$ is integral over $B$ ($w^2=R\in k[x]$), and $\operatorname{Frac}B'=K_0(w)=K_0(y)=L$. So the normalization of $X$ is
\[ X'=\operatorname{Spec}k[x,w]/(w^2-R)\longrightarrow X,\qquad (x,w)\mapsto(x,y=Q(x)w),\]
where $P=Q^2R$ with $R$ squarefree.

\emph{Case $\operatorname{char}k=2$, $k$ perfect.} Write $P=xQ_1^2+Q_2^2$; $Q_1\neq0$ as $P$ is not a square. Put $z=(y+Q_2)/Q_1\in L$. Then $z^2=x$ (as in (c)). The element $z$ is transcendental over $k$ (since $z^2=x$ is), so $B'=k[z]$ is a polynomial ring, hence normal (Lemma~1.4(a)). We have $x=z^2\in B'$ and $y=Q_1(z^2)z+Q_2(z^2)\in B'$, so $B\subseteq B'$; $B'$ is integral over $B$ since $z^2=x\in B$; and $\operatorname{Frac}B'=k(z)\ni x,y$ equals $L$ (as $z\in L$ and $L=k(x,y)$). So the normalization of $X$ is
\[ X'=\mathbb A^1_k=\operatorname{Spec}k[z]\longrightarrow X,\qquad z\mapsto\bigl(x=z^2,\ y=Q_1(z^2)\,z+Q_2(z^2)\bigr). \]
(If $X$ is already normal, $Q_1\in k^*$ and this map is an isomorphism.) \hfill$\square$

\exercise{1.19}{Let $k$ be a field of characteristic $p>0$. Let $f:X\to Y$ be a finite surjective morphism of normal algebraic varieties over $k$. We identify $K(Y)$ with a subfield of $K(X)$ via $f$. Use Lemma~3.2.25 to show the following properties:
\begin{enumerate}[label=(\alph*),leftmargin=*]
\item If $K(X)^p\subseteq K(Y)$, show that $F_{X/k}$ factors into $f:X\to Y$ followed by a morphism $Y\to X^{(p)}$.
\item If $K(Y)\subseteq kK(X)^p$, show that $f$ factors into $F_{X/k}:X\to X^{(p)}$ followed by a morphism $X^{(p)}\to Y$.
\end{enumerate}}
\pf
\textbf{Summary of what is proved.} Part (a) is proved as stated. Part (b) is \emph{false as stated when $k$ is not perfect} (counterexample below). We prove it under the hypothesis $K(Y)\subseteq K(X)^p$, which is the same as the printed hypothesis when $k$ is perfect (then $k=k^p\subseteq K(X)^p$, so $kK(X)^p=K(X)^p$).

\textbf{Notation.} $X,Y$ are normal, hence integral. $X^{(p)}=X\times_kk$, where the second factor $k$ is a $k$-algebra via $a\mapsto a^p$; $F_{X/k}:X\to X^{(p)}$ is the relative Frobenius, $\phi:X^{(p)}\to X$ the first projection, and $\phi\circ F_{X/k}=F_X$ (Defs.~3.2.21, 3.2.23). For an affine open $\operatorname{Spec}B\subseteq X$ we have $X^{(p)}\supseteq\operatorname{Spec}(B\otimes_kk)$, with $B\otimes_kk$ viewed via the second factor, and $F_{X/k}$ corresponds to $\Phi:B\otimes_kk\to B$, $b\otimes a\mapsto ab^p$ (the paragraph before Lemma~3.2.25; Lemma~3.2.25 is the coordinate version). The image of $\Phi$ is the subring $k[B^p]$ generated by $k$ and $B^p$.

For an affine open $Y_0=\operatorname{Spec}A\subseteq Y$, $f^{-1}(Y_0)=\operatorname{Spec}B$ is affine and $A\to B$ is finite (as $f$ is finite). It is injective ($f$ is surjective and $A$ is a domain: the kernel lies in all primes, so it is nilpotent hence $0$). So $A\subseteq B\subseteq K(X)$ and $\operatorname{Frac}A=K(Y)$.

\textbf{(a).} Assume $K(X)^p\subseteq K(Y)$. Let $Y_0=\operatorname{Spec}A$, $f^{-1}(Y_0)=\operatorname{Spec}B$ as above. Each $b\in B$ is integral over $A$, hence so is $b^p$: if $b^n+a_{n-1}b^{n-1}+\dots+a_0=0$ with $a_i\in A$, then $(b^p)^n+a_{n-1}^p(b^p)^{n-1}+\dots+a_0^p=0$ in characteristic $p$. Moreover $b^p\in K(X)^p\subseteq K(Y)$ and $A$ is normal, so $b^p\in A$. Thus $B^p\subseteq A$, and since $k\subseteq A$, $k[B^p]\subseteq A$. The map
\[ \rho:B\otimes_kk\to A,\qquad b\otimes a\mapsto ab^p \]
is a well-defined ring homomorphism ($b\mapsto b^p$ is additive and multiplicative, and $(\lambda b,a)$ and $(b,\lambda^pa)$ have the same image), compatible with the $k$-structure from the second factor, and $\Phi=\iota\circ\rho$, where $\iota:A\subseteq B$. This gives $g_{Y_0}:Y_0\to X^{(p)}$ with $g_{Y_0}\circ f=F_{X/k}$ on $f^{-1}(Y_0)$.

\emph{Gluing.} We use: \textbf{Lemma B.} If $V\subseteq Y$ is open and $g,g':V\to T$ are morphisms of schemes with $g\circ f_V=g'\circ f_V$ ($f_V:f^{-1}V\to V$), then $g=g'$. Indeed, $f$ is surjective so $g,g'$ agree as maps of points. Fix an affine open $T_0\subseteq T$ and an affine open $V_0=\operatorname{Spec}A_0\subseteq g^{-1}(T_0)=g'^{-1}(T_0)$; then $f^{-1}(V_0)=\operatorname{Spec}B_0$ and $A_0\to B_0$ is injective as above. The ring maps $\mathcal O(T_0)\to A_0$ defined by $g,g'$ become equal after composing with $A_0\hookrightarrow B_0$, so they are equal. Hence $g=g'$ on the cover of $V$ by such $V_0$.

For two affine opens $Y_0,Y_1$ the morphisms $g_{Y_0},g_{Y_1}$ restricted to $Y_0\cap Y_1$ satisfy $g\circ f=F_{X/k}=g'\circ f$ there, so they agree by Lemma~B. So the $g_{Y_0}$ glue to $g:Y\to X^{(p)}$ with $g\circ f=F_{X/k}$. This proves (a).

\textbf{(b), under the hypothesis $K(Y)\subseteq K(X)^p$.} We construct $h':X\to Y$ with $h'\circ F_X=f$. Then $h=h'\circ\phi:X^{(p)}\to Y$ satisfies $h\circ F_{X/k}=h'\circ\phi\circ F_{X/k}=h'\circ F_X=f$, as required.

Let $V\subseteq Y$ be a non-empty open set and $c\in\mathcal O_Y(V)\subseteq K(Y)\subseteq K(X)^p$. Let $u\in K(X)$ be the unique element with $u^p=c$ (the Frobenius of the field $K(X)$ is injective). We claim $u\in\mathcal O_X(f^{-1}V)$. Let $x\in f^{-1}(V)$, $y=f(x)$. Then $c\in\mathcal O_{Y,y}\subseteq\mathcal O_{X,x}$, so $u$ is a root of the monic polynomial $T^p-c\in\mathcal O_{X,x}[T]$. As $\mathcal O_{X,x}$ is normal with fraction field $K(X)$, $u\in\mathcal O_{X,x}$. Hence $u\in\bigcap_{x\in f^{-1}V}\mathcal O_{X,x}=\mathcal O_X(f^{-1}V)$ ($X$ integral). Set $\rho_V(c)=u=c^{1/p}$. The maps $\rho_V:\mathcal O_Y(V)\to\mathcal O_X(f^{-1}V)$ are ring homomorphisms compatible with restriction, so they define $\rho:\mathcal O_Y\to f_*\mathcal O_X$. The stalk maps are local: if $c\in\mathfrak m_y$ and $x\in f^{-1}(y)$, then $\rho(c)^p=c\in\mathfrak m_x$ forces $\rho(c)\in\mathfrak m_x$ (prime ideal). So $(f,\rho)$ is a morphism of schemes $h':X\to Y$.

Finally $h'\circ F_X=f$: $F_X$ is the identity on the underlying space (Lemma~3.2.22(b)), so both sides have the same underlying map; and for $c\in\mathcal O_Y(V)$, $F_X^\#(\rho(c))=\rho(c)^p=c=f^\#(c)$. This proves (b) under $K(Y)\subseteq K(X)^p$, in particular for $k$ perfect with the printed hypothesis.

\textbf{(b) is false as printed for imperfect $k$.} Let $p=2$, $k=\mathbb F_2(t)$ and $B=k[x,y]/(y^2-x^3-t)$, $X=\operatorname{Spec}B$, $Y=\mathbb A^1_k=\operatorname{Spec}k[x]$, $f$ induced by $k[x]\subseteq B$.

$B$ is a normal domain. It is a domain as $x^3+t$ is not a square (its derivative $x^2$ is non-zero, whereas squares have derivative $0$). It is a $1$-dimensional Noetherian domain, so it suffices to check that each localization $B_{\mathfrak m}$ at a maximal ideal is a discrete valuation ring, i.e.\ that $\mathfrak mB_{\mathfrak m}$ is principal. Let $G=y^2+x^3+t$ and $M\supseteq(G)$ be a maximal ideal of $k[x,y]$. If $x\notin M$, then $\partial G/\partial x=x^2\notin M$, so $G\notin M^2$ (an element of $M^2$ has partial derivatives in $M$). If $x\in M$, then $y^2+t\in M$ and $M=(x,y^2+t)$ (the latter is maximal: $k[y]/(y^2+t)$ is a field), and $G=(y^2+t)+x^3\notin M^2$ since $x,y^2+t$ form a basis of $M/M^2$. In both cases $G\in M\setminus M^2$; as $\dim k[x,y]_M=2$ and $M/M^2$ has dimension $2$, $\bar G$ is part of a basis of $M/M^2$ and by Nakayama $\mathfrak mB_{\mathfrak m}=M k[x,y]_M/(G)$ is generated by one element. So $B$ is regular of dimension $1$ at all maximal ideals, hence normal. $A=k[x]$ is normal, and $f$ is finite (free of rank $2$) and surjective, so $f$ is a finite surjective morphism of normal varieties.

The hypothesis holds: $K(X)^2\ni x^2$ and $y^2=x^3+t$, so $kK(X)^2\ni x^3+t$, hence $x^3$ (as $t\in k$), hence $x=x^3/x^2$, so $K(Y)=k(x)\subseteq kK(X)^2$.

The conclusion fails: a morphism $h:X^{(2)}\to Y$ with $h\circ F_{X/k}=f$ corresponds to a ring homomorphism $\psi:k[x]\to B\otimes_kk$ with $\Phi\circ\psi=\iota$ ($\iota:k[x]\subseteq B$). Then $x=\Phi(\psi(x))\in\operatorname{Im}\Phi=k[B^2]$. But every element of $B^2$ has the form $(\alpha+\beta y)^2=\alpha^2+\beta^2(x^3+t)\in k[x^2,x^3]$ ($\alpha,\beta\in k[x]$), so $k[B^2]\subseteq k[x^2,x^3]$, which does not contain $x$. A contradiction. (Here the reason is that $X^{(2)}$ is the cuspidal curve $y'^2=x^3$, which is not normal, so $f$ cannot factor through it.) \hfill$\square$

\section*{4.2 Regular schemes}

Conventions for this section. Rings are commutative with $1$. Results of the text are cited as ``Prop.~2.5'', ``Cor.~2.12'', ``Thm.~2.16'' (Chapter 4, Section 2), and results of other chapters with the chapter number in front (``Cor.~2.5.24'' is Chapter 2, Section 5). For a local ring $(\mathcal O,\fm)$ with residue field $k$ we write $\operatorname{emb}(\mathcal O)=\dim_k\fm/\fm^2$; recall $\operatorname{emb}\ge\dim$ (Cor.~2.5.14(b)), with equality exactly when $\mathcal O$ is regular. For a field $k$ and a closed point $x$ of $\mathbb A^n_k$ with maximal ideal $\fm$, the map $D_x$ from $\fm$ to the dual of $k(x)^n$ is $D_xP:(t_i)\mapsto\sum_i\frac{\partial P}{\partial T_i}(x)t_i$, where $\frac{\partial P}{\partial T_i}(x)\in k(x)$ is the image of the partial derivative. From $D_x(PQ)=P(x)D_xQ+Q(x)D_xP$ one sees that $D_x$ kills $\fm^2$ and is $k(x)$-linear on $\fm/\fm^2$.

Honest notes. In Ex.~2.1 the last assertion (``$x$ is discrete in $X_y$'') is false as printed for a point $x$ that is not closed in $X_y$; it is proved for $x$ closed in $X_y$ (which is the case of interest), and what is true in general is also stated. In Ex.~2.4(a) the polynomials $P_i$ are those of Ex.~2.1.5, i.e.\ with $k[T_1,\dots,T_i]\cap\fm=(P_1,\dots,P_i)$; this property is used, and its existence is recalled. In Ex.~2.10 the ``$x_0$'' of the book is read as $x$. Ex.~2.11: the curve is the closed subscheme (possibly non-reduced) defined by $x^n+y^n+z^n$. Ex.~2.13 is proved fully (with $k$ infinite).

\exercise{2.1}{Let $f:X\to Y$ be a morphism of locally Noetherian schemes. Let $x\in X$, $y=f(x)$. Show that $T_{f,x}$ is injective if and only if the maximal ideal $\fm_x$ of $\mathcal O_{X,x}$ is generated by $f^\#_x(\fm_y)$. Show that if $f$ is moreover of finite type, then $x$ is discrete in $X_y$.}
\pf
\emph{First assertion.} By Def.~2.1, $T_{f,x}:T_{X,x}\to T_{Y,y}\otimes_{k(y)}k(x)$ is the $k(x)$-dual of the $k(x)$-linear map
\[ \psi:\ (\fm_y/\fm_y^2)\otimes_{k(y)}k(x)\longrightarrow \fm_x/\fm_x^2 \]
induced by $f^\#_x$ (which is a local homomorphism, so $f^\#_x(\fm_y)\subseteq\fm_x$). Both spaces are finite dimensional because $\fm_x,\fm_y$ are finitely generated ($X,Y$ locally Noetherian), so $T_{Y,y}\otimes k(x)=(\fm_y/\fm_y^2)^\vee\otimes k(x)=((\fm_y/\fm_y^2)\otimes k(x))^\vee$ and $T_{f,x}=\psi^\vee$. For linear maps of finite-dimensional vector spaces, $\psi^\vee$ is injective if and only if $\psi$ is surjective. The image of $\psi$ is the $k(x)$-span of the classes of the elements $f^\#_x(a)$, $a\in\fm_y$, that is $(f^\#_x(\fm_y)\mathcal O_{X,x}+\fm_x^2)/\fm_x^2$. Thus $T_{f,x}$ is injective if and only if $f^\#_x(\fm_y)\mathcal O_{X,x}+\fm_x^2=\fm_x$. Since $\fm_x$ is finitely generated, Nakayama's lemma shows that this holds if and only if $f^\#_x(\fm_y)\mathcal O_{X,x}=\fm_x$ (if $N=\fm_x/\fm_y\mathcal O_{X,x}$ then $N=\fm_xN$ forces $N=0$; the converse is trivial).

\emph{Second assertion.} The statement ``$x$ is discrete in $X_y$'' means that $x$ is an isolated point of the fibre $X_y$. As printed it cannot hold for every $x$: for $f:\mathbb A^1_k\to\operatorname{Spec}k$ and $x$ the generic point, $\fm_x=0$ is generated by $f^\#_x(\fm_y)=0$, but $x$ is not isolated in $X_y=\mathbb A^1_k$. What is true, and what we prove, is:

\emph{(i)} If $T_{f,x}$ is injective, then $x$ is the generic point of an irreducible component of $X_y$. \emph{(ii)} If moreover $f$ is of finite type and $x$ is closed in $X_y$ (equivalently, $k(x)$ is finite over $k(y)$), then $x$ is an isolated point of $X_y$.

The local ring of the fibre is $\mathcal O_{X_y,x}=\mathcal O_{X,x}\otimes_{\mathcal O_{Y,y}}k(y)=\mathcal O_{X,x}/\fm_y\mathcal O_{X,x}$. By the first part, if $T_{f,x}$ is injective this equals $\mathcal O_{X,x}/\fm_x=k(x)$, a field. Take an affine open neighbourhood $\operatorname{Spec}C$ of $x$ in $X_y$, with $x$ corresponding to the prime $\fp$. Then $C_\fp$ is a field, so its only prime is $0$, i.e.\ $\fp$ contains no smaller prime. So $\fp$ is a minimal prime of $C$, and $x$ is the generic point of an irreducible component of $\operatorname{Spec}C$, whose closure $Z$ in $X_y$ is an irreducible component of $X_y$ with generic point $x$. This proves (i).

For (ii), $X_y$ is of finite type over $k(y)$, hence a Noetherian space with finitely many irreducible components. If $x$ is closed in $X_y$ then $Z=\overline{\{x\}}=\{x\}$. Let $Z'$ be the union of the other irreducible components; it is closed. If $x\in Z'$, then $x$ lies in another component $Z_1$, so $Z=\{x\}\subseteq Z_1$ contradicts the maximality of $Z$. Hence $x\notin Z'$ and $\{x\}=X_y\setminus Z'$ is open in $X_y$: $x$ is isolated. \hfill$\square$

\exercise{2.2}{Let $k$ be a field. Show that $\operatorname{Spec}k[x,y,z]/(x^2-yz)$ is normal but singular.}
\pf
Let $A=k[x,y,z]/(x^2-yz)$ and $F=x^2-yz$. The argument works in every characteristic.

\emph{Step 1: $A$ is isomorphic to the subring $R=k[s^2,st,t^2]$ of $k[s,t]$.} Let $\theta:k[x,y,z]\to k[s,t]$ be given by $x\mapsto st$, $y\mapsto s^2$, $z\mapsto t^2$. Then $\theta(F)=0$, so $\theta$ induces $\bar\theta:A\to R$, which is surjective. Using $x^2=yz$ in $A$, the monomials $x^\epsilon y^iz^j$ with $\epsilon\in\{0,1\}$, $i,j\ge0$ span $A$ over $k$. Their images are $s^{2i}t^{2j}$ ($\epsilon=0$) and $s^{2i+1}t^{2j+1}$ ($\epsilon=1$), which are pairwise distinct monomials of $k[s,t]$, hence linearly independent. Therefore these monomials form a basis of $A$ and $\bar\theta$ is injective. So $A\simeq R$. In particular $A$ is an integral domain. Moreover $R$ is the set of polynomials all of whose homogeneous components have even degree: a monomial $s^it^j$ with $i+j$ even has $i,j$ of the same parity, so it is a product of $s^2,st,t^2$.

\emph{Step 2: $A$ is normal.} Let $\alpha\in\operatorname{Frac}(A)\subseteq k(s,t)$ be integral over $A$. Then $\alpha$ is integral over the UFD $k[s,t]$, hence $\alpha\in k[s,t]$. Write $\alpha=a/b$ with $a,b\in A$, $b\ne0$, so $a=b\alpha$. Split $\alpha=\alpha'+\alpha''$ into its even-degree and odd-degree parts. Since $a$ and $b$ have only even-degree components, $b\alpha''$ has only odd-degree components while $a-b\alpha'$ has only even-degree components, and they are equal; so $b\alpha''=0$, whence $\alpha''=0$ because $k[s,t]$ is a domain. Thus $\alpha\in R=A$. So $A$ is an integrally closed domain; all its localizations are integrally closed domains, and $X=\operatorname{Spec}A$ is normal.

\emph{Step 3: $X$ is singular.} The fraction field of $R$ is $k(s^2,st,t^2)$; it is algebraic over $k(s^2,t^2)$ because $(st)^2=s^2t^2$, and $s^2,t^2$ are algebraically independent. So $\operatorname{trdeg}_k=2$ and $\dim X=2$ (Prop.~2.5.19). Let $x_0$ be the rational point with maximal ideal $\fm=(x,y,z)A$; it is a closed point, so $\dim\mathcal O_{X,x_0}=2$ (Cor.~2.5.24). By Lemma~2.3,
\[ \fm\mathcal O_{X,x_0}/\fm^2\mathcal O_{X,x_0}\simeq\fm/\fm^2=(x,y,z)/\bigl((x,y,z)^2+(F)\bigr)=(x,y,z)/(x,y,z)^2, \]
because $F\in(x,y,z)^2$; this has dimension $3>2=\dim\mathcal O_{X,x_0}$. Hence $\mathcal O_{X,x_0}$ is not regular and $X$ is singular (at $x_0$; in characteristic $\neq2$ this is the only singular point, but we do not need it). \hfill$\square$

\exercise{2.3}{Let $\mathcal O_K$ be a discrete valuation ring with uniformizing parameter $t$ and field of fractions $K$. Let $n\ge1$. We suppose that $\operatorname{char}(K)\ne2,3$. (a) Show that $K[x,y]/(y^2+x^3+t^n)$ is regular. (b) Let $A=\mathcal O_K[x,y]/(y^2+x^3+t^n)$. Show that the special fibre of $\operatorname{Spec}A\to\operatorname{Spec}\mathcal O_K$ is reduced. Deduce from this that $A$ is normal (see Lemma~1.18). (c) Let $\fm$ be the maximal ideal of $A$ generated by $x$, $y$ and $t$. Show that $A_\fm$ is regular if and only if $n=1$.}
\pf
Put $F=y^2+x^3+t^n$ and let $k=\mathcal O_K/(t)$ be the residue field.

\emph{(a)} The ring $K[x,y]/(F)$ is Noetherian, so by Cor.~2.17 it suffices to prove regularity at the closed points, i.e.\ at the images of the maximal ideals $\fm$ of $K[x,y]$ containing $F$. Fix such an $\fm$. The ring $S=K[x,y]_\fm$ is a regular local ring of dimension $2$ (Example~2.18), hence a domain (Prop.~2.11), and $F\in\fm S$ is nonzero. The local ring of $\operatorname{Spec}K[x,y]/(F)$ at the point is $S/FS$, which by Cor.~2.12 is regular if and only if $F\notin\fm^2S$. Since $\fm\subset K[x,y]$ is the preimage of $\fm S$ and $\fm S\cap K[x,y]=\fm$, $\fm^2S\cap K[x,y]=\fm^2$, it suffices to show $F\notin\fm^2$.

Suppose $F\in\fm^2$. The partial derivatives $\partial_x,\partial_y$ are $K$-linear derivations of $K[x,y]$, and a derivation $D$ maps $\fm^2$ into $\fm$ (as $D(ab)=aDb+bDa$). Hence $\partial_xF=3x^2\in\fm$ and $\partial_yF=2y\in\fm$. Since $2,3\in K^*$ (here $\operatorname{char}K\ne2,3$ is used), $x^2\in\fm$ and $y\in\fm$, so $x\in\fm$ ($\fm$ is prime). Then $t^n=F-y^2-x^3\in\fm$. But $t^n$ is a unit of $K$, so $\fm$ would be the unit ideal, a contradiction. Thus $F\notin\fm^2$, and $K[x,y]/(F)$ is regular.

\emph{(b)} Since $t^n=0$ in $k$, the special fibre is $\operatorname{Spec}k[x,y]/(y^2+x^3)$. The polynomial $y^2+x^3\in k[x][y]$ is monic of degree $2$ in $y$, so it is reducible only if it has a root $g\in k[x]$, i.e.\ $g^2=-x^3$. This is impossible by comparing degrees ($2\deg g$ is even and $3$ is odd; and $g=0$ gives $x^3=0$). So $y^2+x^3$ is irreducible, hence prime in the UFD $k[x,y]$, and $k[x,y]/(y^2+x^3)$ is a domain; in particular the special fibre is reduced. (This holds whatever the characteristic of $k$.)

To apply Lemma~1.18 we check the hypotheses. As an $\mathcal O_K[x]$-module, $A=\mathcal O_K[x][y]/(y^2+x^3+t^n)$ is free with basis $1,y$ (the relation is monic in $y$), and $\mathcal O_K[x]$ is free over $\mathcal O_K$; so $A$ is free, hence flat, over $\mathcal O_K$. Flatness is local, so $\mathcal O_X(U)$ is flat over $\mathcal O_K$ for every affine open $U\subseteq X=\operatorname{Spec}A$ (all its localizations are localizations of $A$). The generic fibre is $X_K=\operatorname{Spec}(A\otimes_{\mathcal O_K}K)=\operatorname{Spec}K[x,y]/(F)$, which is regular by (a), hence normal (Thm.~2.16(b): regular local rings are factorial, hence normal). The special fibre $X_k$ is reduced. By Lemma~1.18, $X$ is normal, i.e.\ $A$ is normal.

\emph{(c)} Let $M=(t,x,y)$, a maximal ideal of $B=\mathcal O_K[x,y]$ with $B/M=k$, and $F\in M$. The ideal $\fm$ of $A$ is $M/(F)$, and $A_\fm=S/FS$ with $S=B_M$. Since $M\cap\mathcal O_K=(t)$ has height $1$, Lemma~2.5.16 gives $\operatorname{ht}M=1+2=3$, so $\dim S=3$; and $MS$ is generated by the three elements $t,x,y$. Hence $S$ is regular of dimension $3$, a domain (Prop.~2.11), and $t,x,y$ form a basis of the $k$-vector space $MS/M^2S$ (three generators of a $3$-dimensional space). By Cor.~2.12, $A_\fm=S/FS$ is regular if and only if $F\notin M^2S$.

If $n\ge2$, then $t^n,\,x^3,\,y^2\in M^2S$, so $F\in M^2S$ and $A_\fm$ is not regular. If $n=1$, then $F\equiv t\pmod{M^2S}$ because $y^2+x^3\in M^2S$; the class of $t$ is a basis vector of $MS/M^2S$, hence nonzero, so $F\notin M^2S$ and $A_\fm$ is regular. Therefore $A_\fm$ is regular if and only if $n=1$. \hfill$\square$

\exercise{2.4}{Let $k$ be a field and let $\fm$ be a maximal ideal of $k[T_1,\dots,T_n]$ corresponding to a closed point $x$ of $\mathbb A^n_k$. Let us consider the $k(x)$-linear map $\phi_x:\fm/\fm^2\to E_x^\vee$ defined as before Lemma~2.4, with $E_x=k(x)^n$. (a) Let $P_1(T_1),P_2(T_1,T_2),\dots,P_n(T_1,\dots,T_n)$ be a system of generators of $\fm$ (see Exercise 2.1.5). Show that $\phi_x$ is an isomorphism if and only if $(\partial P_i/\partial T_i)(x)\ne0$ for all $i\le n$. (b) Deduce from (a) that $\phi_x$ is an isomorphism if and only if $k(x)$ is separable over $k$.}
\pf
Let $x_i\in k(x)=k[T]/\fm$ be the image of $T_i$, $\fm_i=\fm\cap k[T_1,\dots,T_i]$ and $k_i=k[x_1,\dots,x_i]\subseteq k(x)$. Then $k[T_1,\dots,T_i]/\fm_i\simeq k_i$ is a domain which is finite-dimensional over $k$ (as $k(x)$ is finite over $k$ by Cor.~2.1.12), hence a field: $k_i=k(x_1,\dots,x_i)$, $k_0=k$, $k_n=k(x)$.

In (a) the $P_i$ are polynomials with $P_i\in k[T_1,\dots,T_i]$ generating $\fm$; the proof of (a) uses only this. In (b) we choose them as in Exercise 2.1.5, i.e.\ with $\fm_i=(P_1,\dots,P_i)k[T_1,\dots,T_i]$ (this is harmless: by (a) the criterion holds for every such system, so (b) may use a convenient one). (Existence, for completeness: suppose $P_1,\dots,P_{i-1}$ are found. Then $k[T_1,\dots,T_i]/(P_1,\dots,P_{i-1})\simeq k_{i-1}[T_i]$, a principal ideal ring, and the image of $\fm_i$ is generated by one element; lifting it to $P_i\in k[T_1,\dots,T_i]$ gives $\fm_i=(P_1,\dots,P_i)$.)

\emph{(a)} By Example~2.18 (regularity of $\mathbb A^n_k$ at closed points, Cor.~2.5.24 for the dimension $n$), $\dim_{k(x)}\fm/\fm^2=n$. The classes of $P_1,\dots,P_n$ span $\fm/\fm^2$, hence form a basis. The map $\phi_x$ sends the class of $P_i$ to the $i$-th row of the $n\times n$ matrix $\bigl(\frac{\partial P_i}{\partial T_j}(x)\bigr)_{i,j}$ over $k(x)$. Since $P_i\in k[T_1,\dots,T_i]$, $\frac{\partial P_i}{\partial T_j}=0$ for $j>i$, so the matrix is lower triangular. Now $\phi_x$ is an isomorphism iff it maps a basis to a basis iff the matrix is invertible iff its determinant, the product of the diagonal entries $\frac{\partial P_i}{\partial T_i}(x)$, is nonzero iff all diagonal entries are nonzero.

\emph{(b)} Let $\bar P_i\in k_{i-1}[T_i]$ be the image of $P_i$ under reduction of the coefficients $k[T_1,\dots,T_{i-1}]\to k_{i-1}$. The surjection $k[T_1,\dots,T_i]\to k_{i-1}[T_i]\to k_i$ ($T_i\mapsto x_i$) has kernel $\fm_i=(P_1,\dots,P_i)$, and the first map has kernel $\fm_{i-1}k[T_1,\dots,T_i]=(P_1,\dots,P_{i-1})$. Hence the kernel of $k_{i-1}[T_i]\to k_i$ is generated by $\bar P_i$. This kernel is nonzero, so $\bar P_i=c\,\mu_i$ with $c\in k_{i-1}^*$ and $\mu_i$ the minimal polynomial of $x_i$ over $k_{i-1}$. Since reducing coefficients commutes with $\partial/\partial T_i$, we get $\frac{\partial P_i}{\partial T_i}(x)=\bar P_i'(x_i)=c\,\mu_i'(x_i)$. As $x_i$ is a root of $\mu_i$, it is a simple root iff $\mu_i'(x_i)\ne0$, i.e.\ iff $x_i$ is separable over $k_{i-1}$. By (a), $\phi_x$ is an isomorphism iff $x_i$ is separable over $k_{i-1}$ for all $i$.

If every $x_i$ is separable over $k_{i-1}$, then $k_n=k(x)$ is separable over $k$, since a tower of finite separable extensions is separable. Conversely if $k(x)/k$ is separable, each $x_i$ is separable over $k$, and its minimal polynomial over $k_{i-1}$ divides its (separable) minimal polynomial over $k$, so it is separable over $k_{i-1}$. Hence $\phi_x$ is an isomorphism iff $k(x)/k$ is separable. \hfill$\square$

\exercise{2.5}{Show that the Jacobian criterion (with $J_x$ as a matrix with coefficients in $k(x)$) remains valid at the closed points $x\in X$ for which $k(x)$ is separable over $k$. See also Proposition~3.30. Give an example with $X$ regular at $x$ and $\dim\mathcal O_{X,x}\ne n-\operatorname{rank}J_x$.}
\pf
\emph{Statement.} Let $X=V(I)\subseteq\mathbb A^n_k$ be a closed subscheme, $F_1,\dots,F_r$ generators of $I$, $x\in X$ a closed point with $k(x)$ separable over $k$, and $J_x=\bigl(\frac{\partial F_i}{\partial T_j}(x)\bigr)\in M_{r\times n}(k(x))$. Then $X$ is regular at $x$ iff $\operatorname{rank}J_x=n-\dim\mathcal O_{X,x}$.

\emph{Proof.} Let $\fm\subset k[T]$ be the maximal ideal of $x$ and $\fn=\fm/I$ the maximal ideal of $A=k[T]/I$ at $x$. Then $\fn/\fn^2=\fm/(I+\fm^2)$, and by Lemma~2.3 this is $\fm_x/\fm_x^2$ for $\fm_x$ the maximal ideal of $\mathcal O_{X,x}$. So $X$ is regular at $x$ iff $\dim_{k(x)}\fm/(I+\fm^2)=\dim\mathcal O_{X,x}$.

Since $k(x)/k$ is separable, $\phi_x:\fm/\fm^2\to k(x)^{n\vee}$ is an isomorphism (Ex.~2.4(b)). The subspace $(I+\fm^2)/\fm^2\subseteq\fm/\fm^2$ is the $k(x)$-span of the classes of the elements of $I$. If $G=\sum_iA_iF_i\in I$ ($A_i\in k[T]$), then the class of $G$ in $\fm/\fm^2$ equals $\sum_iA_i(x)\,[F_i]$ because $F_i\in\fm$ (the class of $aF$ is $a(x)[F]$ for $F\in\fm$). So $(I+\fm^2)/\fm^2$ is the span of $[F_1],\dots,[F_r]$, and $\phi_x$ maps it onto the span of the rows of $J_x$, which has dimension $\operatorname{rank}J_x$. As $\phi_x$ is bijective,
\[ \dim_{k(x)}\fm_x/\fm_x^2=n-\dim_{k(x)}(I+\fm^2)/\fm^2=n-\operatorname{rank}J_x. \]
Thus $X$ is regular at $x$ iff $n-\operatorname{rank}J_x=\dim\mathcal O_{X,x}$. (When $x$ is rational this is Thm.~2.19.) \hfill$\square$

\emph{Example where separability fails.} Let $p$ be a prime, $k=\mathbb F_p(s)$, $n=1$ and $X=\operatorname{Spec}k[T]/(T^p-s)$ (the polynomial $T^p-s$ is irreducible, by Eisenstein at the prime $s$ of $\mathbb F_p[s]$). Then $k[T]/(T^p-s)=k(s^{1/p})$ is a field, so $X$ is a single point $x$ with $\mathcal O_{X,x}$ a field: $X$ is regular and $\dim\mathcal O_{X,x}=0$. But $J_x=(pT^{p-1})=(0)$ has rank $0$, so $n-\operatorname{rank}J_x=1\ne0$. (Here $k(x)=k(s^{1/p})$ is purely inseparable over $k$.)

\exercise{2.6}{Let $k$ be a field. Let $F_1,\dots,F_n\in k[T_1,\dots,T_n]$. Let us consider the homomorphism $\phi:k[T_1,\dots,T_n]\to k[T_1,\dots,T_n]$ defined by $\phi(T_i)=F_i(T)$ and let us suppose that the associated morphism $\mathbb A^n_k\to\mathbb A^n_k$ is an isomorphism. Show that $\det\bigl(\frac{\partial F_i}{\partial T_j}\bigr)_{1\le i,j\le n}$ belongs to $k^*$.}
\pf
A morphism of affine schemes is an isomorphism iff the ring homomorphism is (Prop.~2.3.25), so $\phi$ is a $k$-algebra automorphism of $k[T]=k[T_1,\dots,T_n]$. Let $\psi=\phi^{-1}$ and $G_i=\psi(T_i)\in k[T]$. Then $T_i=\phi(\psi(T_i))=\phi(G_i)=G_i(F_1,\dots,F_n)$.

The chain rule holds for composition of polynomials over any ring: for $G\in k[T]$, $\frac{\partial}{\partial T_k}\bigl(G(F_1,\dots,F_n)\bigr)=\sum_j\frac{\partial G}{\partial T_j}(F_1,\dots,F_n)\,\frac{\partial F_j}{\partial T_k}$. (Both sides are $k$-linear in $G$, and for a monomial $G=T_1^{e_1}\cdots T_n^{e_n}$ this is the Leibniz rule.) Applying it to $G=G_i$ and using $G_i(F)=T_i$,
\[ \delta_{ik}=\frac{\partial T_i}{\partial T_k}=\sum_j\frac{\partial G_i}{\partial T_j}(F)\,\frac{\partial F_j}{\partial T_k}. \]
In matrix form, with entries in $k[T]$: $\bigl(\frac{\partial G_i}{\partial T_j}(F)\bigr)_{ij}\cdot\bigl(\frac{\partial F_j}{\partial T_k}\bigr)_{jk}=\operatorname{Id}$. Taking determinants, $\det\bigl(\frac{\partial G_i}{\partial T_j}(F)\bigr)\cdot\det\bigl(\frac{\partial F_i}{\partial T_j}\bigr)=1$ in $k[T]$. So $\det(\partial F_i/\partial T_j)$ is a unit of the polynomial ring $k[T]$, i.e.\ an element of $k^*$. \hfill$\square$

\exercise{2.7}{Let $X$ be a scheme over a field $k$, and let $x\in X(k)$. Show that there exists a canonical bijection between $T_{X,x}$ and the set of $k$-morphisms $\operatorname{Spec}(k[\epsilon]/\epsilon^2)\to X$ with image equal to $x$.}
\pf
Let $D=k[\epsilon]/(\epsilon^2)$; $\operatorname{Spec}D$ has one point, corresponding to the maximal ideal $\epsilon D$. Let $\mathcal O=\mathcal O_{X,x}$, with maximal ideal $\fm$. Since $x\in X(k)$, $k(x)=k$, and the structure $k\to\mathcal O\to\mathcal O/\fm=k$ is the identity; so $\mathcal O=k\oplus\fm$ and we write $\pi:\mathcal O\to k$ for the residue map.

\emph{Step 1: morphisms with image $x$ are local $k$-algebra homomorphisms $\mathcal O\to D$.} Let $g:\operatorname{Spec}D\to X$ be a $k$-morphism with image $\{x\}$. Take an affine open neighbourhood $U=\operatorname{Spec}C$ of $x$; $g$ factors through $U$, so by Prop.~2.3.25 it corresponds to a $k$-algebra homomorphism $h:C\to D$, and ``image $=x$'' means $h^{-1}(\epsilon D)=\fp_x$. Then $h$ maps $C\setminus\fp_x$ into $D^*$ (the elements outside the maximal ideal), so it extends uniquely to $\bar h:C_{\fp_x}=\mathcal O\to D$, and $\bar h(\fm)\subseteq\epsilon D$ (local). Conversely a local $k$-homomorphism $\mathcal O\to D$ composed with $C\to\mathcal O$ gives such an $h$, and the two constructions are inverse to each other. They do not depend on $U$.

\emph{Step 2: local $k$-algebra homomorphisms $\mathcal O\to D$ correspond to $(\fm/\fm^2)^\vee$.} Let $h:\mathcal O\to D$ be $k$-linear, write $h(a)=\lambda(a)+\delta(a)\epsilon$ with $\lambda,\delta:\mathcal O\to k$ $k$-linear. Then $h$ is a local $k$-algebra homomorphism iff $\lambda$ is a $k$-algebra homomorphism with $\lambda(\fm)=0$, and $h(ab)=h(a)h(b)$ holds in the $\epsilon$-part. The first condition forces $\lambda=\pi$ (a $k$-algebra homomorphism $\mathcal O\to k$ killing $\fm$ factors through $\mathcal O/\fm=k$ and is the identity there). The multiplicativity of $h$ then reads
\[ \delta(ab)=\pi(a)\delta(b)+\pi(b)\delta(a)\qquad(a,b\in\mathcal O), \]
i.e.\ $\delta$ is a $k$-derivation $\mathcal O\to k$ (with $k$ an $\mathcal O$-module via $\pi$); in particular $\delta(1)=0$. So local $k$-homomorphisms $\mathcal O\to D$ correspond bijectively to $\operatorname{Der}_k(\mathcal O,k)$.

Now $\operatorname{Der}_k(\mathcal O,k)\simeq(\fm/\fm^2)^\vee$ via $\delta\mapsto\delta|_{\fm}$: if $a,b\in\fm$ then $\delta(ab)=0$, so $\delta|_\fm$ factors through $\fm/\fm^2$. Conversely, given a $k$-linear $\ell:\fm/\fm^2\to k$ define $\delta_\ell(a)=\ell\bigl(\overline{a-\pi(a)}\bigr)$ (note $a-\pi(a)\in\fm$). Writing $a=\lambda+u$, $b=\mu+v$ with $\lambda,\mu\in k$, $u,v\in\fm$, we get $ab=\lambda\mu+(\lambda v+\mu u)+uv$ with $uv\in\fm^2$, so $\delta_\ell(ab)=\lambda\ell(v)+\mu\ell(u)=\pi(a)\delta_\ell(b)+\pi(b)\delta_\ell(a)$: $\delta_\ell$ is a derivation. The two maps are mutually inverse, because a derivation vanishes on $k\cdot1$ and so is determined by its restriction to $\fm$.

Composing Steps 1 and 2 we get a bijection between the set of $k$-morphisms $\operatorname{Spec}D\to X$ with image $x$ and $(\fm/\fm^2)^\vee=T_{X,x}$ (Def.~2.1); it involves no choices, so it is canonical. Explicitly, $\ell$ corresponds to the morphism given by $a\mapsto\pi(a)+\ell(\overline{a-\pi(a)})\,\epsilon$. \hfill$\square$

\exercise{2.8}{Let us once again consider the scheme $\operatorname{Proj}B$ of Exercise 2.3.13 (so $B=A[X,Y,Z]$ with $B_d$ the sub-$A$-module generated by the monomials $X^aY^bZ^c$ with $a+2b+3c=d$). Show that the fibres of $\operatorname{Proj}B\to\operatorname{Spec}A$ are all singular.}
\pf
Here $X,Y,Z$ have weights $1,2,3$, and $Z\in B_3\subseteq B_+$ is homogeneous, so $D_+(Z)=\operatorname{Spec}B_{(Z)}$ is an affine open subscheme of $\operatorname{Proj}B$ (Prop.~2.3.38). The ring $B_{(Z)}$ is the degree-$0$ part of $B_Z=A[X,Y,Z,Z^{-1}]$, i.e.\ the free $A$-module with basis the Laurent monomials $X^aY^bZ^c$ ($a,b\ge0$, $c\in\mathbb Z$) with $a+2b+3c=0$.

\emph{Claim: $B_{(Z)}\simeq A[U,V,W]/(UW-V^3)$ via $U\mapsto X^3/Z$, $V\mapsto XY/Z$, $W\mapsto Y^3/Z^2$.} These three elements lie in $B_{(Z)}$ (their weights are $0$), and $UW-V^3\mapsto X^3Y^3/Z^3-X^3Y^3/Z^3=0$, so we get an $A$-algebra homomorphism $\bar\theta:Q:=A[U,V,W]/(UW-V^3)\to B_{(Z)}$. We have $U^iV^jW^l\mapsto X^{3i+j}Y^{j+3l}Z^{-i-j-2l}$.

\emph{Surjective.} Let $X^aY^bZ^c$ be a basis monomial, so $a+2b=-3c\equiv0\pmod3$, whence $a\equiv b\pmod 3$. If $a\ge b$, write $a-b=3i$; then $c=-(a+2b)/3=-i-b$ and the monomial is the image of $U^iV^b$. If $b>a$, write $b-a=3l$; then $c=-a-2l$ and the monomial is the image of $V^aW^l$.

\emph{Injective.} In $Q$ the relation $UW=V^3$ lets one reduce every monomial, so $Q$ is spanned over $A$ by $\mathcal S=\{U^iV^j\}_{i,j\ge0}\cup\{V^jW^l\}_{j\ge0,\,l\ge1}$. The images of the elements of $\mathcal S$ are $X^aY^bZ^c$ with $(a,b)=(3i+j,j)$ (then $a\ge b$) or $(a,b)=(j,j+3l)$ (then $b>a$); the pair $(a,b)$ determines the element of $\mathcal S$ and the monomial, so distinct elements of $\mathcal S$ go to distinct monomials, which are $A$-linearly independent in $B_Z$. Hence $\mathcal S$ is an $A$-basis of $Q$ and $\bar\theta$ is an isomorphism.

\emph{The fibres.} Let $\mathfrak p\in\operatorname{Spec}A$, $\kappa=k(\mathfrak p)$. The fibre $X_{\mathfrak p}=\operatorname{Proj}B\times_{\operatorname{Spec}A}\operatorname{Spec}\kappa$ contains the open subscheme $D_+(Z)\times_A\kappa=\operatorname{Spec}R$ with $R=B_{(Z)}\otimes_A\kappa=\kappa[U,V,W]/(UW-V^3)$. (Applying the claim with $A=\kappa$ shows that $R$ is isomorphic to a subring of the domain $\kappa[X,Y,Z,Z^{-1}]$, so $R$ is a domain.) Let $o$ be the rational point of $\operatorname{Spec}R$ with maximal ideal $M=(U,V,W)$.

By Lemma~2.3, $\fm_o/\fm_o^2=M/M^2$, and $M/M^2=(U,V,W)/\bigl((U,V,W)^2+(UW-V^3)\bigr)=(U,V,W)/(U,V,W)^2$, since $UW-V^3\in(U,V,W)^2$; so $\dim_\kappa\fm_o/\fm_o^2=3$. On the other hand the domain $R$ is generated over $\kappa$ by $u,v,w$ (the images of $U,V,W$) with $v^3=uw$, and $u\ne0$, so $\operatorname{Frac}(R)=\kappa(u,v)$ has transcendence degree $\le2$, whence $\dim\mathcal O_{X_{\mathfrak p},o}\le\dim R\le2$ (Prop.~2.5.19). Thus $\dim\mathcal O_{X_{\mathfrak p},o}\le2<3=\dim\fm_o/\fm_o^2$, and $\mathcal O_{X_{\mathfrak p},o}$ is not regular. So every fibre $X_{\mathfrak p}$ is singular (at the point $[0:0:1]$). \hfill$\square$

\exercise{2.9}{Let $X=V_+(I)$ be a closed subvariety of $\mathbb P^n_k$ over a field $k$. Let $x\in X(k)$, and let us consider the linear subvariety of $\mathbb P(k^{n+1})$
\[ PT_{X,x}=\Bigl\{(t_0,\dots,t_n)\in\mathbb P(k^{n+1})\ \Bigm|\ \sum_{0\le i\le n}\frac{\partial F}{\partial T_i}(x)\,t_i=0,\ \forall F(T_0,\dots,T_n)\in I\Bigr\}. \]
We will abusively call this linear subvariety the tangent space to $X$ at $x$. Show that $\dim_kT_{X,x}=\dim PT_{X,x}$.}
\pf
Here $I\subseteq k[T_0,\dots,T_n]$ is a homogeneous ideal and $X=\operatorname{Proj}k[T]/I$, and in the formula $x$ stands for a fixed representative $(x_0,\dots,x_n)\in k^{n+1}\setminus\{0\}$ of the point (Lemma~2.3.43). Rescaling $x$ by $\lambda\in k^*$ multiplies the partial derivatives of a homogeneous $F$ of degree $d$ by $\lambda^{d-1}$, so the set $PT_{X,x}$ does not depend on the representative. For $F\in I$ put $L_F(t)=\sum_j\frac{\partial F}{\partial T_j}(x)t_j$, a linear form on $k^{n+1}$, and let
\[ W=\{t\in k^{n+1}:\ L_F(t)=0\ \ \forall F\in I\}. \]
Then $PT_{X,x}=\mathbb P(W)$ has dimension $\dim_kW-1$. As $I$ is homogeneous and $F\mapsto L_F$ is linear, $W$ is cut out by the $L_F$ with $F\in I$ homogeneous.

Choose $i$ with $x_i\ne0$; scaling, we may assume $x_i=1$. Put $S_j=T_j/T_i$ ($j\ne i$); then $D_+(T_i)\simeq\mathbb A^n_k$ with coordinates $S_j$, and $x$ is the rational point $x^a=(x_j)_{j\ne i}$. By Prop.~2.3.38 and Lemma~2.3.36(d),(e), $X\cap D_+(T_i)$ is the closed subscheme of $\mathbb A^n_k$ defined by the ideal $I^a$ generated by the dehomogenizations $F^a(S)=F(S_0,\dots,1,\dots,S_n)$ ($1$ in the $i$-th place) of homogeneous $F\in I$; and $T_{X,x}$ is the tangent space of this affine scheme at $x^a$. By Prop.~2.5,
\[ T_{X,x}\simeq T':=\Bigl\{s\in k^n\ \Bigm|\ \sum_{j\ne i}\frac{\partial F^a}{\partial S_j}(x^a)\,s_j=0\ \ \forall F^a\in I^a\Bigr\}. \]
For $G=\sum_mA_mF_m^a\in I^a$ ($F_m\in I$ homogeneous), $D_{x^a}G=\sum_mA_m(x^a)D_{x^a}F^a_m$ since $F_m^a(x^a)=F_m(x)=0$ ($x\in X$). So it is enough to impose the conditions for $F^a$ with $F\in I$ homogeneous. Moreover $\frac{\partial F^a}{\partial S_j}(x^a)=\frac{\partial F}{\partial T_j}(x)$ for $j\ne i$, because setting $T_i=1$ commutes with $\partial/\partial T_j$, $j\ne i$.

For homogeneous $F\in I$ of degree $d$, Euler's identity $\sum_jT_j\frac{\partial F}{\partial T_j}=dF$ (valid over any ring) at the point $x$ gives $\sum_jx_j\frac{\partial F}{\partial T_j}(x)=dF(x)=0$, i.e.\ $\frac{\partial F}{\partial T_i}(x)=-\sum_{j\ne i}x_j\frac{\partial F}{\partial T_j}(x)$ (using $x_i=1$). Hence for all $t\in k^{n+1}$
\[ L_F(t)=\sum_{j\ne i}\frac{\partial F}{\partial T_j}(x)\,(t_j-x_jt_i). \tag{$*$} \]
Define $\rho:k^{n+1}\to k^n$, $\rho(t)=(t_j-x_jt_i)_{j\ne i}$. By $(*)$ and the description of $T'$, for homogeneous $F\in I$: $L_F(t)=0$ iff $\rho(t)$ satisfies the condition of $F$ in $T'$. So $W=\rho^{-1}(T')$. The map $\rho$ is surjective ($\rho(t)=s$ for $t_i=0$, $t_j=s_j$), with kernel $\{t:\ t_j=x_jt_i\ \forall j\ne i\}=k\cdot x$ (recall $x_i=1$), a line. Hence $\dim_kW=\dim_kT'+1$ and
\[ \dim PT_{X,x}=\dim_kW-1=\dim_kT'=\dim_kT_{X,x}. \]
This proves the assertion. \hfill$\square$

\exercise{2.10}{Let $X$ be a closed subvariety of $\mathbb P^n_k=\operatorname{Proj}k[T_0,\dots,T_n]$ defined by some homogeneous polynomials $F_1,\dots,F_r$. Let $x\in X(k)$. Let $J_x$ denote the matrix $J_x=\bigl(\frac{\partial F_i}{\partial T_j}(x)\bigr)_{1\le i\le r,\ 0\le j\le n}$. Show that $x$ is regular in $X$ if and only if $\operatorname{rank}J_x=n-\dim\mathcal O_{X,x}$. This is the Jacobian criterion for projective varieties.}
\pf
(The point called $x_0$ in the statement is $x$.) Here $I=(F_1,\dots,F_r)$, and $x$ is represented by a fixed vector $x\in k^{n+1}\setminus\{0\}$; rescaling it multiplies the $i$-th row of $J_x$ by $\lambda^{\deg F_i-1}\neq 0$, so the rank of $J_x$ does not depend on the representative.

We use the notation $L_F$, $W$ of the proof of Ex.~2.9. If $F=\sum_iA_iF_i\in I$, then $\frac{\partial F}{\partial T_j}(x)=\sum_iA_i(x)\frac{\partial F_i}{\partial T_j}(x)+\sum_iF_i(x)\frac{\partial A_i}{\partial T_j}(x)=\sum_iA_i(x)\frac{\partial F_i}{\partial T_j}(x)$, since $F_i(x)=0$. So $L_F$ is a combination of $L_{F_1},\dots,L_{F_r}$, and $W=\{t\in k^{n+1}:\ J_xt=0\}$ has dimension $n+1-\operatorname{rank}J_x$. By the proof of Ex.~2.9, $\dim_kT_{X,x}=\dim_kW-1=n-\operatorname{rank}J_x$.

By Def.~2.8, $X$ is regular at $x$ iff $\dim_kT_{X,x}=\dim\mathcal O_{X,x}$ (here $k(x)=k$), i.e.\ iff $n-\operatorname{rank}J_x=\dim\mathcal O_{X,x}$. \hfill$\square$

\exercise{2.11}{Let $k$ be an algebraically closed field, and let $n\ge1$. Show that the curve $\operatorname{Proj}k[x,y,z]/(x^n+y^n+z^n)$ is regular if and only if $n$ is prime to $\operatorname{char}(k)$.}
\pf
Let $F=x^n+y^n+z^n$ and $X=\operatorname{Proj}k[x,y,z]/(F)\subseteq\mathbb P^2_k$ (as a scheme, possibly non-reduced). The scheme $X$ is Noetherian, so it is regular iff it is regular at its closed points (Cor.~2.17). Since $k$ is algebraically closed, every closed point $x$ of $X$ is rational (its residue field is finite over $k$, Cor.~2.1.12). Conversely every rational point of $X$ is closed.

\emph{Local dimension.} A closed point $x$ lies in a chart $D_+(x_i)\cap X=\operatorname{Spec}k[u,v]/(f)$, where $f$ is the dehomogenization of $F$, a non-constant polynomial ($n\ge1$). Then $\mathcal O_{X,x}=S/fS$ with $S=k[u,v]_M$ for a maximal ideal $M\ni f$. The ring $S$ is a regular local ring of dimension $2$ (Example~2.18, Cor.~2.5.24), hence a domain (Prop.~2.11), and $f\ne0$ in $\mathfrak m_S$, so $\dim\mathcal O_{X,x}=\dim S-1=1$ (Thm.~2.5.15, as in the proof of Cor.~2.12). By Ex.~2.10 (with $n=2$, $r=1$), $X$ is regular at $x$ iff $\operatorname{rank}J_x=2-1=1$, i.e.\ iff $J_x\ne0$, where
\[ J_x=\bigl(n\,a^{n-1},\ n\,b^{n-1},\ n\,c^{n-1}\bigr),\qquad x=(a:b:c). \]

\emph{Case $\operatorname{char}(k)\nmid n$ (including $\operatorname{char}k=0$).} Then $n\in k^*$. If $n=1$, $J_x=(1,1,1)\ne0$. If $n\ge2$, $J_x=0$ would force $a^{n-1}=b^{n-1}=c^{n-1}=0$, i.e.\ $a=b=c=0$, impossible for a projective point. So $J_x\ne0$ at every closed point and $X$ is regular.

\emph{Case $p=\operatorname{char}(k)>0$ divides $n$.} Then $n=0$ in $k$, so $J_x=0$ at every point. The scheme $X$ has a rational point: since $k$ is algebraically closed, there is $a\in k$ with $a^n=-1$, and $(a:0:1)\in X$; it is a closed point. At this point $\operatorname{rank}J_x=0\ne1$, so $X$ is not regular there. Hence $X$ is regular if and only if $n$ is prime to $\operatorname{char}(k)$. \hfill$\square$

\exercise{2.12}{Let $X$, $Y$ be algebraic varieties over an algebraically closed field $k$. If $X$, $Y$ are regular, show that $X\times_kY$ is also regular. Show that this is false in general if $k$ is not algebraically closed.}
\pf
\emph{Regularity of the product.} Let $Z=X\times_kY$, with projections $p:Z\to X$, $q:Z\to Y$. It is of finite type over $k$, hence Noetherian, so by Cor.~2.17 it suffices to show that $Z$ is regular at its closed points. Let $z\in Z$ be a closed point; $k(z)$ is finite over $k$, hence equal to $k$ ($k$ algebraically closed). Put $x=p(z)$, $y=q(z)$; then $k\subseteq k(x)\subseteq k(z)=k$, so $k(x)=k=k(y)$: $x,y$ are rational points. Let $A=\mathcal O_{X,x}$, $B=\mathcal O_{Y,y}$, $d=\dim A$, $e=\dim B$ and $C=\mathcal O_{Z,z}$. By hypothesis $A$, $B$ are regular, so there are $f_1,\dots,f_d$ generating $\fm_x=\fm_A$ and $g_1,\dots,g_e$ generating $\fm_B$.

\emph{(1) $\dim C=d+e$.} The morphism $p:Z\to X$ is the base change of $Y\to\operatorname{Spec}k$, which is flat ($k$ is a field), so $p$ is flat (Prop.~3.3(b)). By Thm.~3.12, $\dim\mathcal O_{Z_x,z}=\dim C-\dim A$. The fibre is $Z_x=Z\times_X\operatorname{Spec}k(x)=Y\times_k\operatorname{Spec}k=Y$ (as $k(x)=k$), and $z$ corresponds to $y$, so $\mathcal O_{Z_x,z}=B$. Hence $\dim C=d+e$.

\emph{(2) $\fm_C$ is generated by $d+e$ elements.} Take affine open neighbourhoods $\operatorname{Spec}A'\ni x$, $\operatorname{Spec}B'\ni y$; then $\operatorname{Spec}(A'\otimes_kB')\ni z$, and $z$ corresponds to the kernel of the surjection $A'\otimes B'\to A'/\fm_x\otimes B'/\fm_y=k\otimes_kk=k$. This kernel is generated by $\fm_x\otimes1$ and $1\otimes\fm_y$. Therefore $\fm_C=\fm_xC+\fm_yC$, where $\fm_xC$ (resp.\ $\fm_yC$) is the ideal generated by the image of $\fm_A$ (resp.\ $\fm_B$) under $A\to C$ (resp.\ $B\to C$). Since $\fm_A=(f_1,\dots,f_d)$ and $\fm_B=(g_1,\dots,g_e)$, $\fm_C=(f_1\otimes1,\dots,f_d\otimes1,1\otimes g_1,\dots,1\otimes g_e)$. By Nakayama, $\dim_k\fm_C/\fm_C^2\le d+e$.

By (1), (2) and Cor.~2.5.14(b), $d+e=\dim C\le\dim_k\fm_C/\fm_C^2\le d+e$. So $C$ is regular. Hence $Z$ is regular.

\emph{Counterexample over a non-closed field.} Let $p$ be a prime, $k=\mathbb F_p(s)$, $L=k[\tau]/(\tau^p-s)=k(s^{1/p})$ (a field, since $\tau^p-s$ is irreducible by Eisenstein). Let $X=Y=\operatorname{Spec}L$, a variety over $k$, regular because it is a point whose local ring is a field. Let $\sigma_0\in L$ be the class of $\tau$. Then
\[ L\otimes_kL=L[\tau]/(\tau^p-s)=L[\tau]/\bigl((\tau-\sigma_0)^p\bigr)=L[\eta]/(\eta^p),\qquad\eta=\tau-\sigma_0, \]
because $\tau^p-s=\tau^p-\sigma_0^p=(\tau-\sigma_0)^p$ in characteristic $p$. This is a local Artinian ring with maximal ideal $(\eta)\ne0$ (for $p\ge2$, $\eta\neq 0$), so $X\times_kY$ is a single point with local ring of dimension $0$ and $\dim\fm/\fm^2=1$. So $X\times_kY$ is not regular (it is not even reduced), although $X$ and $Y$ are regular. \hfill$\square$

\exercise{2.13}{Let $(A,\fm)$ be a Noetherian local ring with residue field $k$. We set $\operatorname{gr}_\fm(A)=\bigoplus_{n\ge0}\fm^n/\fm^{n+1}$. This is the graded ring associated to the ideal $\fm$. It naturally has the structure of a $k$-algebra. Let us suppose that $A$ is regular of dimension $d$. We are going to show that $\operatorname{gr}_\fm(A)$ is a polynomial ring in $d$ variables over $k$. To simplify, we suppose that $k$ is infinite. (a) Let $f_1,\dots,f_d$ be a coordinate system for $\fm$. Let $\varphi:k[T_1,\dots,T_d]\to\operatorname{gr}_\fm(A)$ be the $k$-algebra homomorphism defined by $\varphi(T_i)=$ the image of $f_i$ in $\fm/\fm^2$. Show that $\varphi$ is surjective, and that $\varphi$ is an isomorphism if $d=0$. (b) Let $\tilde\fm$ be the maximal ideal of $\operatorname{gr}_\fm(A)$ generated by the $\varphi(T_i)$. Show that if $\dim\operatorname{gr}_\fm(A)_{\tilde\fm}=0$, then $\dim A=0$. (c) Let us suppose $d\ge1$. Show that $\sum_i\varphi(T_i)k$ is not contained in the union of the minimal prime ideals of $\operatorname{gr}_\fm(A)$ (here we use the $k$ infinite hypothesis). (d) Let $f\in\fm$. Show that we have a canonical surjective homomorphism $\operatorname{gr}_\fm(A)\to\operatorname{gr}_{\fm/fA}(A/fA)$. Show by induction on $d$ that $\dim\operatorname{gr}_\fm(A)=d$ and that $\varphi$ is an isomorphism.}
\pf
Write $G=\operatorname{gr}_\fm(A)=\bigoplus_{n\ge0}G_n$, $G_n=\fm^n/\fm^{n+1}$, so $G_0=k$, $G_1=\fm/\fm^2$, and $\bar f_i=\varphi(T_i)\in G_1$ the class of $f_i$.

\emph{(a)} Since $\fm=(f_1,\dots,f_d)$, $\fm^n$ is generated by the monomials of degree $n$ in the $f_i$; so $G_n=\fm^n/\fm^{n+1}$ is spanned over $k=A/\fm$ by the classes of these monomials, which are the monomials of degree $n$ in $\bar f_1,\dots,\bar f_d$. Thus $G_n\subseteq\operatorname{Im}\varphi$ for all $n$, and $\varphi$ is surjective. In particular $G$ is a Noetherian ring and $G_+:=\bigoplus_{n\ge1}G_n=(\bar f_1,\dots,\bar f_d)$. If $d=0$, then $\fm=0$ (Def.~2.7), $A$ is a field, $G=G_0=A=k$, and $\varphi:k\to k$ (the polynomial ring in $0$ variables) is the identity, an isomorphism.

\emph{(b)} The ideal $\tilde\fm=G_+$ is the kernel of $G\to G_0=k$, hence a maximal ideal. Suppose $\dim G_{\tilde\fm}=0$. Then the maximal ideal of the Noetherian local ring $G_{\tilde\fm}$ is nilpotent: there is $N$ with $\bar f_i^N=0$ in $G_{\tilde\fm}$ for all $i$, i.e.\ there are $s_i\in G\setminus\tilde\fm$ with $s_i\bar f_i^N=0$ in $G$. Put $s=\prod s_i\notin\tilde\fm$ (as $\tilde\fm$ is prime). Write $s=c+h$ with $c\in G_0=k$, $h\in G_+$; since $s\notin G_+$, $c\ne0$. Then $0=s\bar f_i^N$ has degree-$N$ component $c\bar f_i^N$, so $\bar f_i^N=0$ in $G$. Hence $G_+=(\bar f_1,\dots,\bar f_d)$ is nilpotent: $G_+^{m}=0$ for $m=d(N-1)+1$, so $G_n=0$ for $n\ge m$. This means $\fm^n=\fm^{n+1}=\fm\cdot\fm^n$ for $n\ge m$, so $\fm^n=0$ by Nakayama. Then every prime ideal of $A$ contains $\fm$, so $\operatorname{Spec}A$ is a point and $\dim A=0$.

\emph{(c)} Let $d\ge1$. $G$ is Noetherian, so it has finitely many minimal primes $P_1,\dots,P_s$. The $k$-vector space $G_1=\sum_ik\bar f_i$ is nonzero (its dimension is $d$). Suppose $G_1\subseteq P_1\cup\dots\cup P_s$. Then $G_1=\bigcup_j(G_1\cap P_j)$ is a finite union of $k$-subspaces. Since $k$ is infinite, a vector space is never a finite union of proper subspaces (take such a union $V=V_1\cup\dots\cup V_s$ with $s$ minimal, $s\geq 2$; pick $u\in V_1\setminus\bigcup_{j\ge2}V_j$ and $w\notin V_1$; the infinitely many vectors $w+\lambda u$ ($\lambda\in k$) are not in $V_1$, so two of them lie in the same $V_j$, $j\ge2$, so their difference $(\lambda-\lambda')u$ is in $V_j$, contradicting the choice of $u$). Hence $G_1\subseteq P_j$ for some $j$, so $G_+=G_1G\subseteq P_j$. As $G_+$ is a prime ideal contained in the minimal prime $P_j$, $G_+=P_j$ is a minimal prime of $G$, so $\dim G_{\tilde\fm}=0$. By (b), $\dim A=0$, contradicting $\dim A=d\ge1$.

\emph{(d)} Let $f\in\fm$, $\bar A=A/fA$ with maximal ideal $\bar\fm=\fm/fA=\fm\bar A$. Then $\bar\fm^n=(\fm^n+fA)/fA$, and the composite $\fm^n\to\bar\fm^n\to\bar\fm^n/\bar\fm^{n+1}$ kills $\fm^{n+1}$. So it induces surjections $\fm^n/\fm^{n+1}\to\bar\fm^n/\bar\fm^{n+1}$ for all $n$, compatible with multiplication (it is induced by the ring homomorphism $A\to\bar A$ mapping $\fm^n$ into $\bar\fm^n$). Their sum is a surjective homomorphism of graded $k$-algebras $\pi:G\to\bar G:=\operatorname{gr}_{\bar\fm}(\bar A)$. If $f\in\fm\setminus\fm^2$, the class of $f$ in $\fm/\fm^2=G_1$ lies in $\ker\pi$, since $f\mapsto0\in\bar A$.

\emph{Induction on $d$.} We prove: for every regular local ring $(A,\fm)$ of dimension $d$ with infinite residue field $k$, $\dim G=d$ and $\varphi$ is an isomorphism (for any coordinate system). For $d=0$ this is (a) ($G=k$ has dimension $0$). Let $d\ge1$ and assume the result in dimension $d-1$.

By (c) there is $\ell\in G_1$ lying in no minimal prime of $G$. Then $\ell\ne0$ ($0$ lies in every prime; $G\ne0$). Write $\ell=\sum c_i\bar f_i$ ($c_i\in k$) and let $f=\sum\tilde c_if_i\in\fm$ with lifts $\tilde c_i\in A$ of $c_i$; its class in $\fm/\fm^2$ is $\ell\ne0$, so $f\in\fm\setminus\fm^2$. Since $A$ is regular, hence a domain (Prop.~2.11), $f\ne0$ and Cor.~2.12 shows that $\bar A=A/fA$ is regular, of dimension $d-1$ (Cor.~2.12), with residue field $k$, still infinite. By the induction hypothesis $\dim\bar G=d-1$.

Take the surjection $\pi:G\to\bar G$ of the first part; $\pi(\ell)=0$. A chain of primes $Q_0\subsetneq\dots\subsetneq Q_{d-1}$ of $\bar G$ pulls back under $\pi$ to a chain $P_0\subsetneq\dots\subsetneq P_{d-1}$ of primes of $G$, all containing $\ker\pi\ni\ell$. Since $\ell$ lies in no minimal prime of $G$, $P_0$ is not minimal, so there is a prime $P_{-1}\subsetneq P_0$. This gives a chain of length $d$, so $\dim G\ge d$.

On the other hand $\varphi$ is surjective by (a), so $G\simeq k[T_1,\dots,T_d]/\ker\varphi$ and $\dim G\le\dim k[T_1,\dots,T_d]=d$ (Cor.~2.5.17). Hence $\dim G=d$. If $\ker\varphi$ were nonzero, then every prime of $k[T_1,\dots,T_d]$ containing $\ker\varphi$ would be nonzero, so any chain of such primes could be extended below by the prime $(0)$ of the domain $k[T_1,\dots,T_d]$, giving $\dim G\le d-1$, a contradiction. So $\ker\varphi=0$ and $\varphi$ is an isomorphism: $\operatorname{gr}_\fm(A)\simeq k[T_1,\dots,T_d]$. \hfill$\square$

\section*{4.3 Flat morphisms and smooth morphisms}

Results of the text are cited by their numbers in Liu (``Prop.~3.8'', ``Thm.~3.12'', \dots), meaning Chapter~4, Section~4.3, unless a Chapter is indicated. Honest notes. Ex.~3.1 needs $\dim\mathcal O_{S,s}=1$ (otherwise $t=0$). Ex.~3.2 is false as printed for a non-Jacobson base; it is proved under the hypothesis that closed points of $X$ map to closed points of $Y$ (automatic for varieties), and a counterexample is given. In Ex.~3.8(b) the assertion ``there is an open $V\ni y$ such that $Z\cap f^{-1}(V)\to f^{-1}(V)$ is an open immersion'' is false for non-proper $f$ of finite type; the correct statement (an open neighbourhood of the fibre $X_y$ in $X$) is proved, and the proper case is proved as stated. Ex.~3.11 is false as printed (take $X_1=Y$), and the correct statements are proved. Ex.~3.9 uses Chevalley's theorem in place of the suggested Ex.~3.2.17. Ex.~3.12 uses Ex.~3.9, and Ex.~3.13(c) uses Ex.~3.2.

\exercise{3.1}{Let $X$ be a flat scheme of finite type over a Dedekind scheme $S$. Let $s\in S$ be a closed point and $x\in X_s$. Let $t$ denote a generator of the maximal ideal of $\mathcal O_{S,s}$. We suppose that $X$ is regular at $x$. Show that $X_s$ is regular at $x$ if and only if $t\in\fm_x\setminus\fm_x^2$, where $\fm_x$ is the maximal ideal of $\mathcal O_{X,x}$.}
\pf
We assume $\dim\mathcal O_{S,s}=1$, i.e.\ $A:=\mathcal O_{S,s}$ is a discrete valuation ring with uniformizing element $t$ (if $\dim\mathcal O_{S,s}=0$ then $t=0$, the statement is empty or false: $X_s$ is regular at $x$ but $t=0\notin\fm_x\setminus\fm_x^2$). Put $B=\mathcal O_{X,x}$, $\fm=\fm_x$, $d=\dim B$, $k=B/\fm$. The point $x$ lies over $s$, so $t\in\fm$. The local ring of the fibre is
\[ \mathcal O_{X_s,x}=B/\fm_sB=B/tB. \]
As $B$ is regular, it is a domain. Since $B$ is flat over $A$ and multiplication by $t$ is injective on $A$, multiplication by $t$ is injective on $B=B\otimes_AA$. Thus $t\neq0$ in $B$, and $B$ is not a field, so $d\ge1$. We have $\dim B/tB=d-1$: the inequality $\ge d-1$ is Krull's principal ideal theorem (Thm.~2.5.15), and $\le d-1$ holds because $B$ is a domain and $t\neq0$ (a chain of primes containing $t$ can be extended by $(0)$ at the bottom). The maximal ideal of $\bar B:=B/tB$ is $\bar\fm=\fm/tB$ and
\[ \bar\fm/\bar\fm^2=\fm/(\fm^2+tB), \]
a quotient of $\fm/\fm^2$ by the line spanned by the class of $t$. Since $B$ is regular, $\dim_k\fm/\fm^2=d$.

If $t\notin\fm^2$, the class of $t$ in $\fm/\fm^2$ is non-zero, so $\dim_k\bar\fm/\bar\fm^2=d-1=\dim\bar B$, and $\bar B$ is regular.

If $t\in\fm^2$, then $\fm^2+tB=\fm^2$, so $\dim_k\bar\fm/\bar\fm^2=d>d-1=\dim\bar B$, and $\bar B$ is not regular.

Hence $X_s$ is regular at $x$ if and only if $t\in\fm\setminus\fm^2$. \hfill$\square$

\exercise{3.2}{Let $f:X\to Y$ be a morphism of finite type. Show that $f$ is flat if and only if for every closed point $y\in Y$ and for every closed point $x$ of $X_y$, $f$ is flat at $x$.}
\pf
\emph{Correction of the statement.} As printed the statement needs a hypothesis. We assume that $Y$ is quasi-compact and that $f$ maps closed points of $X$ to closed points of $Y$; the latter is automatic if $Y$ is a scheme of finite type over a field or over $\mathbb Z$ (Nullstellensatz), in particular for morphisms of algebraic varieties, which is the situation where the statement is used. Without it the statement is false (counterexample at the end). Note that if $y$ is closed in $Y$, then $X_y$ is closed in $X$, so the closed points of $X_y$ are exactly the closed points $x$ of $X$ with $f(x)=y$ closed. Hence, under our hypothesis, the closed points $x$ occurring in the statement are exactly the closed points of $X$. The implication ``$\Rightarrow$'' is trivial.

\emph{Step 1: flatness at $x$ implies flatness at every generization $x'$ of $x$.} Let $y=f(x)$, $y'=f(x')$, $A=\mathcal O_{Y,y}$, $B=\mathcal O_{X,x}$. Then $\mathcal O_{X,x'}=B_{\fq}$ for a prime $\fq$ of $B$, and $\mathcal O_{Y,y'}=A_{\fp}$ where $\fp$ is the inverse image of $\fq$ in $A$. The module $B_{\fq}$ is flat over $B$, hence over $A$. Every element of $A\setminus\fp$ is mapped outside $\fq$, so it is a unit in $B_{\fq}$, and for an $A_{\fp}$-module $N$ we have $B_{\fq}\otimes_{A_{\fp}}N=B_{\fq}\otimes_AN$. The right-hand side is exact in $N$, so $B_{\fq}$ is flat over $A_{\fp}$.

\emph{Step 2.} Suppose $f$ is flat at all closed points of $X$. Since $f$ is of finite type and $Y$ is quasi-compact, $X$ is quasi-compact. Let $x\in X$ be arbitrary. The closure of $\{x\}$ is a non-empty closed subset of a quasi-compact scheme, hence contains a closed point $x_0$ of $X$. Then $f$ is flat at $x_0$ by hypothesis, and $x$ is a generization of $x_0$, so $f$ is flat at $x$ by Step 1. Thus $f$ is flat. \hfill$\square$

\emph{Counterexample without the hypothesis on closed points.} Let $k$ be a field, $A=k[s,t]_{(s,t)}$, $Y=\operatorname{Spec}A$ and $X=\operatorname{Spec}B$ with $B=A[T]/(s,\,tT-1)$. Then $B\simeq(A/sA)[1/t]=k(t)$, so $X\to Y$ is of finite type. The only closed point of $Y$ is $\fm=(s,t)$ and $\fm B=B$ because $t$ is invertible in $B$; so $X_{\fm}=\emptyset$ and the condition of the exercise holds vacuously. But $B\neq0$ is killed by $s\neq0$, so $B$ is not torsion free over the domain $A$, hence not flat.

\exercise{3.3}{Let $f:X\to Y$ be a morphism of locally Noetherian schemes. Show that $\dim X\le\dim Y$ if $f$ is injective, and $\dim X\ge\dim Y$ if $f$ is flat and surjective.}
\pf
\emph{Injective case.} Let $Z_0\subsetneq Z_1\subsetneq\dots\subsetneq Z_n$ be irreducible closed subsets of $X$, with generic points $z_0,\dots,z_n$. Then $z_i\in\overline{\{z_{i+1}\}}$ and $z_i\neq z_{i+1}$. By continuity, $f(z_i)\in\overline{\{f(z_{i+1})\}}$, so $\overline{\{f(z_i)\}}\subseteq\overline{\{f(z_{i+1})\}}$. Since $f$ is injective, $f(z_i)\ne f(z_{i+1})$. A scheme is a $T_0$ space (an irreducible closed subset has a unique generic point), so the two closures are distinct. Thus $\overline{\{f(z_0)\}}\subsetneq\dots\subsetneq\overline{\{f(z_n)\}}$ is a chain of length $n$ of irreducible closed subsets of $Y$. Hence $\dim X\le\dim Y$ (trivially true if $\dim Y=\infty$).

\emph{Flat surjective case.} Recall that $\dim Y=\sup_{y\in Y}\dim\mathcal O_{Y,y}$ and $\dim\mathcal O_{X,x}\le\dim X$ for all $x$. Let $y\in Y$. Since $f$ is surjective, there is $x\in X$ with $f(x)=y$. By Thm.~3.12 (equality for flat $f$),
\[ \dim\mathcal O_{X_y,x}=\dim\mathcal O_{X,x}-\dim\mathcal O_{Y,y}. \]
The left-hand side is $\ge0$, so $\dim X\ge\dim\mathcal O_{X,x}\ge\dim\mathcal O_{Y,y}$. Taking the supremum over $y$ gives $\dim X\ge\dim Y$. \hfill$\square$

\exercise{3.4}{Let $k$ be a field. Let $A$ be the polynomial ring $k[x,y]$. Show, by computing the dimension of the fibers, that the canonical morphism $\operatorname{Proj}A[T_0,T_1]/(yT_0-xT_1)\to\operatorname{Spec}A$ is not flat.}
\pf
Let $X=\operatorname{Proj}A[T_0,T_1]/(yT_0-xT_1)$ and $f:X\to\operatorname{Spec}A=\mathbb A^2_k$. The scheme $X$ is covered by the two affine open sets
\[ D_+(T_0)=\operatorname{Spec}k[x,y,u]/(y-xu)\simeq\operatorname{Spec}k[x,u],\qquad D_+(T_1)=\operatorname{Spec}k[x,y,v]/(yv-x)\simeq\operatorname{Spec}k[y,v] \]
(here $u=T_1/T_0$, $v=T_0/T_1$; the relation allows us to eliminate $y$, respectively $x$). So $X$ is covered by two copies of the affine plane, and for every $x'\in X$ the local ring $\mathcal O_{X,x'}$ is a localization of a polynomial ring in two variables over $k$; hence $\dim\mathcal O_{X,x'}\le2$.

Now consider the origin $o=(x,y)\in\mathbb A^2_k$, with $k(o)=k$. The fibre is $X_o=\operatorname{Proj}k[T_0,T_1]=\mathbb P^1_k$, since the relation $yT_0-xT_1$ vanishes in $k=A/(x,y)$. Take a closed point $x'\in X_o$ (for instance $[1:0]$); then $\dim\mathcal O_{X_o,x'}=1$.

Suppose $f$ were flat. Then Thm.~3.12 gives
\[ 1=\dim\mathcal O_{X_o,x'}=\dim\mathcal O_{X,x'}-\dim\mathcal O_{\mathbb A^2,o}=\dim\mathcal O_{X,x'}-2, \]
so $\dim\mathcal O_{X,x'}=3$, contradicting $\dim\mathcal O_{X,x'}\le2$. Hence $f$ is not flat. (Alternatively, via Cor.~3.14: $f$ is surjective and $X$ is irreducible of dimension $2$, so if $f$ were flat all fibres would have dimension $\dim X-\dim Y=0$, but $\dim X_o=1$.) \hfill$\square$

\exercise{3.5}{Let $A$ be a reduced ring, and let $a\in A$ be a non-zero divisor. Let us consider a ring $B$ containing $A[T]$ and an element $s\in B$ such that $s^n=aT$ for some $n\ge1$. Show that the subalgebra $A[T][s]$ of $B$ is free (and hence flat) over $A[T]$ (consider a linear relation between $1,\dots,s^{n-1}$ and show that the coefficients are zero in $A_{\fp}[T]$ for every minimal prime ideal $\fp$ of $A$).}
\pf
We show that $1,s,\dots,s^{n-1}$ is a basis of $A[T][s]$ over $A[T]$.

\emph{Generation.} For $j\ge0$ write $j=qn+r$ with $0\le r<n$. Then $s^j=(s^n)^qs^r=(aT)^qs^r\in A[T]s^r$. So $A[T][s]=\sum_{r<n}A[T]s^r$.

\emph{Independence.} Let $c_0,\dots,c_{n-1}\in A[T]$ with $\sum_{i<n}c_is^i=0$ in $B$. Let $\fp$ be a minimal prime of $A$, and $S=A\setminus\fp$. Localizing the inclusion $A[T]\subseteq B$ at $S$ (localization is exact) gives an inclusion $K[T]\subseteq S^{-1}B$, where $K=A_{\fp}$. Since $A$ is reduced, $\fp A_{\fp}$ is the unique prime and equals the nilradical, which is $0$; so $K$ is a field. Since $A$ is reduced, every element $c\in\fp$ is a zero divisor (there is $u\notin\fp$ with $uc$ nilpotent, so $uc=0$ and $u\neq0$); hence the non-zero divisor $a$ is not in $\fp$ and is a unit in $K$. In $S^{-1}B$ we have $s^n=aT$ with $a\in K^*$.

Let $K[S]$ be the polynomial ring in one variable and define the $K$-algebra homomorphism $\psi:K[S]\to S^{-1}B$, $S\mapsto s$. Its image is $K[s]\ni aT$, hence contains $T$. Let $I=\ker\psi$. We claim that $I=0$. The subring $K[S^n]\subseteq K[S]$ maps onto $K[s^n]=K[aT]=K[T]$, and $\psi$ restricted to $K[S^n]$ is injective because $K[T]\to S^{-1}B$ is injective and $S^n\mapsto aT$ (the elements $aT$ and $T$ generate the same polynomial ring). So $I\cap K[S^n]=0$. If $I\ne0$, choose $0\ne g\in I$. As $K[S]$ is integral over $K[S^n]$, there is an equation $g^m+e_{m-1}g^{m-1}+\dots+e_0=0$ with $e_i\in K[S^n]$ and $m$ minimal. Then $e_0\neq0$ (otherwise, $K[S]$ being a domain, we could cancel a factor $g$ and lower $m$), and $e_0=-g(g^{m-1}+\dots+e_1)\in I\cap K[S^n]$, a contradiction. So $\psi$ is injective, and $1,s,\dots,s^{n-1}$ are linearly independent over $K[T]$: indeed $K[S]$ is free over $K[S^n]=K[aT]=K[T]$ with basis $1,S,\dots,S^{n-1}$.

Consequently the image of each $c_i$ in $K[T]=A_{\fp}[T]$ is zero, for every minimal prime $\fp$. Since $A$ is reduced, the intersection of its minimal primes is $0$, so $A\to\prod_{\fp}A/\fp$ is injective. The field $A_{\fp}$ is $\operatorname{Frac}(A/\fp)$ and $A\to A_\fp$ factors through $A/\fp\subseteq\operatorname{Frac}(A/\fp)$. Hence $A\to\prod_{\fp}A_{\fp}$ is injective, and so is $A[T]\to\prod_{\fp}A_{\fp}[T]$ (coefficientwise). Therefore all $c_i=0$.

So $A[T][s]$ is a free $A[T]$-module with basis $1,\dots,s^{n-1}$, and in particular flat. \hfill$\square$

\exercise{3.6}{Let $f:X\to Y$ be a morphism of integral schemes. We suppose that the generic fiber of $f$ is geometrically integral (the definition is the same as for algebraic varieties 3.2.8). Show that for any flat morphism $Z\to Y$ with $Z$ integral, we have $X\times_YZ$ integral. This partially generalizes Exercise 3.2.14.}
\pf
Let $\eta$ be the generic point of $Y$, $K=K(Y)=k(\eta)$, $\xi$ the generic point of $X$, $\zeta$ the generic point of $Z$, $g:Z\to Y$ the given flat morphism, and $W=X\times_YZ$.

\emph{Preliminaries.}

(i) The generic fibre $X_\eta$ is geometrically integral, hence non-empty. Since $\xi$ is a generization of every point of $X$, $f(\xi)$ is a generization of $\eta$, so $f(\xi)=\eta$; thus $\xi\in X_\eta$. Similarly $g(\zeta)=\eta$: by Lemma~3.7, $g(Z)$ is dense in $Y$, and $g(Z)\subseteq\overline{\{g(\zeta)\}}$.

(ii) \emph{If $X_\eta$ is geometrically integral over $K$, then $X_\eta\otimes_KL$ is integral for every field extension $L/K$.} First a lemma: \emph{if $\Omega$ is an algebraically closed field, $T$ an integral $\Omega$-scheme (not necessarily of finite type) and $L/\Omega$ a field extension, then $T_L$ is integral.} Let $\operatorname{Spec}B\subseteq T$ be a non-empty affine open, so $B$ is a domain. Write $B$ and $L$ as directed unions of finitely generated $\Omega$-subalgebras $B_0$, $L_0$ (domains). By Ex.~3.2.15 (the case where $\operatorname{Spec}B_0$ is integral over the algebraically closed field $\Omega$, hence geometrically integral, and $\operatorname{Spec}L_0$ is integral) each $B_0\otimes_\Omega L_0$ is a domain. These rings are subrings of $B\otimes_\Omega L$ (tensor products over a field preserve injections), and $B\otimes_\Omega L$ is their directed union; so it is a domain. Thus $T_L$ is covered by the non-empty integral affine opens $U_L$, $U\subseteq T$ non-empty affine open. Two of them meet: $U\cap U'$ is a non-empty open subset of $T$, it contains a non-empty affine open $W$, and $W_L=\operatorname{Spec}(\mathcal O(W)\otimes_\Omega L)$ is non-empty since $\mathcal O(W)\otimes_\Omega L\neq0$. A space covered by irreducible opens which pairwise meet is irreducible, and $T_L$ is reduced, so $T_L$ is integral.

Now choose an algebraic closure $\bar L$ of $L$ and let $\bar K$ be the algebraic closure of $K$ inside $\bar L$. The scheme $(X_\eta)_{\bar K}$ is integral by hypothesis, so by the lemma $(X_\eta)_{\bar L}=(X_\eta)_{\bar K}\times_{\bar K}\bar L$ is integral. Finally $(X_\eta)_{\bar L}\to(X_\eta)_L$ is faithfully flat (base change of $L\to\bar L$), hence surjective, so $(X_\eta)_L$ is irreducible, and it is reduced because its local rings embed in those of $(X_\eta)_{\bar L}$ (flat local homomorphisms are faithfully flat, hence injective).

(iii) The same holds for every non-empty open subscheme $U\subseteq X_\eta$, since $U_L$ is a non-empty open subscheme of the integral scheme $(X_\eta)_L$.

\emph{The affine case.} Let $V=\operatorname{Spec}A\subseteq Y$, $\operatorname{Spec}B\subseteq f^{-1}(V)$, $\operatorname{Spec}C\subseteq g^{-1}(V)$ be non-empty affine open subsets, where $A,B,C$ are domains. The generic points $\xi,\zeta$ lie in them, so $A\to B$ and $A\to C$ are injective (the primes $\ker(A\to B)$ and $\ker(A\to C)$ are the images of the generic points, i.e.\ $0$). The $A$-module $C$ is flat. Let $K=\operatorname{Frac}(A)$, $B_K=B\otimes_AK$, $C_K=C\otimes_AK$ and $L=\operatorname{Frac}(C)\supseteq K$. We have
\[ B\otimes_AC\ \hookrightarrow\ B_K\otimes_AC=B_K\otimes_K C_K\ \hookrightarrow\ B_K\otimes_KL . \]
The first map is injective because $B\to B_K$ is injective ($B$ is a domain containing $A$) and $C$ is $A$-flat. The equality holds because $B_K\otimes_AC=B_K\otimes_K(K\otimes_AC)$. The last map is injective because $C_K\subseteq L$ and $B_K$ is free over the field $K$. Now $\operatorname{Spec}B_K$ is a non-empty open subscheme of $X_\eta$ (it is the part of $X_\eta$ inside $\operatorname{Spec}B$), so $B_K\otimes_KL$ is a domain by (ii) and (iii). Hence $B\otimes_AC$ is a domain: every such affine open piece $U\times_VZ'$ of $W$ is integral.

\emph{Globalization.} The schemes $U\times_VZ'$, where $V\subseteq Y$, $U\subseteq f^{-1}(V)$, $Z'\subseteq g^{-1}(V)$ are non-empty affine opens, form an open cover of $W$ by integral schemes; in particular $W$ is reduced. They are non-empty and pairwise intersect: given two of them, $W_1=U_1\times_{V_1}Z_1$ and $W_2=U_2\times_{V_2}Z_2$, the open subscheme $(U_1\cap U_2)\times_Y(Z_1\cap Z_2)$ of $W$ is contained in both, and it is non-empty: $\xi\in U_1\cap U_2$ and $\zeta\in Z_1\cap Z_2$ both lie over $\eta$, and the ring $k(\xi)\otimes_Kk(\zeta)$ of two field extensions of $K$ is non-zero, so there is a point of $W$ over $(\xi,\zeta)$. A topological space covered by non-empty irreducible open subsets that pairwise intersect is irreducible. So $W$ is irreducible and reduced, i.e.\ integral. \hfill$\square$

\exercise{3.7}{Let $f:X\to Y$ be a morphism of schemes. Let $X_1,X_2$ be two closed subschemes of $X$, flat over $Y$. We suppose that there exists a dense open subscheme $V$ of $Y$ such that $X_1\cap f^{-1}(V)$ and $X_2\cap f^{-1}(V)$ are equal as closed subschemes of $f^{-1}(V)$. Let us suppose that $Y$ is integral. We want to show that $X_1=X_2$ as schemes.

(a) Show that one can suppose that $X=\operatorname{Spec}B$, $Y=\operatorname{Spec}A$, and $V$ are affine.

(b) Let $I,J$ be the ideals of $B$ defining $X_1$ and $X_2$. Tensoring the exact sequence $0\to A\to\mathcal O_Y(V)$ by $B/I$, show that $I=\operatorname{Ker}(B\to\mathcal O_X(f^{-1}(V))/(I))$. Deduce from this that $I=J$.}
\pf
(a) Equality of two closed subschemes of $X$ can be checked on an open cover of $X$. Cover $Y$ by non-empty affine open sets $Y'$ and each $f^{-1}(Y')$ by affine open sets $X'$. Then $V\cap Y'$ is a non-empty open subset of $Y'$ (as $V$ is dense and $Y'$ is open and non-empty), so it contains a non-empty principal open subset $D(a)=\operatorname{Spec}A_a$ of $Y'=\operatorname{Spec}A$; this $D(a)$ is dense in the integral scheme $Y'$, and $A$ is a domain with $a\ne0$ (as $D(a)\ne\emptyset$). The hypotheses are inherited: $X_i\cap X'$ are closed subschemes of $X'$, flat over $Y'$ (flatness is local), and they coincide over $D(a)$. So we may assume that $Y=\operatorname{Spec}A$ with $A$ a domain, $X=\operatorname{Spec}B$, and $V=D(a)=\operatorname{Spec}A_a$ with $a\ne0$.

(b) The closed subschemes $X_1,X_2$ are $\operatorname{Spec}B/I$ and $\operatorname{Spec}B/J$, and $B/I$, $B/J$ are flat over $A$. Since $A$ is a domain and $a\ne0$, the map $A\to A_a$ is injective. Tensoring with the flat $A$-module $B/I$ gives an injection
\[ B/I\ \hookrightarrow\ (B/I)\otimes_AA_a=(B/I)_a=B_a/IB_a . \]
Note $B_a=\mathcal O_X(f^{-1}(V))$, so the right-hand side is $\mathcal O_X(f^{-1}(V))/(I)$. The injectivity of the displayed map says: the kernel of the composition $B\to B/I\to B_a/IB_a$ is exactly $I$, that is,
\[ I=\operatorname{Ker}\bigl(B\to B_a/IB_a\bigr). \]
(Concretely: $b\mapsto0$ iff $a^nb\in I$ for some $n$, and flatness gives $b\in I$.) The same holds for $J$. By hypothesis the closed subschemes of $\operatorname{Spec}B_a$ defined by $IB_a$ and $JB_a$ coincide, i.e.\ $IB_a=JB_a$. Hence
\[ I=\operatorname{Ker}(B\to B_a/IB_a)=\operatorname{Ker}(B\to B_a/JB_a)=J. \]
So $X_1=X_2$ on each piece, hence globally. \hfill$\square$

\exercise{3.8}{Let $f:X\to Y$ be a morphism of locally Noetherian schemes. Let $Z$ be a closed subscheme of $X$. We suppose that there exists a point $y\in Y$ such that $Z_y=X_y$ as schemes.

(a) Show that if $Z$ is flat over $Y$ at $z\in X_y$, then $Z$ equals $X$ in a neighborhood of $z$ (use Exercise 1.2.14 and Nakayama's lemma).

(b) Show that if $Z$ is flat over $Y$ and if $f$ is moreover of finite type (resp. proper) over $Y$, then there exists an open neighborhood $V\ni y$ such that $Z\cap f^{-1}(V)\to f^{-1}(V)$ is an open immersion (resp. an isomorphism) (use Exercise 3.2.5).}
\pf
(a) Let $A=\mathcal O_{Y,y}$ with maximal ideal $\fm_y$ and residue field $k(y)$, $B=\mathcal O_{X,z}$, and let $I\subseteq B$ be the ideal of $Z$ at $z$ (so $\mathcal O_{Z,z}=B/I$). The ring $B$ is Noetherian, so $I$ is a finitely generated $B$-module. The local ring at $z$ of the fibre $Z_y\subseteq X_y$ is $B/(I+\fm_yB)$, while that of $X_y$ is $B/\fm_yB$. The hypothesis $Z_y=X_y$ means $I\subseteq\fm_yB$.

Since $B/I$ is flat over $A$, $\operatorname{Tor}_1^A(B/I,k(y))=0$, so tensoring $0\to I\to B\to B/I\to0$ with $k(y)=A/\fm_y$ gives an exact sequence
\[ 0\to I/\fm_yI\to B/\fm_yB\to (B/I)/\fm_y(B/I)\to0 . \]
The last map is the quotient $B/\fm_yB\to B/(I+\fm_yB)$, an isomorphism because $I\subseteq\fm_yB$. So the map $I/\fm_yI\to B/\fm_yB$ is both zero (its image is the kernel of the iso, which is $0$) and injective; hence $I=\fm_yI$. As $\fm_yB\subseteq\fm_z$, we get $I=\fm_zI$, and Nakayama's lemma gives $I=0$. Thus the coherent ideal sheaf $\mathcal I_Z$ of $Z$ has zero stalk at $z$; being of finite type, it vanishes on an open neighbourhood $U_z$ of $z$ in $X$. So $Z\cap U_z=U_z$ as closed subschemes.

(b) \emph{What is true.} Since $Z_y=X_y$, every point of $X_y$ lies in $Z$, and $Z$ is flat over $Y$ at it. By (a) there is for each $z\in X_y$ an open set $U_z\ni z$ of $X$ with $Z\cap U_z=U_z$. Let $U=\bigcup_{z\in X_y}U_z$. It is an open subset of $X$ containing $X_y$ and $Z\cap U=U$; in other words $Z\cap U\to U$ is an isomorphism (in particular an open immersion). This holds for any locally Noetherian $f$.

\emph{The proper case.} Suppose $f$ is proper. The set $C=X\setminus U$ is closed, so $f(C)$ is closed in $Y$, and $y\notin f(C)$ because $X_y\subseteq U$. Put $V=Y\setminus f(C)$, an open neighbourhood of $y$. Then $f^{-1}(V)\subseteq U$, so $Z\cap f^{-1}(V)=f^{-1}(V)$, i.e.\ $Z\cap f^{-1}(V)\to f^{-1}(V)$ is an isomorphism.

\emph{The finite type case: correction.} For $f$ merely of finite type, the neighbourhood $V$ of $y$ in $Y$ cannot be found in general, so the statement must be read as: there is an open neighbourhood $U$ of $X_y$ in $X$ such that $Z\cap U\to U$ is an open immersion (even an isomorphism); this was proved above. Here is a counterexample to the version with $V$. Let $\mathcal V$ be a discrete valuation ring with fraction field $K$, closed point $y$ and residue field $k$, and $Y=\operatorname{Spec}\mathcal V$. Let $X=\operatorname{Spec}\mathcal V[x]\sqcup\operatorname{Spec}K[x]$ (finite type over $Y$), and let $Z=\operatorname{Spec}\mathcal V[x]\sqcup V(x)$, where $V(x)=\operatorname{Spec}K\subset\operatorname{Spec}K[x]$. Then $Z$ is a closed subscheme of $X$, flat over $Y$ (a disjoint union of the flat schemes $\operatorname{Spec}\mathcal V[x]$ and $\operatorname{Spec}K$), and $Z_y=\operatorname{Spec}k[x]=X_y$. The only open neighbourhood of $y$ in $Y$ is $Y$ itself, and $Z\to X$ is not an open immersion, since $\{x=0\}$ is a closed non-open point of the irreducible curve $\operatorname{Spec}K[x]$. \hfill$\square$

\exercise{3.9}{Let $f:X\to Y$ be a flat morphism of finite type with $X,Y$ Noetherian. We will show that $f$ is open.

(a) Use Exercise 3.2.17 to show that $f(X)$ contains a non-empty open subset $V$ of $Y$.

(b) By considering the morphism $X\times_Y(Y\setminus V)\to Y\setminus V$, show that $f(X)\setminus V$ contains a non-empty open subset of $Y\setminus V$.

(c) Show that $Y\setminus f(X)$ is closed in $Y$ by using the fact that the topological space $Y$ is Noetherian; that is, every descending chain of closed subsets in $Y$ is stationary.}
\pf
We prove the statement by a direct argument which contains the content of the steps (a)--(c) (the finiteness of the situation being used through Chevalley's theorem, of which Ex.~3.2.17 is the special case needed for (a), and through the Noetherian property of $Y$, as suggested in (c)). It suffices to prove that $f(X)$ is open for every flat morphism $f:X\to Y$ of finite type of Noetherian schemes: if $W\subseteq X$ is open, then $W$ is Noetherian and $W\to Y$ is flat of finite type, and then $f(W)$ is open.

\emph{Step 1: $f(X)$ is stable under generization.} Let $x\in X$, $y=f(x)$, and let $y'$ be a generization of $y$. The local homomorphism $\mathcal O_{Y,y}\to\mathcal O_{X,x}$ is flat, hence faithfully flat, hence $\operatorname{Spec}\mathcal O_{X,x}\to\operatorname{Spec}\mathcal O_{Y,y}$ is surjective. The points of $\operatorname{Spec}\mathcal O_{Y,y}$ (resp.\ $\operatorname{Spec}\mathcal O_{X,x}$) are the generizations of $y$ in $Y$ (resp.\ of $x$ in $X$), so $y'=f(x')$ for some generization $x'$ of $x$. In particular $y'\in f(X)$.

\emph{Step 2: $f(X)$ is constructible} by Chevalley's theorem ([EGA IV, 1.8.4]; for $f$ of finite type between Noetherian schemes it is the content of Ex.~3.2.17 together with a Noetherian induction as in (b), (c)).

\emph{Step 3: a constructible subset $E$ of a Noetherian scheme $Y$ stable under generization is open.} Let $C=Y\setminus E$; it is constructible and stable under specialization (if $c\in C$ and $c'\in\overline{\{c\}}$, then $c'\in C$, else $c'\in E$ would give $c\in E$). Write $C=\bigcup_{i=1}^rU_i\cap F_i$ with $U_i$ open, $F_i$ closed irreducible and $U_i\cap F_i\neq\emptyset$ (decompose each closed part into its finitely many irreducible components and discard empty pieces). The non-empty open subset $U_i\cap F_i$ of the irreducible space $F_i$ contains the generic point $\xi_i$ of $F_i$. As $C$ is stable under specialization, $F_i=\overline{\{\xi_i\}}\subseteq C$. Therefore $C=\bigcup F_i$ is closed, i.e.\ $E$ is open.

By Steps 1--3 with $E=f(X)$, $f(X)$ is open, hence $f$ is open. (Step 3 is the Noetherian induction of (c), and the existence of a non-empty open subset in a dense constructible set is (a), (b).) \hfill$\square$

\exercise{3.10}{Let $k$ be a field of characteristic $\operatorname{char}(k)\ne2$. Let $Y=\operatorname{Spec}k[u,v]/(v^2-u^2(u+1))$, and let $f:X\to Y$ be the normalization morphism. Show that $f$ is unramified, surjective, but not etale.}
\pf
Let $R=k[u,v]/(v^2-u^2(u+1))$.

\emph{$R$ is a domain.} If $F=v^2-u^2(u+1)$ were reducible in $k[u][v]$ (monic in $v$), it would have a root $g\in k[u]$, so $g^2=u^2(u+1)$; then $u\mid g$ and $u+1=(g/u)^2$ would be a square in $k[u]$, impossible by degree. So $F$ is irreducible, and $R$ is a domain of dimension $1$.

\emph{The normalization.} Put $t=v/u\in\operatorname{Frac}(R)$. Then $t^2=v^2/u^2=u+1$. Define $\varphi:R\to k[t]$ by $u\mapsto t^2-1$, $v\mapsto t(t^2-1)$; this is well defined since $v^2=t^2(t^2-1)^2=u^2(u+1)$. It is injective: its kernel is a prime ideal of the one-dimensional domain $R$; if non-zero it would be maximal, and the image would be a field, but the image contains $k[t^2-1]$, a polynomial ring, which is not a field. Moreover $\varphi(v)/\varphi(u)=t$, so $\operatorname{Frac}(R)=k(t)$ and $k[t]$ is integral over $R$ ($t^2=u+1$). Since $k[t]$ is integrally closed, it is the integral closure of $R$. Therefore $X=\operatorname{Spec}k[t]=\mathbb A^1_k$ and $f$ is given by $\varphi$. It is a finite morphism.

\emph{Structure of $f$.} Localizing at $u$: the map $R_u\to k[t]_{t^2-1}$ is injective, and surjective since $t=\varphi(v)/\varphi(u)$ and $1/(t^2-1)=1/\varphi(u)$ lie in the image. So $f$ restricts to an isomorphism $f^{-1}(D(u))=D(t^2-1)\to D(u)$. The complement of $D(u)$ in $Y$ is $V(u)=V(u,v^2)=\{o\}$, the point $o=(u,v)$ with $k(o)=k$ (note $R/uR=k[v]/(v^2)$). Its fibre is
\[ X_o=\operatorname{Spec}k[t]/(t^2-1,\,t^3-t)=\operatorname{Spec}k[t]/(t^2-1). \]
As $\operatorname{char}k\ne2$, $t^2-1=(t-1)(t+1)$ with distinct roots in $k$, so $X_o=\operatorname{Spec}k\sqcup\operatorname{Spec}k$.

\emph{Surjective.} $f$ maps $f^{-1}(D(u))$ isomorphically onto $D(u)$, and the fibre over $o$ is non-empty. So $f$ is surjective.

\emph{Unramified.} By Lemma~3.20 it suffices to check that every fibre is finite, reduced, with separable residue fields. Over $y\in D(u)$ the fibre is $\operatorname{Spec}k(y)$ (as $f$ is an isomorphism there). Over $o$ the fibre is $\operatorname{Spec}k\sqcup\operatorname{Spec}k$, reduced with residue fields equal to $k(o)=k$. So $f$ is unramified.

\emph{Not etale.} Let $x$ be the point $t=1$ of $X$; $f(x)=o$. The ring $\mathcal O_{Y,o}=R_{(u,v)}$ has dimension $1$. Its maximal ideal $\fm_o=(u,v)$ satisfies $\fm_o/\fm_o^2=\fm/(\fm^2+(F))$ computed in $k[u,v]_{(u,v)}$, and $F=v^2-u^3-u^2\in(u,v)^2$, so $\dim_k\fm_o/\fm_o^2=2>1$. Hence $Y$ is not regular at $o$, whereas $X=\mathbb A^1_k$ is regular at $x$. If $f$ were etale at $x$, then by Cor.~3.24 $X$ regular at $x$ would be equivalent to $Y$ regular at $f(x)=o$, a contradiction. So $f$ is not etale at $x$ (equivalently, as $f$ is unramified, $f$ is not flat at $x$). \hfill$\square$

\exercise{3.11}{Let $f_1:X_1\to Y$, $f_2:X_2\to Y$ be morphisms of locally Noetherian schemes of finite type. Let us suppose that for every $y\in Y$, there exists $i=1$ or $2$ such that $f_i$ is unramified (resp.\ etale; resp.\ smooth) at all points of $f_i^{-1}(y)$. Show that $X_1\times_YX_2\to Y$ is unramified (resp.\ etale; resp.\ smooth).}
\pf
\emph{The statement is false as printed.} Take $X_1=Y$, $f_1=\operatorname{id}$ (etale, smooth and unramified at all points), and $f_2:X_2\to Y$ any morphism of finite type, e.g.\ $\mathbb A^1_Y\to Y$ with $Y$ the spectrum of a field. Then $X_1\times_YX_2=X_2$, and the claim would say that $\mathbb A^1_k\to\operatorname{Spec}k$ is unramified, which is false (the fibre is not finite) and, for ``etale'', too. The correct statements are the two following ones, which I prove: with $W=X_1\times_YX_2$, $h:W\to Y$ the structure morphism and $p_i:W\to X_i$ the projections, and where ``$P$'' stands for unramified, etale, or smooth and a morphism has $P$ over $y$ if it has $P$ at all points of the fibre over $y$:

(A) if for a given $y\in Y$ \emph{both} $f_1$ and $f_2$ have $P$ over $y$, then $h$ has $P$ over $y$ (so if the hypothesis holds with ``both'', $h$ has $P$; this is the pointwise version of Prop.~3.22(c) and Prop.~3.38);

(B) if for a given $y\in Y$, $f_i$ has $P$ over $y$ for one index $i$, then the projection $p_j:W\to X_j$ ($j\ne i$) has $P$ at all points of $h^{-1}(y)$.

For all of them, $W$ is locally Noetherian and $h,p_j$ are of finite type.

\emph{Fibres.} For $y\in Y$ the fibre of $h$ is $W_y=X_{1,y}\times_{k(y)}X_{2,y}$. Moreover, if $x_2\in X_2$ lies over $y$, the fibre of $p_2$ over $x_2$ is $X_{1,y}\times_{k(y)}\operatorname{Spec}k(x_2)=(X_{1,y})_{k(x_2)}$, and similarly for $p_1$.

\emph{Unramified.} By Lemma~3.20 and its proof, $f_i$ is unramified over $y$ if and only if $X_{i,y}$ is finite, reduced, with separable residue field extensions $k(x)/k(y)$; this means $X_{i,y}=\bigsqcup_m\operatorname{Spec}L_m$ with $L_m/k(y)$ finite separable. We use the following fact. If $L/\kappa$ is finite separable and $M$ is a field extension of $\kappa$, then $L\otimes_\kappa M$ is a finite product of finite separable field extensions of $M$: write $L=\kappa[T]/(P)$ with $P$ separable; over $M$ the polynomial $P$ is a product of distinct irreducible factors with distinct roots, so $L\otimes_\kappa M=M[T]/(P)$ is the product of the corresponding finite separable extensions of $M$.

(A), unramified: $X_{1,y}=\bigsqcup\operatorname{Spec}L_m$ and $X_{2,y}=\bigsqcup\operatorname{Spec}L'_n$ with $L_m,L'_n$ finite separable over $k(y)$. So $W_y=\bigsqcup\operatorname{Spec}(L_m\otimes_{k(y)}L'_n)$, and each $L_m\otimes L'_n$ is a finite product of fields, finite and separable over $L'_n$ and hence over $k(y)$ (a tower of finite separable extensions is separable). Thus $W_y$ is finite, reduced and has separable residue fields, i.e.\ $h$ is unramified over $y$.

(B), unramified: the fibre of $p_2$ over $x_2$ is $(X_{1,y})_{k(x_2)}=\bigsqcup_m\operatorname{Spec}(L_m\otimes_{k(y)}k(x_2))$, finite, reduced with separable residue fields over $k(x_2)$ by the same computation. Hence $p_2$ is unramified at all points over $x_2$, by Lemma~3.20 applied to $p_2$ (and its fibre over $x_2$; the local-ring condition $\fm_{x_2}\mathcal O_{W,w}=\fm_w$ is a property of the fibre).

\emph{Flatness.} Let $w\in W$ with images $x_1,x_2,y$, and put $A=\mathcal O_{Y,y}$, $B_i=\mathcal O_{X_i,x_i}$. Then $\mathcal O_{W,w}$ is a localization of $B_1\otimes_AB_2$. If $A\to B_1$ and $A\to B_2$ are flat, then $B_1\otimes_AB_2$ is flat over $B_2$ (base change of $A\to B_1$) and $B_2$ is flat over $A$, so $B_1\otimes_AB_2$ is flat over $A$; localizing, $\mathcal O_{W,w}$ is flat over $A$. Similarly, if $A\to B_1$ is flat then $B_2\to B_1\otimes_AB_2$ is flat, so $p_2$ is flat at $w$. Hence: if $f_1,f_2$ are flat over $y$, then $h$ is flat over $y$ (needed in (A)); if $f_1$ is flat over $y$, $p_2$ is flat at all points over $y$ (needed in (B)).

\emph{Etale.} Etale means unramified and flat, so (A) and (B) for etale follow from the unramified and flat cases.

\emph{Smooth.} Smooth at $x$ means flat at $x$ and $X_y\to\operatorname{Spec}k(y)$ smooth at $x$. For (A): $h$ is flat over $y$ by the above, and $W_y=X_{1,y}\times_{k(y)}X_{2,y}$ is smooth over $k(y)$ at all its points by Prop.~3.38 (stability under fibered products), since $X_{i,y}$ is smooth over $k(y)$ (a smooth-at-$x$ statement at all points of the fibre is the same as smoothness of $X_{i,y}\to\operatorname{Spec}k(y)$, and over a field flatness is automatic). For (B): $p_2$ is flat at points over $y$ by the above, and the fibre of $p_2$ over $x_2$ is $(X_{1,y})_{k(x_2)}$, smooth over $k(x_2)$ by stability of smoothness under base change (Prop.~3.38). \hfill$\square$

\exercise{3.12}{Let $f:X\to Y$ be a flat morphism of algebraic varieties. We suppose that $X_y$ is irreducible for every closed point $y\in Y$ and that $Y$ is irreducible. Show that $X$ is irreducible.}
\pf
Algebraic varieties are of finite type over a field $k$, so $X$ and $Y$ are Noetherian and the closed points of $Y$ are dense in every non-empty open subset of $Y$ (Nullstellensatz). Irreducible includes non-empty: $Y\ne\emptyset$ has a closed point $y$, and $X_y\neq\emptyset$, so $X\ne\emptyset$.

Let $X_1,\dots,X_r$ be the irreducible components of $X$ ($X$ is Noetherian). Suppose $r\ge2$. For $i=1,2$ let
\[ U_i=X\setminus\bigcup_{j\ne i}X_j . \]
This is a non-empty open subset of $X$ (it is non-empty since no $X_i$ is contained in the union of the others), and $U_1\cap U_2=\emptyset$ because a point of $U_1$ lies in no $X_j$ with $j\ne1$, in particular not in $X_2$, whereas $U_2\subseteq X_2$. The morphisms $f|_{U_i}:U_i\to Y$ are flat and of finite type, so $f(U_i)$ is open in $Y$ by Ex.~3.9. Moreover, $f(U_i)\neq\emptyset$. Since $Y$ is irreducible, $f(U_1)\cap f(U_2)$ is a non-empty open subset of $Y$ (two non-empty opens of an irreducible space meet), and so it contains a closed point $y$ of $Y$. Then $X_y$ meets both $U_1$ and $U_2$, i.e.\ $U_1\cap X_y$ and $U_2\cap X_y$ are non-empty open subsets of $X_y$. As $X_y$ is irreducible, they meet, which contradicts $U_1\cap U_2=\emptyset$. Hence $r=1$ and $X$ is irreducible. \hfill$\square$

\exercise{3.13}{Let $X$ be a smooth morphism over a scheme $S$ of characteristic $p>0$.

(a) Show that $X^{(p)}\to S$ is smooth (use Proposition 3.38).

(b) Let us suppose that $S$ is the spectrum of a field $k$. Let $F_{X/k}:X\to X^{(p)}$ be the relative Frobenius. Show that $F_{X/k}$ is flat at the rational points of $X$ (use Proposition 2.27).

(c) Show that $F_{X/k}$ is faithfully flat, that is flat and surjective.

(d) Let us suppose $p=2$. Let $Y=\operatorname{Spec}k[x,y]/(y^2+x^3)$. Show that $F_{Y/k}$ is not flat.}
\pf
Recall the definitions: $X^{(p)}=X\times_{S,F_S}S$, where $F_S:S\to S$ is the absolute Frobenius ($S$ has characteristic $p$, and $F_S$ is the identity on the underlying space and $a\mapsto a^p$ on $\mathcal O_S$), viewed as an $S$-scheme by the second projection; and $F_{X/S}:X\to X^{(p)}$ is the morphism induced by the absolute Frobenius $F_X$ and the structure map $X\to S$. We assume $S$ locally Noetherian, so that smoothness is defined.

(a) The morphism $X^{(p)}\to S$ is the base change of the smooth morphism $X\to S$ by $F_S:S\to S$. Smooth morphisms are stable under base change (Prop.~3.38), so $X^{(p)}\to S$ is smooth.

For (b)--(d) let $S=\operatorname{Spec}k$, $\operatorname{char}k=p$, and $k'$ denote the ring $k$ viewed as a $k$-algebra via $F:\lambda\mapsto\lambda^p$. Then for an affine open $\operatorname{Spec}P$ of $X$, the corresponding open subscheme of $X^{(p)}$ is $\operatorname{Spec}(P\otimes_kk')$, a $k$-algebra via the second factor, and $F_{X/k}$ corresponds to the $k$-algebra map $\Phi:P\otimes_kk'\to P$, $a\otimes\mu\mapsto\mu a^p$.

(b) Let $x\in X(k)$ be a rational point, $A=\mathcal O_{X,x}$ with maximal ideal $\fm$, $A/\fm=k$, $d=\dim A$. Since $X$ is smooth over $k$, $X$ is regular at $x$ (Prop.~3.30). Let $f_1,\dots,f_d$ be a coordinate system of $A$. (Since $x$ is a rational point, $x$ is closed, and we may replace $X$ by an affine open neighbourhood of $x$, which is an algebraic variety, so that Prop.~3.30 and Prop.~2.27 apply.) The $k$-algebra map $k[[T_1,\dots,T_d]]\to\widehat A$, $T_i\mapsto f_i$, is an isomorphism (this is Prop.~2.27, with the given coordinate system): it exists as $\widehat A$ is complete and contains $k$, it is surjective because $\widehat{\fm}=(f_1,\dots,f_d)$ and $\widehat A/\widehat\fm=k$, and its kernel is a prime ideal of the $d$-dimensional domain $k[[T]]$ with quotient of dimension $\dim\widehat A=\dim A=d$ (Lemma~2.26), hence it is $0$.

Put $B=A\otimes_kk'$ and $\Phi:B\to A$. We have $B/(\fm\otimes1)B=k\otimes_kk'=k'$, a field, so $\fn:=(\fm\otimes1)B=(s_1,\dots,s_d)B$ with $s_i=f_i\otimes1$ is a maximal ideal. Moreover $\Phi(\fn)\subseteq\fm$, so $\Phi^{-1}(\fm)=\fn$, and $B_{\fn}=\mathcal O_{X^{(p)},x'}$ where $x'=F_{X/k}(x)$, which is a rational point of $X^{(p)}$ (its residue field is $B/\fn=k$ as a $k$-algebra via the second factor). The local homomorphism $\Phi_x:B_{\fn}\to A$ is $F^\#_{X/k,x}$.

\emph{$B_{\fn}$ is regular of dimension $d$ with coordinate system $s_1,\dots,s_d$.} Since $\fn$ is generated by $d$ elements, $d':=\dim B_{\fn}\le d$. On the other hand $A/\fn A=A/(f_1^p,\dots,f_d^p)$ is of dimension $0$, and Thm.~3.12 (inequality) for $\operatorname{Spec}A\to\operatorname{Spec}B_{\fn}$ gives $d\le d'+0$. So $d'=d$, and $B_{\fn}$ is a Noetherian local ring of dimension $d$ with maximal ideal generated by $s_1,\dots,s_d$; it is regular (this also follows from (a) and Prop.~3.30).

Applying the argument of Prop.~2.27 to $X^{(p)}$ at the rational point $x'$ (regular by (a) and Prop.~3.30) with the coordinate system $s_i$, we obtain an isomorphism $k[[S_1,\dots,S_d]]\to\widehat{B_{\fn}}$, $S_i\mapsto s_i$. The induced map of completions $\widehat\Phi:\widehat{B_{\fn}}\to\widehat A$ is continuous and $k$-linear, and $\widehat\Phi(s_i)=f_i^p$. Under the identifications with power series rings it is
\[ k[[S_1,\dots,S_d]]\to k[[T_1,\dots,T_d]],\qquad S_i\mapsto T_i^p. \]
Now $k[[T_1,\dots,T_d]]$ is a free module over $k[[T_1^p,\dots,T_d^p]]\simeq k[[S_1,\dots,S_d]]$ with basis the monomials $T^\alpha$, $0\le\alpha_i<p$ (write every exponent as $pq+r$ with $0\le r<p$). So $\widehat\Phi$ is flat.

\emph{Descent of flatness to the local rings.} $\widehat{B_{\fn}}$ is flat over $B_{\fn}$ and $\widehat A$ is flat over $\widehat{B_{\fn}}$, so $\widehat A$ is flat over $B_{\fn}$. Also $A\to\widehat A$ is faithfully flat. For any $B_{\fn}$-module $N$ and $i\ge1$: $\operatorname{Tor}_i^{B_{\fn}}(A,N)\otimes_A\widehat A=\operatorname{Tor}_i^{B_{\fn}}(\widehat A,N)=0$, hence $\operatorname{Tor}_i^{B_{\fn}}(A,N)=0$. So $A$ is flat over $B_{\fn}$, i.e.\ $F_{X/k}$ is flat at $x$.

(c) \emph{Flatness.} Flatness of $F_{X/k}$ at a point is local on $X$, and $F_{X/k}^{-1}(U^{(p)})=U$ for $U\subseteq X$ open (since $F_X$ is the identity on points and $W\circ F_{X/k}=F_X$, with $W:X^{(p)}\to X$ the first projection, which is a homeomorphism by Lemma~2.24); so we may replace $X$ by a non-empty affine open subscheme $U$, hence assume $X$ of finite type over $k$ (so $X$, $X^{(p)}$ are quasi-compact). First let $K=\bar k$ and $x$ be a closed point of $X$. The relative Frobenius commutes with base change (Ex.~3.2.20(a)), $F_{X_K/K}=F_{X/k}\times_kK$ (with $(X_K)^{(p)}=(X^{(p)})_K$). Let $\pi:X_K\to X$ and $\pi':(X^{(p)})_K\to X^{(p)}$ be the projections, flat, with $\pi'\circ F_{X_K/K}=F_{X/k}\circ\pi$. Choose a point $x'\in X_K$ over $x$ with $k(x')=K$: the fibre $\operatorname{Spec}(k(x)\otimes_kK)$ is non-empty and finite over $K$, and its points have residue field $K$ (as $k(x)$ is finite over $k$ and $K$ is algebraically closed). Then $X_K$ is smooth over $K$ (base change), and $x'$ is a rational point, so by (b) $F_{X_K/K}$ is flat at $x'$. Let $y=F_{X/k}(x)$, $y'=F_{X_K/K}(x')$, $A=\mathcal O_{X^{(p)},y}$, $A'=\mathcal O_{(X^{(p)})_K,y'}$, $B=\mathcal O_{X,x}$, $B'=\mathcal O_{X_K,x'}$. The maps $A\to A'$ and $B\to B'$ are flat local homomorphisms (so $B\to B'$ is faithfully flat), and $A'\to B'$ is flat. Hence $B'$ is flat over $A$ (composition $A\to A'\to B'$). Then for any $A$-module $N$, $\operatorname{Tor}_i^A(B,N)\otimes_BB'=\operatorname{Tor}_i^A(B',N)=0$, and faithful flatness of $B\to B'$ gives $\operatorname{Tor}_i^A(B,N)=0$. So $F_{X/k}$ is flat at every closed point $x$ of $X$.

Next, $F_{X/k}$ is a morphism of finite type between $k$-varieties, and maps closed points to closed points: if $x$ is closed, $k(x)$ is finite over $k$, so $k(F(x))\subseteq k(x)$ is finite over $k$, so $F(x)$ is closed. Thus Ex.~3.2 (in the corrected form proved above: $X^{(p)}$ is quasi-compact and closed points of $X$ map to closed points) shows that $F_{X/k}$ is flat everywhere.

\emph{Surjectivity.} (This also follows from Lemma~2.24, which says that $F_{X/k}$ is a homeomorphism; here is a direct argument.) Let $W:X^{(p)}\to X$ be the first projection. Then $W\circ F_{X/k}=F_X$, the absolute Frobenius of $X$, which is the identity on the underlying topological space; so $F_{X/k}$ is injective and $W$ is surjective on points. We claim $W$ is injective on points. The fibre of $W$ over $x\in X$ is $\operatorname{Spec}(k(x)\otimes_kk')$. In the ring $C=k(x)\otimes_kk'$, for any element $z=\sum a_i\otimes\mu_i$ we have $z^p=\sum a_i^p\otimes\mu_i^p=\sum a_i^p\mu_i\otimes1\in k(x)\otimes1$ (because $\mu^p=F(\mu)$ is the image in $k'$ of $\mu\in k$). So for a prime $\fq$ of $C$ and $z\in\fq$ we get $z^p\in\fq\cap(k(x)\otimes1)=0$, as $k(x)\otimes1$ is a field. Hence every prime of $C$ equals the nilradical of $C$, and $C$ has a single prime. So $W$ is injective. Now for $y\in X^{(p)}$ choose $x$ with $F_X(x)=W(y)$ (as $F_X$ is bijective); then $W(F_{X/k}(x))=F_X(x)=W(y)$ and injectivity of $W$ gives $y=F_{X/k}(x)$. So $F_{X/k}$ is surjective, hence faithfully flat.

(d) Here $p=2$, $A=k[x,y]/(y^2+x^3)$ and $Y=\operatorname{Spec}A$. Applying $F$ to the coefficients of the defining equation (which lie in $\mathbb F_2$), $A\otimes_kk'=k[x',y']/(y'^2+x'^3)=:B$ with $x'=x\otimes1$, $y'=y\otimes1$, and $\Phi:B\to A$, $x'\mapsto x^2$, $y'\mapsto y^2$ (consistent: $y^4+x^6=(y^2+x^3)^2=0$). The ring $A$ is isomorphic to $k[\tau^2,\tau^3]\subset k[\tau]$ via $x\mapsto\tau^2$, $y\mapsto\tau^3$ (characteristic $2$ so signs are irrelevant; $A$ is a domain, the cuspidal cubic). Then the image of $\Phi$ is $k[\tau^4,\tau^6]\simeq B$ (so $\Phi$ is injective), $A$ is a finite $B$-module (generated by $1,x,y,xy$; indeed $x^2$, $y^2$ lie in the image of $\Phi$), and the fraction fields are $\operatorname{Frac}A=k(\tau)$, $\operatorname{Frac}B=k(\tau^4,\tau^6)=k(\tau^2)$, so $[\operatorname{Frac}A:\operatorname{Frac}B]=2$.

Suppose $F_{Y/k}$ were flat at the origin $o$, $\fn=(x',y')\subset B$. The only prime of $A$ above $\fn$ is $\fm=(x,y)$, because $A/\fn A=k[x,y]/(y^2+x^3,x^2,y^2)=k[x,y]/(x^2,y^2)$ is local (this uses $x^3\in(x^2)$). So $A_{\fn}:=A\otimes_BB_{\fn}=A_{\fm}$ is a finite flat module over the Noetherian local ring $B_{\fn}$, hence free, say of rank $r$. Computing the rank at the closed point: $r=\dim_kA_{\fn}/\fn A_{\fn}=\dim_kk[x,y]/(x^2,y^2)=4$. Computing the rank at the generic point: $A_{\fn}\otimes_{B_{\fn}}\operatorname{Frac}B$ is a domain finite over the field $\operatorname{Frac}B$, hence a field containing $A$, and equals $\operatorname{Frac}A$; so $r=2$. Contradiction. Therefore $F_{Y/k}$ is not flat at $o$. (Of course $Y$ is not smooth at $o$, in accordance with (c).) \hfill$\square$

Citations without a chapter prefix (``Thm.~3.12'', ``Cor.~2.12'', \dots) refer to Chapter~4 of Liu; results of other chapters are cited as ``Thm.~1.3.16'' (Chapter~1, Section~3) and so on. ``Homomorphism of local rings'' always means a local homomorphism.

\exercise{3.14}{Let us keep the notation and hypotheses of Lemma 3.16 (so $A\to B$ is a flat homomorphism of Noetherian local rings and $b\in B$ is an element whose image in $B/\fm_AB$ is not a zero divisor). Let $J=\Ann(b)$.
(a) Show that $bB$ is flat over $A$.
(b) Using the exact sequence $0\to J\to B\to B$, where the last map is multiplication by $b$, show that $J\otimes_A A/\fm_A=0$. Deduce from this that $J=0$ and that $b$ does not divide zero in $B$.}
\pf
Put $k=A/\fm_A$. The homomorphism $A\to B$ is local, so $\fm_AB\subseteq\fm_B$.

(a) Consider the exact sequence of $A$-modules $0\to bB\to B\to B/bB\to0$. By Lemma 3.16, $B/bB$ is flat over $A$, and $B$ is flat by hypothesis. For every $A$-module $N$, the long exact sequence of $\Tor^A(-,N)$ contains the segment
\[ \Tor_2^A(B/bB,N)\to\Tor_1^A(bB,N)\to\Tor_1^A(B,N). \]
The outer terms vanish (a flat module has vanishing higher $\Tor$), hence $\Tor_1^A(bB,N)=0$ for all $N$. A module $P$ is flat if and only if $\Tor_1^A(N,P)=0$ for all $A$-modules $N$; hence $bB$ is flat over $A$.

(b) Multiplication by $b$ factors as $B\xrightarrow{\,u\,}bB\xrightarrow{\,\iota\,}B$, where $u$ is onto, $\iota$ is the inclusion, and $\Ker u=J$. So we have an exact sequence $0\to J\to B\xrightarrow{u}bB\to0$. Tensoring with $k$ and using $\Tor_1^A(bB,k)=0$ (part (a)) gives an exact sequence
\[ 0\to J\otimes_Ak\to B\otimes_Ak\xrightarrow{\,u\otimes 1\,}bB\otimes_Ak. \]
The composite $(\iota\otimes1)\circ(u\otimes1)$ is multiplication by $b$ on $B\otimes_Ak=B/\fm_AB$, which is injective by hypothesis. Hence $u\otimes1$ is injective, so its kernel, which is the image of the injective map $J\otimes_Ak\to B\otimes_Ak$, is $0$. Therefore $J\otimes_Ak=J/\fm_AJ=0$.

Thus $J=\fm_AJ\subseteq\fm_BJ$. As $B$ is Noetherian, $J$ is a finitely generated $B$-module, and Nakayama's lemma gives $J=0$. That is, $\Ann(b)=0$, so $b$ is not a zero divisor of $B$. \hfill$\square$

\exercise{3.15}{Let $X,Y$ be integral schemes, separated and of finite type over a locally Noetherian scheme $S$. Let $f:X\dashrightarrow Y$ be a birational map. Let $T\to S$ be a faithfully flat morphism such that $X_T$ and $Y_T$ are integral. Suppose that $f_T:X_T\dashrightarrow Y_T$ is defined everywhere and that $Y$ is affine. We want to show that $f$ is also defined everywhere.
(a) Show that we can suppose $S,T$ are affine (and even local).
(b) Let us identify the function fields $K(X)$ and $K(Y)$. Show that $f$ is defined everywhere if and only if $\mathcal O_Y(Y)\subseteq\mathcal O_X(X)$.
(c) Let $M=(\mathcal O_X(X)+\mathcal O_Y(Y))/\mathcal O_X(X)$. Show that $M\otimes_{\mathcal O_S(S)}\mathcal O_T(T)=0$. Deduce from this that $M=0$ and that $f$ is defined everywhere.}
\pf
\textbf{Conventions.} The birational map $f$ is an isomorphism of $S$-schemes $U\to V$ between dense open subschemes of $X$ and $Y$; through it we identify $K:=K(X)=K(Y)$. For a nonempty open $W\subseteq X$ (resp.\ $W\subseteq Y$) the ring $\mathcal O_X(W)$ (resp.\ $\mathcal O_Y(W)$) is a subring of $K$, since $X,Y$ are integral. The map $f$ is \emph{defined at} $x\in X$ if some morphism $W\to Y$, with $W\ni x$ open, agrees with $f$ on a dense open subset of $W\cap U$; it is \emph{defined everywhere} if it extends to a morphism $X\to Y$. Since $Y$ is separated over $S$ and $X$ is reduced, two such local extensions agree on overlaps (they agree on a dense open and $Y$ is separated), so they glue: $f$ is defined everywhere if and only if it is defined at every point. The base change $f_T$ is the isomorphism $U_T\to V_T$; it is a map of $T$-schemes, and ``defined everywhere'' for $f_T$ means that it extends to a $T$-morphism $X_T\to Y_T$. Since $X_T$, $Y_T$ are integral and nonempty, $U_T$ and $V_T$ are dense open in them.

We prove (b), then (c) in the affine case, then the reduction (a). This order is the logical one.

\textbf{(b)} We prove it for any integral schemes $X',Y'$ with $Y'$ affine and any birational map $f':X'\dashrightarrow Y'$ (an isomorphism $U'\to V'$ of dense opens), $K$ being the common function field.

$(\Rightarrow)$ Let $g:X'\to Y'$ be an extension of $f'$ and $C=\mathcal O_{Y'}(Y')$. For $c\in C$ the element $g^\#(c)\in\mathcal O_{X'}(X')\subseteq K$ restricts on a dense open subset $U''\subseteq U'$ to $(g|_{U''})^\#(c)=(f'|_{U''})^\#(c|_{V'})$, and the latter is the element $c\in K$ by the very definition of the identification $K(Y')=K(X')$. As $\mathcal O_{X'}(U'')\to K$ is injective, $c=g^\#(c)\in\mathcal O_{X'}(X')$. So $C\subseteq\mathcal O_{X'}(X')$.

$(\Leftarrow)$ Suppose $C\subseteq\mathcal O_{X'}(X')$. The inclusion is a ring homomorphism $C\to\mathcal O_{X'}(X')$, hence defines a morphism $g:X'\to\operatorname{Spec}C=Y'$. On $U'$, both $g|_{U'}$ and $f'$ are morphisms into the affine scheme $Y'$, so they correspond to ring homomorphisms $C\to\mathcal O_{X'}(U')\subseteq K$. Both are $c\mapsto c$ (the inclusion $C\subseteq K$). Hence $g|_{U'}=f'$ and $f'$ is defined everywhere.

\textbf{(c), the affine case.} Let $V_0=\operatorname{Spec}R$ be an affine open subset of $S$, and put $X_0=X_{V_0}$, $Y_0=Y_{V_0}$. We assume that $Y_0$ is nonempty and affine. Put $A_X=\mathcal O(X_0)$, $A_Y=\mathcal O(Y_0)$ (subrings of $K$), $N=A_X+A_Y\subseteq K$ and $M=N/A_X$; these are $R$-modules. Let $T'=\operatorname{Spec}R'$ be an affine open subset of $T$ lying over $V_0$ such that $X_{T'}\neq\emptyset$, and assume that $f_{T'}$ is defined everywhere. We claim that
\[ M\otimes_RR'=0. \]
$X_{T'}$ and $Y_{T'}$ are nonempty open subschemes of the integral schemes $X_T,Y_T$, hence integral; $K':=K(X_{T'})=K(Y_{T'})$ via $f_{T'}$. The schemes $X_0$ and $Y_0$ are of finite type over the Noetherian affine scheme $V_0$, hence quasi-compact and quasi-separated, and so are all their open subsets. Therefore flat base change for global sections gives, for every open $W\subseteq X_0$,
\[ \mathcal O(W_{T'})=\mathcal O(W)\otimes_RR',\qquad \mathcal O(Y_{T'})=A_Y\otimes_RR'. \]
(For $X_0$ write $\mathcal O(W)$ as the kernel of a map between finite products of rings of affine opens, and use that $R'$ is flat over $R$.)

\emph{Step 1: an injection $K\otimes_RR'\hookrightarrow K'$.} We have $K=\varinjlim_W\mathcal O(W)$, over the nonempty opens $W\subseteq X_0$, so $K\otimes_RR'=\varinjlim_W\mathcal O(W_{T'})$. Each $W_{T'}$ is a nonempty open subscheme of the integral scheme $X_{T'}$. (It is nonempty because $X_{T'}\to X_0$ is flat; its image is stable under generization and so contains the generic point of $X_0$, which lies in $W$.) So the maps $\mathcal O(W_{T'})\to K'$ are injective and compatible, and they induce an injection
\[ j:K\otimes_RR'\hookrightarrow K'. \]
By construction $j(A_X\otimes R')=\mathcal O(X_{T'})$. Moreover, $j(A_Y\otimes R')=\mathcal O(Y_{T'})$: an element $c\in A_Y$ corresponds, under $f:U\to V$, to an element of $\mathcal O(U\cap X_0)\subseteq K$, and the element $c\otimes r$ of $\mathcal O(Y_{T'})$ restricts on $V_{T'}$ to the element of $\mathcal O(V_{T'})$ that corresponds to $c\otimes r\in\mathcal O((U\cap X_0)_{T'})$ under the isomorphism $f_{T'}$. These are the identifications used to define $K'$.

\emph{Step 2.} Since $R'$ is flat, $N\otimes_RR'\to K\otimes_RR'$ is injective, hence so is $\psi=j\circ(\text{this}):N\otimes_RR'\to K'$. Because $N=A_X+A_Y$, the image of $\psi$ is $\mathcal O(X_{T'})+\mathcal O(Y_{T'})$. By hypothesis $f_{T'}$ is defined everywhere and $Y_{T'}$ is affine, so by (b) $\mathcal O(Y_{T'})\subseteq\mathcal O(X_{T'})$. Hence $\psi(N\otimes R')=\mathcal O(X_{T'})=\psi(A_X\otimes R')$. As $\psi$ is injective, the injection $A_X\otimes R'\to N\otimes R'$ is onto. By flatness, $M\otimes_RR'=(N\otimes R')/(A_X\otimes R')=0$. This proves the claim.

\textbf{(a) The reduction, and the end of the proof.} Here we use that the preimage in the affine scheme $Y$ of an affine open subset of $S$ is again affine; this is automatic if $S$ is separated (e.g.\ affine, or quasi-projective over a field or $\mathbf Z$), which we assume for this step. Let $x\in X$ with image $s\in S$ and take an affine open neighborhood $V_0=\operatorname{Spec}R$ of $s$, with $X_0,Y_0,A_X,A_Y,M$ as above. Let $\sigma\in S$ be the image of the generic point of $X$; it is also the image of the generic point of $Y$ (as $f$ is an $S$-isomorphism $U\cong V$), and $\sigma$ is a generization of $s$. Hence $\sigma\in V_0$, so $X_0,Y_0$ are nonempty, and $f$ restricts to a birational map $f_0:X_0\dashrightarrow Y_0$. It suffices to show $A_Y\subseteq A_X$, i.e.\ $M=0$: indeed, then (b) applied to $f_0$ shows that $f_0$ is defined everywhere, in particular at $x$, so that $f$ is defined at $x$; and $x$ is arbitrary.

Let $\fp=\Ker(R\to K)$, a prime of $R$ (the point $\sigma$). Both $A_X$ and $A_Y$ are $R/\fp$-modules, so $\fp M=0$ and $M_\fq=0$ for every prime $\fq\not\supseteq\fp$. Let $\fq\supseteq\fp$ be a prime of $R$. By surjectivity of $T\to S$ there is $t\in T$ over $\fq$; let $T'=\operatorname{Spec}R'$ be an affine open neighborhood of $t$ lying over $V_0$. Since $T'\to V_0$ is flat, going down gives $t'\in T'$ over $\fp$; the generic point $\xi$ of $X_0$ lies over $\fp$, and $k(\xi)\otimes_{k(\fp)}k(t')\neq0$, so $X_{T'}\neq\emptyset$. Moreover $f_{T'}$ is defined everywhere: the $T$-morphism $X_T\to Y_T$ extending $f_T$ maps $X_{T'}$ into $Y_{T'}$. By the affine case, $M\otimes_RR'=0$, hence $M\otimes_RO_{T,t}=0$. Now $R_\fq\to O_{T,t}$ is a flat local homomorphism, hence faithfully flat, and $M\otimes_RO_{T,t}=M_\fq\otimes_{R_\fq}O_{T,t}$. Therefore $M_\fq=0$. So $M_\fq=0$ for all primes $\fq$ of $R$, i.e.\ $M=0$ (Lemma~1.2.12). \hfill$\square$

\exercise{3.16}{Let $B$ be a semi-local ring, that is, having only a finite number of maximal ideals $\fm_1,\dots,\fm_n$. Let $\widehat B$ be the formal completion of $B$ for the topology defined by $\bigcap_i\fm_i$. Show that $\widehat B\simeq\bigoplus_i\widehat B_i$, where $\widehat B_i$ is the formal completion of $B_{\fm_i}$ for the topology defined by its maximal ideal.}
\pf
No Noetherian hypothesis is needed. Let the $\fm_i$ be pairwise distinct, $J=\bigcap_i\fm_i$ and $P=\fm_1\cdots\fm_n$. By definition $\widehat B=\varprojlim_kB/J^k$.

\emph{Step 1: the $J$-adic and $P$-adic topologies agree.} Since $P\subseteq\bigcap_i\fm_i=J$, we have $P^k\subseteq J^k$. Conversely $J^n\subseteq P$: a product $a_1\cdots a_n$ of $n$ elements of $J$ lies in $\fm_1\cdots\fm_n$, because $a_i\in\fm_i$. Therefore $J^{nk}\subseteq P^k\subseteq J^k$ for all $k\ge1$. These inclusions give mutually inverse isomorphisms between $\varprojlim_kB/J^k$ and $\varprojlim_kB/P^k$ (a map of inverse systems in each direction after reindexing, composing to the canonical transition maps). So $\widehat B=\varprojlim_kB/P^k$.

\emph{Step 2: Chinese remainder theorem.} For $i\ne j$ we have $\fm_i+\fm_j=B$, hence $\fm_i^k+\fm_j^k=B$ for all $k$. For pairwise comaximal ideals the intersection equals the product, so $\bigcap_i\fm_i^k=\prod_i\fm_i^k=P^k$, and the Chinese remainder theorem gives a canonical isomorphism, compatible with $k$,
\[ B/P^k\;\xrightarrow{\ \sim\ }\;\prod_{i=1}^nB/\fm_i^k. \]
Passing to inverse limits, $\widehat B\simeq\prod_i\varprojlim_kB/\fm_i^k$.

\emph{Step 3: $B/\fm_i^k=B_{\fm_i}/\fm_i^kB_{\fm_i}$.} The ring $B/\fm_i^k$ is local with maximal ideal $\fm_i/\fm_i^k$: if $s\notin\fm_i$ there are $a\in B$ and $m\in\fm_i$ with $as+m=1$, so $as\equiv1-m$, which is a unit modulo $\fm_i^k$ (its inverse is $1+m+\dots+m^{k-1}$). Hence every element of $B\setminus\fm_i$ is already invertible in $B/\fm_i^k$, and localization at $\fm_i$ does not change this ring: $B/\fm_i^k\to B_{\fm_i}/\fm_i^kB_{\fm_i}$ is an isomorphism (compatible with $k$). Hence $\varprojlim_kB/\fm_i^k=\widehat B_i$, and $\widehat B\simeq\bigoplus_i\widehat B_i$ (a finite product is a finite direct sum). \hfill$\square$

\exercise{3.17}{Let $A$ be a Noetherian local ring, and let $B$ be a finite $A$-algebra with maximal ideals $\fm_1,\dots,\fm_n$. Show that $B\otimes_A\widehat A\simeq\bigoplus_i\widehat{B_{\fm_i}}$, where $\widehat{\ }$ denotes the formal completion for the topology defined by the maximal ideal.}
\pf
Let $\fm=\fm_A$ and $J=\bigcap_i\fm_i$. Since $B$ is a finitely generated $A$-module and $A$ is Noetherian, Cor.~1.3.14 gives a canonical isomorphism
\[ B\otimes_A\widehat A\;\simeq\;\varprojlim_kB/\fm^kB, \]
the completion of the $A$-module $B$ for the $\fm$-adic topology.

\emph{The topologies $(\fm^kB)_k$ and $(J^k)_k$ coincide.} Since $B$ is integral over $A$, each maximal ideal $\fm_i$ contracts to a maximal ideal of $A$ (if $B/\fm_i$ is a field integral over the subring $A/(\fm_i\cap A)$, the latter is a field), hence $\fm_i\cap A=\fm$ and $\fm B\subseteq\fm_i$ for all $i$. So $\fm B\subseteq J$ and $\fm^kB\subseteq J^k$. Conversely, $B/\fm B$ is a finite-dimensional algebra over $A/\fm$, hence an Artinian ring, and its Jacobson radical $J/\fm B$ is nilpotent. So $J^N\subseteq\fm B$ for some $N$, and $J^{Nk}\subseteq\fm^kB$. As in Exercise 3.16, Step~1, $\varprojlim_kB/\fm^kB=\varprojlim_kB/J^k$, the formal completion of the semi-local ring $B$ for the topology defined by $J=\bigcap_i\fm_i$.

By Exercise 3.16 this is $\bigoplus_i\widehat{B_{\fm_i}}$. The resulting isomorphism $B\otimes_A\widehat A\to\bigoplus_i\widehat{B_{\fm_i}}$ is canonical: it sends $b\otimes a$ to $(a\cdot\hat b_i)_i$, where $\hat b_i$ is the image of $b$ in $\widehat{B_{\fm_i}}$ and $\widehat A\to\widehat{B_{\fm_i}}$ is induced by $A\to B_{\fm_i}$. \hfill$\square$

\exercise{3.18}{Let $X$ be a quasi-projective scheme over a locally Noetherian scheme $S$. Let $G$ be a finite group acting on $X$. Hence the quotient $Y:=X/G$ exists (Exercise 3.3.23).
(a) Let $G$ act on an $A$-algebra $B$, and $C$ be a flat $A$-algebra. Then $G$ acts on $B\otimes_AC$ by the identity on the second component. Let us denote the $A$-linear map $g-\operatorname{Id}_B:B\to B$ by $\varphi_g$. Show that $\Ker(\varphi_g\otimes\operatorname{Id}_C)=\Ker(\varphi_g)\otimes_AC$ and that
$\Ker(\varphi_g\otimes\operatorname{Id}_C)\cap\Ker(\varphi_{g'}\otimes\operatorname{Id}_C)=(\Ker(\varphi_g)\cap\Ker(\varphi_{g'}))\otimes_AC$ in $B\otimes_AC$. Deduce from this that $(B\otimes_AC)^G=B^G\otimes_AC$. Show that the quotient morphism $f:X\to Y$ commutes with flat base change.
(b) Let us fix a point $x\in X$. Let us set $D:=\{g\in G\mid gx=x\}$. This is the decomposition group of $x$. Show that $D$ acts canonically on $\mathcal O_{X,x}$ as well as on the formal completion $\widehat{\mathcal O}_{X,x}$. Let $y=f(x)$. Show that $(\widehat{\mathcal O}_{X,x})^D=\widehat{\mathcal O}_{Y,y}$ (use Exercise 3.17 and (a)). Show that $\mathcal O_{Y,y}=(\mathcal O_{X,x})^D$ if $G=D$.
(c) Show that the morphism $f:X\to Y$ factors into $X\to X/D$ and $X/D\to Y$. Let $z$ be the image of $x$ in $X/D$. Show that $X/D\to Y$ is etale at $z$.}
\pf
\textbf{Recollections (Exercises 2.3.21 and 3.3.23).} We use the following facts about $f:X\to Y=X/G$ .
(Q1) Every $G$-orbit of $X$ is contained in a $G$-stable affine open subset $U=\operatorname{Spec}B$. Then $f(U)=\operatorname{Spec}B^G$ is open in $Y$, $f^{-1}(f(U))=U$, and such $f(U)$ cover $Y$.
(Q2) $f$ is finite and surjective, $G$ acts transitively on the fibers of $f$, and $Y$ is of finite type over $S$. In the situation of (Q1), $B^G$ is Noetherian and $B$ is a finite $B^G$-module (each $b\in B$ is a root of $\prod_{g}(T-gb)\in B^G[T]$; Artin--Tate).
(Q3) (universal property) every $G$-invariant morphism $X\to Z$ factors uniquely through $f$.
(Q4) For every subgroup $H\subseteq G$ the quotient $X/H$ exists and has the same properties.
We let $G$ act on rings of $G$-stable affine opens by $g\cdot$ so that the prime of $B$ attached to $gx$ is $g$ applied to the prime attached to $x$.

\textbf{(a)} Since $C$ is $A$-flat, $-\otimes_AC$ preserves kernels. The sequence $0\to\Ker\varphi_g\to B\xrightarrow{\varphi_g}B$ is exact, so $0\to\Ker\varphi_g\otimes C\to B\otimes C\xrightarrow{\varphi_g\otimes1}B\otimes C$ is exact, i.e.\ $\Ker(\varphi_g\otimes1)=\Ker(\varphi_g)\otimes_AC$ (a submodule of $B\otimes C$ by flatness). More generally let $g_1,\dots,g_r\in G$ and consider $\Phi=(\varphi_{g_1},\dots,\varphi_{g_r}):B\to B^r$. Then $\Ker\Phi=\bigcap_i\Ker\varphi_{g_i}$ and $\Phi\otimes1:B\otimes C\to(B\otimes C)^r$ is $(\varphi_{g_i}\otimes1)_i$, so $\Ker(\Phi\otimes1)=\bigcap_i\Ker(\varphi_{g_i}\otimes1)$. Exactness of $0\to\Ker\Phi\to B\to B^r$ and flatness give $\bigcap_i\Ker(\varphi_{g_i}\otimes1)=(\bigcap_i\Ker\varphi_{g_i})\otimes_AC$. For $r=2$ this is the second assertion. Taking for $g_1,\dots,g_r$ all the elements of $G$, we get $(B\otimes C)^G=\bigcap_g\Ker(\varphi_g\otimes1)=(\bigcap_g\Ker\varphi_g)\otimes C=B^G\otimes_AC$.

\emph{Flat base change.} Let $T\to S$ be a flat morphism. $G$ acts on $X_T=X\times_ST$ through $X$, and $f_T:X_T\to Y_T$ is $G$-invariant. Let $U=\operatorname{Spec}B$ be as in (Q1) for a point $x\in X$ with image $s\in S$. Choose an affine open neighborhood $\operatorname{Spec}A_0$ of $s$ and let $W\subseteq U$ be the preimage of $\operatorname{Spec}A_0$ in $U$; it is open, $G$-stable ($G$ acts by $S$-automorphisms) and contains $Gx$. Let $J$ be the ideal of $B$ defining $U\setminus W$. By prime avoidance there is $h\in J$ not in any of the finitely many primes of $Gx$; then $N(h)=\prod_{g\in G}g(h)\in B^G$ and $D(N(h))\subseteq D(h)\subseteq W$ is a $G$-stable principal open subset of $U$ containing $Gx$. Replacing $U$ by it, we may assume that $U=\operatorname{Spec}B$ lies over $\operatorname{Spec}A_0$ (it is still a $G$-stable affine open, with the properties of (Q1)). Let $\operatorname{Spec}C$ be an affine open of $T$ over $\operatorname{Spec}A_0$. Then $C$ is a flat $A_0$-algebra, $U\times_S\operatorname{Spec}C=\operatorname{Spec}(B\otimes_{A_0}C)$ is a $G$-stable affine open subset of $X_T$, and by the above its quotient is $\operatorname{Spec}((B\otimes C)^G)=\operatorname{Spec}(B^G\otimes_{A_0}C)=f(U)\times_S\operatorname{Spec}C$. These $G$-stable affine opens cover $X_T$ and every point has a $G$-stable affine neighborhood, so by Exercise 2.3.21(c) the quotient $X_T/G$ exists and is obtained by gluing these quotients, i.e.\ $X_T/G=Y_T$ with quotient morphism $f_T$. This says that the formation of the quotient commutes with flat base change.

\textbf{(b)} \emph{Action.} An element $g\in D$ is an automorphism of $X$ fixing $x$, so it induces an automorphism of the local ring $\mathcal O_{X,x}$, preserving $\fm_x$. It is therefore continuous for the $\fm_x$-adic topology and extends to $\widehat{\mathcal O}_{X,x}$. Because $G$ acts by $S$-automorphisms, $D$ also acts on $k(x)$, trivially on $k(y)$ (the action on $k(x)$ is the one used in Exercise 3.19).

\emph{Invariants.} Take $U=\operatorname{Spec}B\supseteq Gx$ as in (Q1). Let $\fq\subset B^G$ be the prime ideal of $y$, $\Sigma=B^G\setminus\fq$, $A'=\Sigma^{-1}B^G=\mathcal O_{Y,y}$ (a Noetherian local ring) and $B'=\Sigma^{-1}B$. The group $G$ acts on $B'$ by $A'$-algebra automorphisms and $(B')^G=A'$: if $b/s$ is $G$-invariant with $s\in\Sigma$, then for each $g$ there is $t_g\in\Sigma$ with $t_g(gb-b)=0$; with $t=\prod_gt_g\in\Sigma$ we get $tb\in B^G$, and $b/s=tb/ts\in A'$. By (Q2), $B'$ is a finite $A'$-algebra. Its maximal ideals are the primes of $B$ lying over $\fq$, i.e.\ the points of $f^{-1}(y)=Gx$ (they lie over $\fm_{A'}$ since $B'$ is finite over the local ring $A'$), and the localization of $B'$ at the maximal ideal of $x'$ is $\mathcal O_{X,x'}$. By Exercise 3.17 there is a canonical isomorphism
\[ \Phi:\ B'\otimes_{A'}\widehat{A'}\ \xrightarrow{\ \sim\ }\ \bigoplus_{x'\in Gx}\widehat{\mathcal O}_{X,x'},\qquad b\otimes a\mapsto\bigl(a\,\hat b_{x'}\bigr)_{x'}. \]
It is $G$-equivariant, where $G$ acts on the left side through $B'$ and on the right side by the isomorphisms $g:\widehat{\mathcal O}_{X,x'}\to\widehat{\mathcal O}_{X,gx'}$, permuting the summands (this is clear on pure tensors, as $g$ is $A'$-linear). The ring $\widehat{A'}=\widehat{\mathcal O}_{Y,y}$ is flat over $A'$ (Thm.~1.3.15), so by (a)
\[ (B'\otimes_{A'}\widehat{A'})^G=(B')^G\otimes_{A'}\widehat{A'}=\widehat{A'}. \]
On the right side, a family $(u_{x'})_{x'}$ is $G$-invariant if and only if $u_{gx}=g\,u_x$ for all $g$; this forces $u_x\in(\widehat{\mathcal O}_{X,x})^D$, and conversely every $u\in(\widehat{\mathcal O}_{X,x})^D$ gives the invariant family $(g\,u)_{gx}$, which is well defined because $u$ is $D$-invariant. So the projection to the $x$-component is a bijection from the $G$-invariants of $\bigoplus_{x'}\widehat{\mathcal O}_{X,x'}$ onto $(\widehat{\mathcal O}_{X,x})^D$. Composing with $\Phi$ and with the identification $\widehat{A'}=(B'\otimes\widehat{A'})^G$, the element $a\in\widehat{A'}$ goes to the image of $a$ under the map $\rho:\widehat{\mathcal O}_{Y,y}\to\widehat{\mathcal O}_{X,x}$ induced by $f$. We have shown:
\par\emph{$\rho$ is injective, and its image is exactly $(\widehat{\mathcal O}_{X,x})^D$.}\par
This is the statement $(\widehat{\mathcal O}_{X,x})^D=\widehat{\mathcal O}_{Y,y}$ (the equality being realized by $\rho$).

\emph{The case $G=D$.} Then $Gx=\{x\}$, so $B'$ is local, and $B'=\mathcal O_{X,x}$ (a local ring is its own localization at its maximal ideal). Hence $\mathcal O_{Y,y}=A'=(B')^G=(\mathcal O_{X,x})^D$; moreover $\mathcal O_{X,x}$ is finite over $\mathcal O_{Y,y}$.

\textbf{(c)} $D$ is a finite group acting on $X$, so $X/D$ exists (Q4), and $f$ is $D$-invariant because it is $G$-invariant. By (Q3) it factors uniquely as $X\xrightarrow{\pi}X/D\xrightarrow{h}Y$. Both $X/D$ and $Y$ are of finite type over $S$, so locally Noetherian, and $h$ is of finite type (indeed finite). Let $z=\pi(x)$.

Apply (b) to the group $D$ acting on $X$: the decomposition group of $x$ in $D$ is $D$ itself, so $\widehat{\mathcal O}_{X/D,z}\xrightarrow{\ \rho_D\ }\widehat{\mathcal O}_{X,x}$ is injective with image $(\widehat{\mathcal O}_{X,x})^D$. Apply (b) to $G$: $\widehat{\mathcal O}_{Y,y}\xrightarrow{\rho}\widehat{\mathcal O}_{X,x}$ is injective with the same image. Since $\rho=\rho_D\circ\hat h^\#$ (with $\hat h^\#:\widehat{\mathcal O}_{Y,y}\to\widehat{\mathcal O}_{X/D,z}$ induced by $h$), the map $\hat h^\#$ is injective and has the same image as $\rho$ after composing with the injection $\rho_D$; hence $\hat h^\#$ is surjective. So
\[ \hat h^\#:\widehat{\mathcal O}_{Y,y}\to\widehat{\mathcal O}_{X/D,z}\ \text{ is an isomorphism}. \]
The residue field of a Noetherian local ring and of its completion coincide, so $k(z)=k(y)$. By Prop.~3.26 (applied to $h$ at $z$), $h$ is etale at $z$. \hfill$\square$

\emph{Remark} (used in Exercise 3.19). We proved a bit more: $k(z)=k(y)$, and $\widehat{\mathcal O}_{Y,y}\to\widehat{\mathcal O}_{X/D,z}$ is an isomorphism; with etaleness this gives $\fm_y\mathcal O_{X/D,z}=\fm_z$ and flatness of $\mathcal O_{Y,y}\to\mathcal O_{X/D,z}$.

\exercise{3.19}{We keep the notation and hypotheses of the preceding exercise. Let $x\in X$. Then the decomposition group $D$ naturally acts on $k(x)$. This induces a group homomorphism $D\to\operatorname{Aut}(k(x)/k(y))$. The inertia group $I$ of $x$ is by definition the kernel of this homomorphism. We want to show that $f:X\to X/G$ is etale at $x$ if and only if $I$ acts trivially on $\mathcal O_{X,x}$.
(a) Show that $X\to X/G$ is etale at $x$ if and only if $X\to X/D$ is etale at $x$. In what follows we will therefore reduce to the case that $G=D$, that is, $f^{-1}(y)=\{x\}$. Let us write $A=\mathcal O_{Y,y}$ and $B=\mathcal O_{X,x}$. Show that $B^D=A$ and that $B$ is finite over $A$.
(b) Let $k'$ be the separable closure of $k(y)$ in $k(x)$. For any element $b\in B$, let $\bar b$ denote its image in $k(x)$. Let $\theta\in B$ be such that $k'=k(y)[\bar\theta]$. Show that $P(T):=\prod_{g\in D}(T-g(\theta))\in A[T]$ and that its image in $k(y)[T]$ vanishes in $\bar\theta$. Let $\tau\in\operatorname{Aut}(k'/k(y))$. Show that there exists a $g\in D$ such that $g(\bar\theta)=\tau(\bar\theta)$. Deduce from this that $\tau=\bar g$ and that $D\to\operatorname{Aut}(k(x)/k(y))=\operatorname{Aut}(k'/k(y))$ is surjective. Show that $k'/k(y)$ is Galois, and that we have an exact sequence of groups $1\to I\to D\to\operatorname{Gal}(k'/k(y))\to1$.
(c) Let $C=A[T]/(P(T))$. This is a finite etale local $A$-algebra. Let $D$ act on $C\otimes_AB$ as the identity on the first component. Show that $\operatorname{Spec}C=\operatorname{Spec}(C\otimes_AB)/D$, and that $B$ is etale over $A$ if and only if $C\otimes_AB$ is etale over $C$. Show that the decomposition group of every point of $\operatorname{Spec}(C\otimes_AB)_z$, where $z$ is the closed point of $\operatorname{Spec}C$, is equal to $I$. Hence, by (c), we have reduced to the case when $D=I$, and $k(x)$ is purely inseparable over $k(y)$.
(d) Using Proposition 3.26, show that $X\to Y$ is etale at $x$ if and only if $I$ acts trivially on $\mathcal O_{X,x}$. Show that in general $X/I\to Y$ is etale at the image of $x$ in $X/I$.}
\pf
\textbf{Conventions.} We use the facts (Q1)--(Q4) recalled in Exercise 3.18 and its results (a), (b), (c). Put $Y=X/G$, $y=f(x)$. For a local homomorphism $A\to B$ of Noetherian local rings (always the local rings of a morphism of finite type of locally Noetherian schemes) we say that $A\to B$ is \emph{etale} if (F) $B$ is flat over $A$, (U1) $\fm_AB=\fm_B$, and (U2) $k(B)/k(A)$ is a (finite) separable extension. By Def.~3.17, $X\to Y$ is etale at $x$ if and only if $\mathcal O_{Y,y}\to\mathcal O_{X,x}$ is etale in this sense; ``unramified'' means (U1)+(U2). The group $D$ acts on $k(x)$ through its action on $\mathcal O_{X,x}$ (Exercise 3.18(b)); it acts trivially on $k(y)$.

\textbf{Two lemmas.}

\emph{Lemma L.} Let $A\to C\to B$ be local homomorphisms of Noetherian local rings such that $A\to C$ is flat, $\fm_AC=\fm_C$ and $k(A)=k(C)$ (via $A\to C$). Then (i) $B$ satisfies (U1), (U2) over $A$ if and only if it does over $C$; (ii) $B$ is flat over $A$ if and only if it is flat over $C$. Hence $A\to B$ is etale if and only if $C\to B$ is etale.

\emph{Proof.} (i) $\fm_AB=(\fm_AC)B=\fm_CB$ and $k(B)/k(A)=k(B)/k(C)$. (ii) If $B$ is $C$-flat it is $A$-flat. Conversely let $B$ be $A$-flat. Put $\kappa=A/\fm_A$; then $C/\fm_C=C/\fm_AC=C\otimes_A\kappa$. If $F_\bullet\to\kappa$ is a free resolution over $A$, then $C\otimes_AF_\bullet$ is a free resolution of $C\otimes_A\kappa$ over $C$ ($C$ is $A$-flat), so
\[ \Tor_1^C(C/\fm_C,B)=H_1\bigl((C\otimes_AF_\bullet)\otimes_CB\bigr)=H_1(F_\bullet\otimes_AB)=\Tor_1^A(\kappa,B)=0. \]
By the local flatness criterion (the one used for Lemma~3.16: for a local homomorphism $C\to B$ of Noetherian local rings, $B$ is $C$-flat if and only if $\Tor_1^C(k(C),B)=0$), $B$ is flat over $C$. $\square$

\emph{Lemma E.} Let $A\to B$ be a local homomorphism of Noetherian local rings with $B$ finite over $A$; put $k_A=k(A)$, $k_B=k(B)$. Let $k'\subseteq k_B$ be a finite Galois extension of $k_A$, $k'=k_A(\bar\theta)$, with minimal polynomial $\bar m$ of $\bar\theta$ over $k_A$. Let $Q\in A[T]$ be any monic lift of $\bar m$, and put $C=A[T]/(Q)$, $B_1=C\otimes_AB=B[T]/(Q)$. Then:
(i) $C$ is a finite free $A$-module (hence faithfully flat), a local ring with maximal ideal $\fm_AC$ and residue field $k'$.
(ii) The maximal ideals of $B_1$ are the ideals $\fn_\sigma=\Ker(B_1\to k_B,\ b\mapsto\bar b,\ T\mapsto\sigma(\bar\theta))$, $\sigma\in\operatorname{Gal}(k'/k_A)$; they are pairwise distinct, lie over $\fm_B$, and have residue field $k_B$.
(iii) $A\to B$ is etale if and only if $C\to(B_1)_{\fn_\sigma}$ is etale for every $\sigma$.

\emph{Proof.} (i) $C$ is free with basis $1,T,\dots,T^{r-1}$, $r=\deg Q$, and $C/\fm_AC=k_A[T]/(\bar m)=k'$ is a field. Every maximal ideal of the finite $A$-algebra $C$ contains $\fm_A$, so $\fm_AC$ is the only maximal ideal.
(ii) $B_1$ is finite over $B$, so its maximal ideals contain $\fm_BB_1$ and correspond to the maximal ideals of $B_1/\fm_BB_1=k_B[T]/(\bar m)$. As $k'/k_A$ is Galois, $\bar m=\prod_\sigma(T-\sigma\bar\theta)$ with distinct roots in $k'\subseteq k_B$, so $k_B[T]/(\bar m)\simeq\prod_\sigma k_B$, with the evident maximal ideals.
(iii) Put $E=B/\fm_AB$ and note $B_1/\fm_AB_1=E\otimes_{k_A}k'$ and $\fm_AB_1=\fm_CB_1$.

\emph{Flatness.} If $B$ is $A$-flat, then $B_1$ is $C$-flat by base change, hence so are its localizations. Conversely suppose $(B_1)_{\fn_\sigma}$ is $C$-flat for every $\sigma$. These are all the localizations of $B_1$ at maximal ideals, so $B_1$ is $C$-flat by Cor.~1.2.15. Let $N'\hookrightarrow N$ be an injection of $A$-modules. Then $C\otimes N'\to C\otimes N$ is injective ($C$ is $A$-flat), hence so is $B_1\otimes_C(C\otimes N')\to B_1\otimes_C(C\otimes N)$, i.e.\ $(B\otimes N')\otimes_AC\to(B\otimes N)\otimes_AC$. As $C$ is faithfully flat over $A$, $B\otimes N'\to B\otimes N$ is injective. So $B$ is $A$-flat.

\emph{Unramifiedness.} Suppose $B$ is unramified over $A$. Then $E=k_B$ and $B_1/\fm_CB_1=k_B\otimes_{k_A}k'\simeq\prod_\sigma k_B$ (the same computation as in (ii), with the extension $k'/k_A$). A localization of this product of fields at a maximal ideal is the field $k_B$, so $\fm_C(B_1)_{\fn_\sigma}=\fn_\sigma(B_1)_{\fn_\sigma}$, and the residue extension is $k_B/k'$, separable because $k_B/k_A$ is. Conversely suppose all $(B_1)_{\fn_\sigma}$ are unramified over $C$. The Artinian ring $E\otimes_{k_A}k'=B_1/\fm_CB_1$ is the product of its localizations at its maximal ideals, namely $(B_1)_{\fn_\sigma}/\fm_C(B_1)_{\fn_\sigma}=k(\fn_\sigma)=k_B$, which are fields. So $E\otimes_{k_A}k'$ is reduced; since $E\to E\otimes_{k_A}k'$ is injective, $E$ is reduced. A reduced local Artinian ring is a field: $E=k_B$, i.e.\ $\fm_AB=\fm_B$. Finally $k_B=k(\fn_\sigma)$ is separable over $k'$ by hypothesis, and $k'/k_A$ is separable, so $k_B/k_A$ is separable. $\square$

\textbf{(a)} Let $z$ be the image of $x$ in $X/D$ and $h:X/D\to Y$ the induced morphism. Put $A_0=\mathcal O_{Y,y}$, $C_0=\mathcal O_{X/D,z}$, $B_0=\mathcal O_{X,x}$ (this $B_0$ is called $B$ below, once $G=D$). By Exercise 3.18(c) and its Remark, $A_0\to C_0$ is flat, $\fm_{A_0}C_0=\fm_{C_0}$ and $k(C_0)=k(A_0)$. By Lemma L, $A_0\to B_0$ is etale if and only if $C_0\to B_0$ is. In other words, $X\to X/G$ is etale at $x$ if and only if $X\to X/D$ is etale at $x$. The group $D$ acts on $X$, the point $x$ has decomposition group $D$ (and the same inertia group $I$) for this action, and the fiber of $X\to X/D$ over $z$ is the $D$-orbit $\{x\}$ of $x$. So it suffices to prove the theorem when $G=D$.

Assume now $G=D$ (that is, we replace $Y$ by $X/D$ and $y$ by $z$, which is legitimate by the above). By Exercise 3.18(b), $A:=\mathcal O_{Y,y}=B^D$ with $B=\mathcal O_{X,x}$, and $B$ is finite over $A$. In particular $k(x)$ is a finite extension of $k(y)$.

\textbf{(b)} Let $k'$ be the separable closure of $k(y)$ in the finite extension $k(x)$ and $\theta\in B$ a lift of a primitive element $\bar\theta$ of the finite separable extension $k'/k(y)$. The coefficients of $P(T)=\prod_{g\in D}(T-g\theta)\in B[T]$ are $D$-invariant, hence lie in $B^D=A$. Since $g=\operatorname{id}$ is one of the factors, $P(\theta)=0$. Reducing modulo $\fm_B$: the image $\bar P\in k(y)[T]$ of $P$ satisfies $\bar P=\prod_{g}(T-\overline{g\theta})$ in $k(x)[T]$, and $\bar P(\bar\theta)=0$; here $\overline{g\theta}=\bar g(\bar\theta)$, where $\bar g$ is the automorphism of $k(x)$ induced by $g$.

Let $\bar m\in k(y)[T]$ be the minimal polynomial of $\bar\theta$. It divides $\bar P$, so it splits in $k(x)[T]$ into linear factors, and each root $\rho$ is of the form $\bar g(\bar\theta)$. Since $\bar m$ is separable and its roots lie in $k(x)$, they lie in $k'$. Hence $k'=k(y)(\bar\theta)$ contains all roots of the separable polynomial $\bar m$: $k'/k(y)$ is Galois.

Let $\tau\in\operatorname{Aut}(k'/k(y))$. Then $\tau(\bar\theta)$ is a root of $\bar m$, so a root of $\bar P=\prod_g(T-\bar g\bar\theta)$; hence $\tau(\bar\theta)=\bar g(\bar\theta)$ for some $g\in D$. Every automorphism of $k(x)/k(y)$ preserves $k'$ (the separable closure), so $\bar g|_{k'}\in\operatorname{Aut}(k'/k(y))$ agrees with $\tau$ on the generator $\bar\theta$; thus $\bar g|_{k'}=\tau$. Consequently $D\to\operatorname{Aut}(k'/k(y))$, $g\mapsto\bar g|_{k'}$, is surjective. The restriction map $\operatorname{Aut}(k(x)/k(y))\to\operatorname{Aut}(k'/k(y))$ is injective: $k(x)/k'$ is purely inseparable, so an automorphism of $k(x)$ that is the identity on $k'$ fixes the elements $u$ with $u^{p^e}\in k'$ (because $p^e$-th power is injective), i.e.\ everything. Hence it is an isomorphism (the composite $D\to\operatorname{Aut}(k(x)/k(y))\to\operatorname{Aut}(k'/k(y))$ being onto), $D\to\operatorname{Aut}(k(x)/k(y))$ is surjective, and its kernel $I$ is the kernel of $D\to\operatorname{Gal}(k'/k(y))$. This gives the exact sequence $1\to I\to D\to\operatorname{Gal}(k'/k(y))\to1$. In particular $I$ is normal in $D$, and $k(x)^I\supseteq k'$.

\emph{A correction to (c) as printed.} Each $\tau\in\operatorname{Gal}(k'/k(y))$ equals $\bar g|_{k'}$ for exactly $|I|$ elements $g\in D$, so $\bar P=\prod_\tau(T-\tau\bar\theta)^{|I|}=\bar m^{|I|}$. Hence $A[T]/(P)$ is etale over $A$ only if $I=1$ (its special fiber is $k(y)[T]/(\bar m^{|I|})$, nonreduced for $|I|>1$). So in (c) one has to take instead $C=A[T]/(Q)$, where $Q\in A[T]$ is a monic lift of the minimal polynomial $\bar m$ of $\bar\theta$ (no map $C\to B$ is needed, only $C\otimes_AB$). This is what we do.

\textbf{(c)} Let $C=A[T]/(Q)$, $B_1=C\otimes_AB=B[T]/(Q)$ and $X_1=\operatorname{Spec}B_1$, which is affine and of finite type over $\operatorname{Spec}A$. By Lemma E(i), $C$ is a finite free local $A$-algebra, and it is etale over $A$ since its special fiber is the separable field extension $k'$ of $k(y)$. The group $D$ acts on $B_1$ through $B$, by $C$-algebra automorphisms. Since $C$ is $A$-flat, Exercise 3.18(a) gives $B_1^D=B^D\otimes_AC=C$, i.e.\ $\operatorname{Spec}C=X_1/D$.

By Lemma E(iii), $B$ is etale over $A$ if and only if $(B_1)_{\fn_\sigma}$ is etale over $C$ for all $\sigma\in\operatorname{Gal}(k'/k(y))$. These are all the local rings of $X_1$ at closed points, and the fiber of $X_1\to\operatorname{Spec}C$ over the closed point $z$ of $\operatorname{Spec}C$ consists of the $\fn_\sigma$ (the fiber is $\operatorname{Spec}(E\otimes_{k(y)}k')$, a spectrum of an Artinian ring, so its points are the maximal ideals $\fn_\sigma$).

\emph{Decomposition groups.} For $g\in D$ and $\sigma\in\operatorname{Gal}(k'/k(y))$, the ideal $g\fn_\sigma$ is the kernel of $\bar g\circ\operatorname{ev}_\sigma\circ g^{-1}:B_1\to k(x)$, which is $b\mapsto\bar b$, $T\mapsto\bar g(\sigma\bar\theta)$. So $g\fn_\sigma=\fn_{\bar g|_{k'}\circ\sigma}$. Since $D\to\operatorname{Gal}(k'/k(y))$ is onto, $D$ permutes the $\fn_\sigma$ transitively, and the stabilizer of every $\fn_\sigma$ is $\{g:\bar g|_{k'}=\operatorname{id}\}=I$. Hence the decomposition group of every point of $X_1$ over $z$ is $I$. (In particular, the $(B_1)_{\fn_\sigma}$ are all isomorphic over $C$, via the elements of $D$, so ``for all $\sigma$'' can be replaced by ``for $\sigma=\operatorname{id}$''.)

\emph{Reduction.} Let $\fn=\fn_{\operatorname{id}}$ and $R=(B_1)_\fn=\mathcal O_{X_1,\fn}$. The group $I$ (the decomposition group of $\fn$) acts on $R$, and it acts trivially on the residue field $k(\fn)=k(x)$, by the very definition of $I$. So for the action of the group $I$ on $X_1$ at $\fn$ the decomposition group and the inertia group both equal $I$. By the exact sequence of (b) applied to this action, $k(x)$ is purely inseparable over the residue field of $R^I$ (Step~3 below). Thus the question for $B$ over $A$ has been reduced, via $B_1$ over $C$, to the situation ``$D=I$ and $k(x)$ purely inseparable over $k(y)$'', as claimed. The proof is completed in (d).

\textbf{(d)} We prove, with $G=D$:

\par\emph{$B$ is etale over $A$ if and only if $I$ acts trivially on $B=\mathcal O_{X,x}$.}\par

\emph{Step 1.} Let $\fn$, $R$ be as in (c). By (c), $B$ is etale over $A$ iff $C\to R$ is etale.

\emph{Step 2.} Apply Exercise 3.18 to $D$ acting on $X_1$ at the point $\fn$, whose decomposition group is $I$. Let $\fn'$ be the image of $\fn$ in $X_1/I$ and $C'=\mathcal O_{X_1/I,\fn'}$. By 3.18(c) and the Remark, $C\to C'$ is flat, $\fm_CC'=\fm_{C'}$ and $k(C')=k(C)=k'$. By Exercise 3.18(b), applied to $I$ acting on $X_1$ (the decomposition group of $\fn$ in $I$ is $I$), $C'=R^I$, $R$ is finite over $C'$, and $\widehat{C'}\to\widehat R$ is injective with image $(\widehat R)^I$. By Lemma L, $C\to R$ is etale iff $C'\to R$ is etale.

\emph{Step 3.} (Here $D=I$.) Apply the sequence of (b) to the action of $I$ on $R$ (that is, on $X_1$ at $\fn$): the inertia group is all of $I$, so $\operatorname{Gal}(k''/k(C'))=1$ where $k''$ is the separable closure of $k(C')$ in $k(\fn)=k(x)$. Hence $k''=k(C')$, i.e.\ $k(x)$ is purely inseparable over $k(C')$.

\emph{Step 4: $C'\to R$ is etale iff $I$ acts trivially on $R$.} If $I$ acts trivially on $R$, then $C'=R^I=R$ and $C'\to R$ is the identity, which is etale. Conversely suppose $C'\to R$ is etale. Then $k(\fn)/k(C')$ is separable (U2) and purely inseparable (Step 3), so $k(\fn)=k(C')$. This is the hypothesis of Prop.~3.26 for the morphism $X_1\to X_1/I$ at $\fn$ (a finite morphism of finite type between locally Noetherian schemes). By Prop.~3.26, $\widehat{C'}\to\widehat R$ is an isomorphism. By Step 2 its image is $(\widehat R)^I$, so $I$ acts trivially on $\widehat R$, hence on $R\subseteq\widehat R$ (Thm.~1.3.16(b); $R\to\widehat R$ is injective, being faithfully flat).

\emph{Step 5: $I$ acts trivially on $R$ iff it acts trivially on $B$.} If $I$ acts trivially on $B$, it acts trivially on $B_1=B[T]/(Q)$ (it acts on coefficients), hence on its localization $R$. Conversely, $B\to B_1$ is injective ($B_1$ is free over $B$) and $B_1\to R$ is a localization, so $B\to R$ is flat; it is local ($\fn\cap B=\fm_B$), hence faithfully flat (Cor.~1.2.20), hence injective (Exercise 1.2.19(a)). It is $I$-equivariant, so if $I$ acts trivially on $R$ it acts trivially on $B$.

Steps 1--5 prove the boxed statement. With (a) we get the theorem for general $G$: \emph{$f:X\to X/G$ is etale at $x$ iff $I$ acts trivially on $\mathcal O_{X,x}$.}

\emph{$X/I\to Y$ is etale at the image of $x$.} First let $G=D$. Replacing $X$ by a $G$-stable affine open neighborhood $\operatorname{Spec}B_X$ of $x$ (see (Q1)) we may assume $X$ affine, so that $X'=X/I=\operatorname{Spec}B_X^I$ is again affine. Let $x'$ be the image of $x$ in $X'$. The group $\bar G=D/I$ acts on $X'$, $X'/\bar G=X/D=Y$, and the fiber of $X'\to Y$ over $y$ is $\{x'\}$, so the decomposition group of $x'$ is all of $\bar G$. By Exercise 3.18(b) applied to $I$ acting on $X$, $\mathcal O_{X',x'}=\mathcal O_{X,x}^I$, and by Step 3 (the same argument, for $I$ acting on $X$ at $x$) $k(x)$ is purely inseparable over $k(x')$. The elements of $k'$ are separable over $k(y)\subseteq k(x')$ and lie in $k(x)$, hence they lie in $k(x')$ (an element which is both separable and purely inseparable over $k(x')$ belongs to $k(x')$): $k'\subseteq k(x')$. The inclusion $k(x')\subseteq k(x)$ is $D$-equivariant, so the composite $D/I\to\operatorname{Aut}(k(x')/k(y))\to\operatorname{Aut}(k'/k(y))$ is the isomorphism $D/I\simeq\operatorname{Gal}(k'/k(y))$ of (b). Hence $D/I\to\operatorname{Aut}(k(x')/k(y))$ is injective: the inertia group of $x'$ for the action of $\bar G$ is trivial. In particular it acts trivially on $\mathcal O_{X',x'}$, so by the theorem just proved (for the group $\bar G$ acting on $X'$, with $G=D$-type decomposition group), $X'\to X'/\bar G=Y$ is etale at $x'$.

For general $G$: $X/I\to X/D$ is etale at the image of $x$ by the case just proved, applied to the group $D$ acting on $X$ (the decomposition group of $x$ in $D$ is $D$, and its inertia group is $I$). Also $X/D\to X/G$ is etale at $z$ by Exercise 3.18(c). A composition of etale morphisms is etale (Prop.~3.22(c)). Hence $X/I\to X/G$ is etale at the image of $x$. \hfill$\square$

\exercise{3.20}{Let $X=\operatorname{Spec}k[T_1,\dots,T_n]/(F_1,\dots,F_r)$ be an affine variety over a field $k$, and $x$ a closed point of $X$. Let us consider $J_x=\bigl((\partial F_i/\partial T_j)(x)\bigr)_{i,j}$ as a matrix with coefficients in $k(x)$. Show that $X$ is smooth at $x$ if and only if $\operatorname{rank}J_x=n-\dim\mathcal O_{X,x}$. Hence the Jacobian criterion is a criterion for smoothness.}
\pf
Let $\bar k$ be an algebraic closure of $k$ and $\pi:X_{\bar k}\to X$ the projection; $X_{\bar k}=V(F_1,\dots,F_r)\subseteq\mathbb A^n_{\bar k}$ is defined by the same polynomials. Since $x$ is closed, $k(x)$ is a finite extension of $k$, so the fiber $\pi^{-1}(x)=\operatorname{Spec}(k(x)\otimes_k\bar k)$ is the spectrum of a finite $\bar k$-algebra. Its points $x'$ are closed in $X_{\bar k}$, with residue field $\bar k$, and correspond to the $k$-embeddings $\sigma:k(x)\to\bar k$ (the maximal ideal of $x'$ is the kernel of $k(x)\otimes\bar k\to\bar k$, $a\otimes b\mapsto\sigma(a)b$). As a rational point of $\mathbb A^n_{\bar k}$, $x'=(\sigma(t_1),\dots,\sigma(t_n))$, where $t_j\in k(x)$ is the image of $T_j$.

Fix such an $x'$. Since the $F_i$ have coefficients in $k$,
\[ \frac{\partial F_i}{\partial T_j}(x')=\sigma\Bigl(\frac{\partial F_i}{\partial T_j}(x)\Bigr),\quad\text{i.e. } J_{x'}=\sigma(J_x),\]
and applying a field homomorphism does not change the rank: $\operatorname{rank}J_{x'}=\operatorname{rank}J_x$.

The morphism $\pi$ is flat (a base change of $\operatorname{Spec}\bar k\to\operatorname{Spec}k$) and the local ring of the fiber $\pi^{-1}(x)$ at $x'$ is Artinian, of dimension $0$. By Thm.~3.12 (the case of equality for flat morphisms),
\[ \dim\mathcal O_{X_{\bar k},x'}=\dim\mathcal O_{X,x}+0=\dim\mathcal O_{X,x}. \]
By the Jacobian criterion (Thm.~2.19) for the rational point $x'$ of $X_{\bar k}\subseteq\mathbb A^n_{\bar k}$, $X_{\bar k}$ is regular at $x'$ if and only if $\operatorname{rank}J_{x'}=n-\dim\mathcal O_{X_{\bar k},x'}$. By what we just proved, this is equivalent to $\operatorname{rank}J_x=n-\dim\mathcal O_{X,x}$ --- a condition which does not depend on $x'$. By Def.~3.28, $X$ is smooth at $x$ if and only if all the points $x'$ of $X_{\bar k}$ above $x$ are regular, i.e.\ if and only if $\operatorname{rank}J_x=n-\dim\mathcal O_{X,x}$. \hfill$\square$

\exercise{3.21}{Let $F\in k[T_0,\dots,T_n]$ be a homogeneous polynomial of degree $d$ over a field $k$. We suppose that $d$ is prime to the characteristic of $k$. Show that the closed subscheme $V_+(F)$ of $\mathbb P^n_k$ is smooth over $k$ if and only if the $\partial F/\partial T_i$ are without common zero in $\bar k$.}
\pf
\emph{Hypotheses.} We assume $F\neq0$ and $d\ge1$, $n\ge1$ (for $F=0$ the scheme is $\mathbb P^n_k$, which is smooth although all partial derivatives vanish; for $d=0$, $F$ is a nonzero constant and $V_+(F)=\emptyset$). ``Common zero in $\bar k$'' means a common zero in $\mathbb P^n(\bar k)$, i.e.\ a nonzero vector of $\bar k^{n+1}$ (the partial derivatives are homogeneous, so this is a projective notion).

\textbf{Lemma A} (valid for any $d\ge1$ and any characteristic). Let $F\neq0$ be homogeneous of degree $d\ge1$ over $k$, $X=V_+(F)\subseteq\mathbb P^n_k$, and let $x\in X(\bar k)$. Then $X_{\bar k}$ is regular at $x$ if and only if the numbers $\frac{\partial F}{\partial T_j}(x)$, $0\le j\le n$, are not all zero.

\emph{Proof.} Choose $i$ with $x_i\neq0$ and normalize $x_i=1$. On the chart $D_+(T_i)\simeq\mathbb A^n_{\bar k}$ with coordinates $t_j=T_j/T_i$ ($j\ne i$), $X_{\bar k}$ is $V(f)$ with $f=F(t_0,\dots,1,\dots,t_n)\neq0$, and $x$ is a rational point, with $f(x)=0$. The local ring $\mathcal O=\mathcal O_{\mathbb A^n,x}$ is regular of dimension $n$ and $f\neq0$ is a nonunit in the domain $\mathcal O$, so $\dim\mathcal O/f\mathcal O=n-1$ (Thm.~2.5.15, as in the proof of Thm.~3.12). By the Jacobian criterion (Thm.~2.19), $X_{\bar k}$ is regular at $x$ if and only if the $1\times n$ matrix $(\partial f/\partial t_j(x))_{j\neq i}$ has rank $n-(n-1)=1$, i.e.\ if and only if it is nonzero. Now $\partial f/\partial t_j(x)=\frac{\partial F}{\partial T_j}(x)$ for $j\neq i$. Euler's identity $\sum_jT_j\frac{\partial F}{\partial T_j}=dF$ at $x$ (where $F(x)=0$, $x_i=1$) gives $\frac{\partial F}{\partial T_i}(x)=-\sum_{j\neq i}x_j\frac{\partial F}{\partial T_j}(x)$. So all $\frac{\partial F}{\partial T_j}(x)$ vanish if and only if those with $j\ne i$ do. $\square$

\emph{Proof of the exercise.} By Cor.~2.17, $X_{\bar k}$ is regular if and only if it is regular at its closed points, which are the points $x\in X(\bar k)$ ($\bar k$ is algebraically closed). By Lemma A, $X$ is smooth over $k$ if and only if there is no $x\in X(\bar k)$ at which all $\partial F/\partial T_j$ vanish. Now let $x\in\mathbb P^n(\bar k)$ be a common zero of all the partial derivatives. Euler's identity gives $d\,F(x)=\sum_jx_j\frac{\partial F}{\partial T_j}(x)=0$, and $d$ is invertible in $k$, so $F(x)=0$, i.e.\ $x\in X(\bar k)$. Hence ``there is no point of $X(\bar k)$ where all partials vanish'' is equivalent to ``the partials have no common zero in $\mathbb P^n(\bar k)$''. \hfill$\square$

\exercise{3.22}{Let $C$ be a conic over a field $k$, that is $C=V_+(F)\subset\mathbb P^2_k$, where $F(x,y,z)$ is a homogeneous polynomial of degree $2$.
(a) Show that $C$ is reduced (resp.\ irreducible) if and only if $F$ is not a square modulo $k^*$ (resp.\ $F$ is irreducible).
(b) Let us suppose that $C$ contains a rational point $P$.
(1) Show that if $C$ is singular at $P$, then (up to isomorphism) $C$ is of the form $V_+(G)$, with $G(x,y)\in k[x,y]$ homogeneous of degree $2$.
(2) Let us suppose $C$ is regular at $P$. Let $D$ be a line in $\mathbb P^2_k$ passing through $P$ and different from $PT_{C,P}$ (see the notation in Exercise 2.9). Show that $D$ meets $C$ at another rational point $Q$. Show that $C\simeq V_+(ax^2+yz)$ with $a\in k$ (let $R$ be the intersection point of the lines $PT_{C,P}$ and $PT_{C,Q}$, and reduce to $P=(1,0,0)$, $Q=(0,1,0)$ and $R=(0,0,1)$). Deduce from this that either $C$ is the union of two distinct lines, or $C\simeq\mathbb P^1_k$.
(c) Let us suppose $\operatorname{char}(k)\neq2$. Show that $C\simeq V_+(ax^2+by^2+cz^2)$, and that $C$ is smooth if and only if $abc\neq0$.
(d) Let us suppose $\operatorname{char}(k)=2$, $C(k)=\emptyset$, and that $C$ is not smooth.
(1) Show that there exist $a,b,c\in k$ such that $C\simeq V_+(ax^2+by^2+cz^2)$. Show that $C$ does not contain any smooth point.
(2) Show that $\mathcal O_C(D_+(z))$ is equal to the integral closure of $k[x/z]$ in $K(C)$. Deduce from this that $C$ is normal, and therefore regular.}
\pf
Throughout, $F\neq0$. ``Regular at $P$'' is meant for the scheme $C$ over $k$ (the ring $\mathcal O_{C,P}$ is regular). Recall (Exercise 2.9) that for $P\in C(k)$ the ``tangent line'' is $PT_{C,P}=\{t:\sum_i\frac{\partial F}{\partial T_i}(P)t_i=0\}$.

\textbf{(a)} Let $S=k[x,y,z]$. \emph{Possible shapes of $F$.} Since $\deg F=2$, either $F$ is irreducible over $k$, or $F=L_1L_2$ with $L_1,L_2$ linear forms over $k$. In the second case either $L_1,L_2$ are not proportional, or $F=\lambda L^2$ with $\lambda\in k^*$ and $L$ a linear form (this is the case ``$F$ is a square modulo $k^*$'').

\emph{Case $F$ irreducible.} $S/(F)$ is a graded domain, so $C=\operatorname{Proj}S/(F)$ is integral (in particular reduced and irreducible).

\emph{Case $F=L_1L_2$, $L_1,L_2$ not proportional.} Then $F$ is squarefree, $S/(F)$ is reduced, and so are the rings of the charts $D_+(T_i)\cap C$ (degree-$0$ parts of the localizations $(S/F)_{T_i}$). So $C$ is reduced. Set-theoretically $C=V_+(L_1)\cup V_+(L_2)$, a union of two distinct lines, neither of which contains the other: $C$ is reducible.

\emph{Case $F=\lambda L^2$.} After a linear change of coordinates we may take $L=z$. In the chart $D_+(x)$ the coordinate ring of $C$ is $k[Y,Z]/(\lambda Z^2)$ ($Y=y/x$, $Z=z/x$), in which $Z\neq0$ but $Z^2=0$. So $C$ is not reduced; but $C_{\operatorname{red}}=V_+(z)\simeq\mathbb P^1_k$ is irreducible, so $C$ is irreducible.

\emph{Conclusion.} $C$ is reduced if and only if $F$ is not a square modulo $k^*$ (i.e.\ not of the form $\lambda L^2$). Moreover $C$ is integral if and only if $F$ is irreducible. \emph{Note on the printed statement:} the equivalence ``$C$ irreducible $\iff$ $F$ irreducible'' is false as stated for $F=\lambda L^2$, where $C$ is irreducible (but not reduced) although $F$ is not irreducible. The correct statements are: $C$ is irreducible if and only if $F$ is irreducible or of the form $\lambda L^2$; $C$ is integral if and only if $F$ is irreducible. This is what is proved above.

\textbf{(b)} Choose coordinates such that $P=(1:0:0)$. Since $F(P)=0$, $F$ has no $x^2$ term:
\[ F=x\,L(y,z)+q(y,z),\qquad L=L_yy+L_zz,\ q\ \text{a binary quadratic form}. \]
In the chart $x=1$ the point $P$ is the origin of $\mathbb A^2$ and $f=L+q$. Let $\fm=(y,z)$ in the regular local ring $\mathcal O=k[y,z]_\fm$. The element $f=L+q\in\fm$ is nonzero, and $f\notin\fm^2$ if and only if $L\neq0$. By Cor.~2.12, $\mathcal O/f\mathcal O=\mathcal O_{C,P}$ is regular if and only if $L\neq0$.

(1) If $C$ is singular at $P$ then $L=0$ and $F=q(y,z)$ is a binary quadratic form: $C=V_+(G)$ with $G(y,z)=q$ (renaming the variables).

(2) Now $L\neq0$. The gradient of $F$ at $P$ is $(0,L_y,L_z)$, so $T:=PT_{C,P}=\{L_yy+L_zz=0\}$ is a line through $P$. A line $D$ through $P$ has the form $\{\lambda y+\mu z=0\}$, with $(\lambda,\mu)\neq0$, and its points are $(s:\mu t:-\lambda t)$, $(s:t)\in\mathbb P^1(k)$. One has $D=T$ iff $(\lambda,\mu)$ is proportional to $(L_y,L_z)$, iff $\ell:=L(\mu,-\lambda)=L_y\mu-L_z\lambda=0$. We have
\[ F(s,\mu t,-\lambda t)=s\,t\,L(\mu,-\lambda)+t^2q(\mu,-\lambda)=t\,(\ell s+q_0t),\qquad q_0=q(\mu,-\lambda). \]
Suppose $D\neq T$, i.e.\ $\ell\neq0$. The zeros of this form on $D$ are $t=0$ (the point $P$) and $(s:t)=(-q_0:\ell)$, i.e.\ the rational point $Q=(-q_0:\mu\ell:-\lambda\ell)$, which differs from $P$ since $\ell\neq0$. So $D$ meets $C$ in $P$ and in another rational point $Q$.

\emph{Normal form.} Choose new coordinates with $P=(1:0:0)$, $Q=(0:1:0)$ (complete the vectors representing $P,Q$ to a basis of $k^3$). Since $F(P)=F(Q)=0$,
\[ F=uxy+vxz+wyz+ez^2. \]
The line $D=PQ$ is $\{z=0\}$. The gradient at $P$ is $(0,u,v)$, so $T=\{uy+vz=0\}$, and $D\neq T$ means $u\neq0$. (In particular $(u,v)\neq0$ and $(u,w)\neq0$, so $P$ and $Q$ are regular points.) At $Q$ the gradient is $(u,0,w)$, so $PT_{C,Q}=\{ux+wz=0\}$. The lines $PT_{C,P}$ and $PT_{C,Q}$ meet at $R=(-w:-v:u)$, whose third coordinate is $u\neq0$. Hence $P,Q,R$ are linearly independent and we may take them as the new coordinate points: put
\[ x=x'-wz',\qquad y=y'-vz',\qquad z=uz'. \]
A direct computation gives
\begin{align*}
F&=u(x'-wz')(y'-vz')+v(x'-wz')\,uz'+w(y'-vz')\,uz'+e\,u^2z'^2\\
&=u\,x'y'+u(eu-vw)\,z'^2 .
\end{align*}
Thus $C=V_+(x'y'+e'z'^2)$ with $e'=eu-vw\in k$ ($u$ is a unit). Renaming $(z',x',y')\to(x,y,z)$ and $a=e'$: $C\simeq V_+(ax^2+yz)$.

\emph{Conclusion.} If $a=0$, $C=V_+(yz)$ is the union of the two distinct lines $y=0$ and $z=0$. If $a\neq0$, $C\cap V_+(y,z)=\operatorname{Proj}k[x]/(ax^2)=\emptyset$, so $C$ is covered by the opens $C\cap D_+(y)$ and $C\cap D_+(z)$. Dividing the equation by $y^2$, resp.\ $z^2$:
\[ C\cap D_+(y)=\operatorname{Spec}k[X,Z]/(aX^2+Z)\simeq\mathbb A^1_k\ (\text{coordinate }X=x/y),\qquad C\cap D_+(z)\simeq\mathbb A^1_k\ (\text{coordinate }X'=x/z). \]
On the overlap $\{y\neq0,z\neq0\}$, where $X\neq0$, we have $z/y=-aX^2$ and so $X'=x/z=X/(z/y)=-1/(aX)$. So $C$ is obtained by gluing two copies of $\mathbb A^1_k$ along $\mathbb G_{m,k}$ via $X'=-1/(aX)$; this is $\mathbb P^1_k$ (replace the coordinate $X'$ by $-aX'$, which makes the gluing $X'\cdot X=1$). Hence $C\simeq\mathbb P^1_k$.

\textbf{(c)} Let $\operatorname{char}k\neq2$. Let $B(u,v)=\frac12\bigl(F(u+v)-F(u)-F(v)\bigr)$, a symmetric bilinear form on $k^3$ with $F(v)=B(v,v)$. We show that $k^3$ has a $B$-orthogonal basis; then in the corresponding coordinates $F=ax^2+by^2+cz^2$ ($a,b,c\in k$). If $B(v,v)=0$ for all $v$ then $B=0$ by polarization ($B(u,v)=\frac12(B(u+v,u+v)-B(u,u)-B(v,v))$) and every basis is orthogonal. Otherwise take $v_1$ with $B(v_1,v_1)\neq0$; then $k^3=kv_1\oplus v_1^\perp$ (since $v-\frac{B(v,v_1)}{B(v_1,v_1)}v_1\in v_1^\perp$), and we apply the same argument to the restriction of $B$ to the plane $v_1^\perp$, and then to a line. So $C\simeq V_+(ax^2+by^2+cz^2)$, with $(a,b,c)\neq0$ since $F\neq0$.

Smoothness: the degree $2$ is prime to $\operatorname{char}k$. By Exercise 3.21, $C$ is smooth if and only if $2ax,2by,2cz$ have no common zero in $\mathbb P^2(\bar k)$. If $abc\neq0$ this holds. If, say, $c=0$, then $(0:0:1)$ is a common zero, and similarly for $a=0$ or $b=0$. So $C$ is smooth if and only if $abc\neq0$.

\textbf{(d)} Let $\operatorname{char}k=2$. Write $F=ax^2+by^2+cz^2+uxy+vxz+wyz$. Then
\[ \frac{\partial F}{\partial x}=uy+vz,\qquad\frac{\partial F}{\partial y}=ux+wz,\qquad\frac{\partial F}{\partial z}=vx+wy, \]
and by Lemma A (Exercise 3.21; it holds for $d=2$ in characteristic $2$) a point $x\in C(\bar k)$ is a singular point of $C_{\bar k}$ if and only if all three partials vanish at $x$.

(1) \emph{Diagonal form.} Suppose $(u,v,w)\neq(0,0,0)$. The common zeros in $\mathbb P^2(\bar k)$ of the three partials are given by the kernel of the matrix $N=\begin{pmatrix}0&u&v\\u&0&w\\v&w&0\end{pmatrix}$. We have $N\,(w,v,u)^t=(2uv,2uw,2vw)^t=0$, and $(w,v,u)\neq0$, so $\operatorname{rank}N\le2$. On the other hand $\operatorname{rank}N\ge2$: if $u\neq0$ the first two rows are independent; if $u=0$ the rows are $(0,0,v),(0,0,w),(v,w,0)$ with $(v,w)\neq0$, and the third row is not a multiple of $(0,0,1)$. So there is exactly one common zero, namely $p=(w:v:u)$, a point of $\mathbb P^2(k)$. As $C$ is not smooth, $C_{\bar k}$ has a singular point (Cor.~2.17), which must be a common zero of the partials, i.e.\ $p$, so $p\in C(\bar k)$ and $F(p)=0$ with $p\in\mathbb P^2(k)$: $C(k)\neq\emptyset$, a contradiction. Hence $u=v=w=0$ and
\[ F=ax^2+by^2+cz^2. \]
Then all three partials vanish identically, so every point of $C_{\bar k}$ is singular (Lemma A): $C$ has no smooth point.

(2) \emph{The coefficients.} $C(k)=\emptyset$ gives $F(1,0,0)=a\neq0$, $b\neq0$, $c\neq0$. Fix an algebraic closure $\bar k$ and square roots $\sqrt a,\sqrt b,\sqrt c\in\bar k$. We claim that $\sqrt a,\sqrt b,\sqrt c$ are linearly independent over $k$. Indeed, if $\lambda\sqrt a+\mu\sqrt b+\nu\sqrt c=0$ with $\lambda,\mu,\nu\in k$ not all zero, squaring (characteristic $2$) gives $a\lambda^2+b\mu^2+c\nu^2=0$, i.e.\ $(\lambda:\mu:\nu)\in C(k)$. Put $\rho=\sqrt a/\sqrt b$ and $\delta=\sqrt c/\sqrt b$ (so $\rho^2=a/b$, $\delta^2=c/b$). Then $1,\rho,\delta$ are linearly independent over $k$ (divide the relation by $\sqrt b$). Let $k_0=k(\rho,\delta)=k[\rho,\delta]$, a finite (purely inseparable) extension of $k$.

Let $X=x/z$, $Y=y/z$. Then
\[ A:=\mathcal O_C(D_+(z))=k[X,Y]/(aX^2+bY^2+c)=k[X,Y]/(Y^2+\alpha X^2+\gamma),\qquad \alpha=\rho^2,\ \gamma=\delta^2. \]
$A$ is a free $k[X]$-module with basis $1,Y$. Consider the $k$-algebra homomorphism
\[ \iota:A\to k_0[X],\qquad X\mapsto X,\ Y\mapsto\theta:=\rho X+\delta; \]
it is well defined since $\theta^2=\alpha X^2+\gamma$. Choose a $k$-basis $e_1=1,e_2=\rho,e_3=\delta,e_4,\dots$ of $k_0$. Then $k_0[X]=\bigoplus_ie_ik[X]$, and the $e_i$ are also linearly independent over $k(X)$ in $k_0(X)$ (indeed $k_0\otimes_kk(X)$ is a localization of $k_0[X]$, hence a subring of $k_0(X)$). For $r,s\in k(X)$,
\[ r+s\theta=r\cdot e_1+sX\cdot e_2+s\cdot e_3\qquad(*) \]
has coordinates $(r,sX,s,0,\dots)$. In particular $\iota(r+sY)=0$ forces $r=s=0$ for $r,s\in k[X]$: $\iota$ is injective. So $A$ is a domain, identified with $k[X]+k[X]\theta\subseteq k_0[X]$, and $K(C)=\operatorname{Frac}(A)=k(X)+k(X)\theta\subseteq k_0(X)$.

\emph{$A$ is the integral closure of $k[X]$ in $K(C)$.} $A$ is finite over $k[X]$, so it is contained in this integral closure $A'$. Conversely let $u=r+s\theta\in K(C)$ ($r,s\in k(X)$) be integral over $k[X]$. Then $u$ is integral over the normal domain $k_0[X]$ (a PID), whose fraction field contains $u$, so $u\in k_0[X]$. By $(*)$, the coordinates of $u$ with respect to the basis $(e_i)$ over $k(X)$ are $r$ and $s$ (on $e_1$ and $e_3$), and they lie in $k[X]$ since $u\in k_0[X]=\bigoplus e_ik[X]$. Hence $u=r+s\theta\in A$. So $A=A'$.

\emph{Normality and regularity of $C$.} $F$ is irreducible (a reducible $F$ would be a product of two linear forms, and a line has a rational point), so $C$ is integral, and the above shows that $\mathcal O_C(D_+(z))$ is a normal domain. The hypotheses ($C(k)=\emptyset$, $\operatorname{char}k=2$, $F$ diagonal) are symmetric in $x,y,z$ (permuting the variables permutes $a,b,c$), so the same holds for $D_+(x)$ and $D_+(y)$, which together with $D_+(z)$ cover $C$. Hence $C$ is a normal scheme. A normal Noetherian local ring of dimension $1$ is a discrete valuation ring (Prop.~1.12), and the generic point is regular (a field). So every local ring of $C$ is regular: $C$ is regular. (By (1), $C$ is nowhere smooth: it is a regular, non-smooth conic, cf.\ Rem.~3.34.) \hfill$\square$

\exercise{3.23}{Let $k$ be a field and let $F(x,y,z)\in k[x,y,z]$ be a homogeneous polynomial of degree $2$ defining a smooth conic $V_+(F)$ over $k$.
(a) If $\operatorname{char}(k)\neq2$, show that after a change of variables, one can write $F=ax^2+by^2+cz^2$ for some $a,b,c\in k^*$. Let us suppose moreover that $k$ is finite with $q$ elements. Show that the sets $\{ax^2\mid x\in k\}$ and $\{-c-by^2\mid y\in k\}$ both have $(q+1)/2$ elements. Deduce from this that there exists a rational point in $V_+(F)$.
(b) Let us suppose that $\operatorname{char}(k)=2$ and that $k$ is perfect. Show that after a suitable change of variables, one has $F=xy+z^2$. Deduce from this that $V_+(F)$ contains a rational point.}
\pf
(Here $F\neq0$, since the zero polynomial does not define a conic; a smooth conic is also nonempty over $\bar k$.)

\textbf{(a)} By Exercise 3.22(c), after a linear change of variables $F=ax^2+by^2+cz^2$, and since the conic is smooth, $abc\neq0$ (Exercise 3.22(c)).

Let now $k=\mathbb F_q$, $q$ odd. The squaring map $k^*\to k^*$ is a $2$-to-$1$ homomorphism (its kernel is $\{\pm1\}$, of order $2$ as $q$ is odd), so $k^*$ contains $(q-1)/2$ nonzero squares, and $k$ contains $(q-1)/2+1=(q+1)/2$ squares (including $0$). The map $x\mapsto ax^2$ has as image $a\cdot\{\text{squares}\}$, which has $(q+1)/2$ elements since $a\neq0$. Likewise $y\mapsto-c-by^2$ has image $-c-b\cdot\{\text{squares}\}$, the image of the set of squares under the bijection $t\mapsto-c-bt$ of $k$, hence of cardinality $(q+1)/2$. The two sets are subsets of $k$, which has $q$ elements, and $\frac{q+1}2+\frac{q+1}2=q+1>q$. So they intersect: there are $x_0,y_0\in k$ with $ax_0^2=-c-by_0^2$, i.e.\ $F(x_0,y_0,1)=0$. The point $(x_0:y_0:1)$ is a rational point of $V_+(F)$ (for the new coordinates; hence also for the original ones, by the linear change of variables).

\textbf{(b)} Let $\operatorname{char}k=2$, $k$ perfect, $F=ax^2+by^2+cz^2+uxy+vxz+wyz$. By the computation of Exercise 3.22(d), the singular points of $V_+(F)_{\bar k}$ are the points of $V_+(F)(\bar k)$ at which $uy+vz$, $ux+wz$, $vx+wy$ vanish (Lemma A of Exercise 3.21).

\emph{Step 1: $(u,v,w)\neq0$.} If $u=v=w=0$, all three partials vanish identically and every point of the nonempty $V_+(F)_{\bar k}$ is singular, contradicting smoothness.

\emph{Step 2: the point $p$.} By Step 1 and the computation in Exercise 3.22(d), the partials have exactly one common zero $p=(w:v:u)\in\mathbb P^2(k)$ in $\mathbb P^2(\bar k)$. Since $V_+(F)$ is smooth, $p\notin V_+(F)$, i.e.\ $F(p)\neq0$.

\emph{Step 3: normal form.} Choose a linear change of variables such that $p=(0:0:1)$. The partial derivatives of a form transform contragrediently under linear changes of variables, so the property ``all partial derivatives vanish at $p$'' is invariant under linear changes of variables; in the new coordinates $F$ is still smooth, $F(p)\neq0$, and all three partials vanish at $p$. Write the new form again as $F=a'x^2+b'y^2+c'z^2+u'xy+v'xz+w'yz$. Then $\frac{\partial F}{\partial x}(p)=v'$ and $\frac{\partial F}{\partial y}(p)=w'$ vanish, so $v'=w'=0$ and $c'=F(p)\neq0$. Also $u'\neq0$, because otherwise the new form would be diagonal, and by Step 1 (applied in the new coordinates) it would not be smooth. Thus
\[ F=a'x^2+u'xy+b'y^2+c'z^2,\qquad u'c'\neq0. \]
Since $k$ is perfect, every element of $k$ is a square. Replace $z$ by $\sqrt{c'}\,z$ to get $c'=1$. Then
\[ z^2+a'x^2+b'y^2=(z+\sqrt{a'}\,x+\sqrt{b'}\,y)^2\quad(\text{characteristic }2), \]
and the substitution $z_1=z+\sqrt{a'}x+\sqrt{b'}y$ is invertible. In the variables $x,y,z_1$ we get $F=u'xy+z_1^2$, and putting $x_1=u'x$, $F=x_1y+z_1^2$. This is the form $xy+z^2$ in the new variables.

\emph{A rational point.} In these coordinates $(1:0:0)$ satisfies $xy+z^2=0$. Hence $V_+(F)$ has a rational point (the corresponding point of the original coordinates). \hfill$\square$

\exercise{3.24}{Let $X$ be an algebraic variety over a field $k$. Show that the following properties are equivalent:
(i) $X$ is smooth;
(ii) for every finite extension $l$ of $k$, $X_l$ is regular;
(iii) for every purely inseparable extension $l$ of $k$, $X_l$ is regular;
(iv) for every (not necessarily algebraic) extension $l$ of $k$, $X_l$ is regular.}
\pf
We first prove a lemma on regular varieties over algebraically closed fields.

\textbf{Lemma B.} Let $k_1\subseteq k_2$ be algebraically closed fields and $Z$ a regular algebraic variety over $k_1$. Then $Z_{k_2}$ is regular. (This is a special case of the stability of smoothness under base change, Prop.~3.38, whose proof relies on exactly this fact.)

\emph{Proof.} Regularity is local, so we may assume $Z$ affine, $Z=V(I)\subseteq\mathbb A^n_{k_1}$. As $Z$ is regular, its local rings are domains (Prop.~2.11), so its irreducible components are pairwise disjoint; hence $Z=Z_1\sqcup\dots\sqcup Z_s$ with $Z_j$ open and closed in $Z$, integral, of dimension $d_j$. Each $Z_j$ is affine, $Z_j=V(I_j)\subseteq\mathbb A^n_{k_1}$ with $I_j$ prime, generated by $G_{j1},\dots,G_{jm}$; and $Z_{k_2}=\bigsqcup(Z_j)_{k_2}$, $(Z_j)_{k_2}=V(I_jk_2[T])\subseteq\mathbb A^n_{k_2}$. Fix $j$ and write $d=d_j$, $G_i=G_{ji}$, $M=(\partial G_i/\partial T_l)$, an $m\times n$ matrix over $k_1[T]$.

For a closed point $w\in Z_j$ (a $k_1$-rational point) we have $\dim\mathcal O_{Z_j,w}=d$ (Prop.~2.5.23(a)). By Prop.~2.5, $\dim T_{Z_j,w}=n-\operatorname{rank}M(w)$, and $\dim T_{Z_j,w}\ge\dim\mathcal O_{Z_j,w}=d$ (Prop.~2.2(a)), so $\operatorname{rank}M(w)\le n-d$ with equality if and only if $Z_j$ is regular at $w$ (Thm.~2.19). Let $\delta$ be the ideal of $k_1[T]/I_j$ generated by the minors of $M$ of order $n-d$. Since $Z_j$ is regular at all its closed points, $\delta$ is contained in no maximal ideal, so $\delta=(1)$: there are $c_\alpha\in k_1[T]$ and minors $\Delta_\alpha$ of order $n-d$ with $1=\sum_\alpha c_\alpha\Delta_\alpha$ in $k_1[T]/I_j$. This identity holds a fortiori in $k_2[T]/I_jk_2[T]$. Hence for every $k_2$-point $w'$ of $(Z_j)_{k_2}$ some minor of order $n-d$ of $M(w')$ is nonzero, i.e.\ $\operatorname{rank}M(w')\ge n-d$, so $\dim T_{(Z_j)_{k_2},w'}=n-\operatorname{rank}M(w')\le d$.

On the other hand $Z_j$ is integral over the algebraically closed field $k_1$, hence geometrically integral (cf.\ Exercise 3.2.14), so $(Z_j)_{k_2}$ is integral; its function field is generated over $k_2$ by a transcendence basis of $K(Z_j)/k_1$ and an algebraic extension of it, so $\dim(Z_j)_{k_2}=\operatorname{trdeg}_{k_2}=d$. Hence $\dim\mathcal O_{(Z_j)_{k_2},w'}=d$ for every closed point $w'$ (Prop.~2.5.23(a)). We get $d=\dim\mathcal O\le\dim T_{w'}\le d$, so $(Z_j)_{k_2}$ is regular at $w'$. The closed points of $(Z_j)_{k_2}$ are the $k_2$-points ($k_2$ is algebraically closed), so $(Z_j)_{k_2}$ is regular by Cor.~2.17. $\square$

\textbf{(i)$\Rightarrow$(iv).} Let $X$ be smooth, so $X_{\bar k}$ is regular, and let $l\supseteq k$ be any extension. Choose an algebraic closure $\bar l$ of $l$ and a $k$-embedding $\bar k\hookrightarrow\bar l$. Then $X_{\bar l}=(X_{\bar k})_{\bar l}$ is regular by Lemma B. Thus $X_l$ is smooth over $l$ (Def.~3.28), hence regular by Cor.~3.32.

\textbf{(iv)$\Rightarrow$(ii) and (iv)$\Rightarrow$(iii)} are trivial (finite and purely inseparable extensions are extensions).

\textbf{(iii)$\Rightarrow$(i).} Let $l\subseteq\bar k$ be the perfect closure of $k$ (the union of the $k^{1/p^e}$ if $\operatorname{char}k=p>0$, and $l=k$ if $\operatorname{char}k=0$). It is a purely inseparable extension of $k$ and a perfect field. By (iii), $X_l$ is a regular algebraic variety over the perfect field $l$, hence smooth over $l$ (Cor.~3.33). Since $\bar k$ is an algebraic closure of $l$, $X_{\bar k}=(X_l)_{\bar k}$ is regular, i.e.\ $X$ is smooth.

\textbf{(ii)$\Rightarrow$(i).} Assume (ii). Since smoothness means regularity of $X_{\bar k}$, which is local on $X$, we may assume $X=V(F_1,\dots,F_r)\subseteq\mathbb A^n_k$ affine. By Cor.~2.17 it suffices to show that $X_{\bar k}$ is regular at its closed points $x'=(t_1,\dots,t_n)\in\bar k^n$. Let $l=k(t_1,\dots,t_n)$, a finite extension of $k$, and $w$ the $l$-rational point $(t_1,\dots,t_n)$ of $X_l$, so $x'$ lies over $w$. By (ii), $X_l$ is regular at $w$, so by the Jacobian criterion (Thm.~2.19) $\operatorname{rank}J_w=n-\dim\mathcal O_{X_l,w}$. The morphism $X_{\bar k}\to X_l$ is flat and its fiber over $w$ is $\operatorname{Spec}(l\otimes_l\bar k)=\operatorname{Spec}\bar k$, of dimension $0$, so $\dim\mathcal O_{X_{\bar k},x'}=\dim\mathcal O_{X_l,w}$ by Thm.~3.12. The Jacobian matrices of $X_{\bar k}$ at $x'$ and of $X_l$ at $w$ are the same matrix (entries in $l\subseteq\bar k$). So $\operatorname{rank}J_{x'}=n-\dim\mathcal O_{X_{\bar k},x'}$, and $X_{\bar k}$ is regular at $x'$ by Thm.~2.19. \hfill$\square$

\exercise{3.25}{Let $f:X\to Y$ be a morphism of schemes, let $x\in X$, and $y=f(x)$.
(a) Show that we have an exact sequence of $k(x)$-vector spaces: $0\to T_{X_y,x}\to T_{X,x}\to T_{Y,y}\otimes_{k(y)}k(x)$.
(b) Let us suppose that $f$ admits a section $g:Y\to X$, and that $x=g(y)$. Show that $k(x)=k(y)$, and that we have a decomposition $T_{X,x}\simeq T_{X_y,x}\oplus T_{Y,y}$.
(c) Let us suppose that $X,Y$ are locally Noetherian and regular, and that $f$ is of finite type. Show that under the hypothesis of (b), $X_y$ is smooth at $x$.
(d) Let us keep the hypothesis above. Let $d=\dim_xX_y$. Show that we have an isomorphism $\widehat{\mathcal O}_{X,x}\simeq\widehat{\mathcal O}_{Y,y}[[T_1,\dots,T_d]]$, and that $f$ is smooth at $x$.}
\pf
Put $A=\mathcal O_{Y,y}$, $B=\mathcal O_{X,x}$, with maximal ideals $\fm_y,\fm_x$ and residue fields $k(y),k(x)$, and let $\fn=\fm_x/\fm_yB$ be the maximal ideal of $\mathcal O_{X_y,x}=B/\fm_yB$.

\textbf{(a)} We have $\fn=\fm_x/\fm_yB$ and $\fn^2=(\fm_x^2+\fm_yB)/\fm_yB$, so $\fn/\fn^2=\fm_x/(\fm_x^2+\fm_yB)$. Hence the sequence of $k(x)$-vector spaces
\[ \fm_y/\fm_y^2\otimes_{k(y)}k(x)\ \xrightarrow{\ \alpha\ }\ \fm_x/\fm_x^2\ \longrightarrow\ \fn/\fn^2\ \longrightarrow\ 0 \]
(where $\alpha(a\otimes\lambda)=\lambda f^\#(a)$) is exact: the image of $\alpha$ is $(\fm_yB+\fm_x^2)/\fm_x^2$ and the cokernel is $\fm_x/(\fm_x^2+\fm_yB)=\fn/\fn^2$. The functor $\operatorname{Hom}_{k(x)}(-,k(x))$ is exact, so the dual sequence
\[ 0\to T_{X_y,x}\to T_{X,x}\xrightarrow{\ \alpha^\vee\ }\operatorname{Hom}_{k(x)}\bigl(\fm_y/\fm_y^2\otimes_{k(y)}k(x),k(x)\bigr)=\operatorname{Hom}_{k(y)}(\fm_y/\fm_y^2,k(x)) \]
is exact. The right-hand side contains $T_{Y,y}\otimes_{k(y)}k(x)$ (equality if $\dim_{k(y)}\fm_y/\fm_y^2<\infty$, e.g.\ if $Y$ is locally Noetherian), and $\alpha^\vee$ is the tangent map $T_{f,x}$ of Def.~2.1. (In general the target of $T_{f,x}$ in Def.~2.1 is to be read as $\operatorname{Hom}_{k(y)}(\fm_y/\fm_y^2,k(x))$; the exercise is stated with $T_{Y,y}\otimes k(x)$, which is the same space in the Noetherian situation of (b)--(d).) This is the required exact sequence.

\textbf{(b)} Let $g^\#:k(x)=k(g(y))\to k(y)$ and $f^\#:k(y)\to k(x)$ be the maps induced by $g$ and $f$ on residue fields. Since $f\circ g=\operatorname{id}_Y$, we have $g^\#\circ f^\#=\operatorname{id}_{k(y)}$. The homomorphism $g^\#$ of fields is injective, and surjective by this relation, so it is an isomorphism; hence $f^\#=(g^\#)^{-1}$ and $k(x)=k(y)$ (via $f^\#$). We identify the two fields; then $T_{Y,y}\otimes_{k(y)}k(x)=T_{Y,y}$.

By (a), $0\to T_{X_y,x}\to T_{X,x}\xrightarrow{T_{f,x}}T_{Y,y}$ is exact. Applying Prop.~2.2(b) to $f\circ g=\operatorname{id}_Y$ gives $T_{f,x}\circ T_{g,y}=\operatorname{id}_{T_{Y,y}}$ (the identification $k(x)=k(y)$ being used for the base change of scalars). So $T_{f,x}$ is surjective and has the section $T_{g,y}$, and the exact sequence $0\to T_{X_y,x}\to T_{X,x}\to T_{Y,y}\to0$ splits: $T_{X,x}\simeq T_{X_y,x}\oplus T_{Y,y}$.

\textbf{(c)} Now $X,Y$ are locally Noetherian and regular, $f$ is of finite type, so $X_y$ is an algebraic variety over $k(y)$. Since $k(x)=k(y)$, the point $x$ is a rational point of $X_y$, hence closed, and $k(x)$ is separable over $k(y)$. By Prop.~3.30, $X_y$ is smooth at $x$ if and only if it is regular at $x$. By (b), $\dim T_{X_y,x}=\dim T_{X,x}-\dim T_{Y,y}$, and since $X,Y$ are regular at $x,y$ this equals $\dim B-\dim A$. By Thm.~3.12, $\dim\mathcal O_{X_y,x}\ge\dim B-\dim A$, and by Prop.~2.2(a), $\dim\mathcal O_{X_y,x}\le\dim T_{X_y,x}$. Hence
\[ \dim B-\dim A\le\dim\mathcal O_{X_y,x}\le\dim T_{X_y,x}=\dim B-\dim A, \]
so $\dim\mathcal O_{X_y,x}=\dim T_{X_y,x}$: $X_y$ is regular at $x$, hence smooth at $x$. (Here $f$ need not be flat.)

\textbf{(d)} Let $d=\dim\mathcal O_{X_y,x}=\dim B-\dim A$ (by (c) the fiber ring is regular, of dimension $\dim T_{X_y,x}$), $e=\dim A$. Choose $a_1,\dots,a_e$ generating $\fm_y$. Since $\mathcal O_{X_y,x}$ is regular of dimension $d$, $\fn$ is generated by $d$ elements $\bar t_1,\dots,\bar t_d$; lift them to $t_1,\dots,t_d\in\fm_x$. Then $\fm_x=\fm_yB+(t_1,\dots,t_d)=(f^\#a_1,\dots,f^\#a_e,t_1,\dots,t_d)$, with $e+d=\dim B$ generators; so $(f^\#a_i,t_j)$ is a coordinate system of $B$.

The local homomorphism $f^\#:A\to B$ induces $\widehat A\to\widehat B$ (continuity), and $\widehat B$ is complete, so there is a well defined homomorphism
\[ \varphi:\ P:=\widehat A[[T_1,\dots,T_d]]\longrightarrow\widehat B,\qquad T_j\mapsto t_j,\]
(the series $\sum a_\alpha t^\alpha$ converge since $t_j\in\fm_x\widehat B$). $P$ is a Noetherian complete local ring (Prop.~1.3.7) with maximal ideal $(\fm_y\widehat A,T_1,\dots,T_d)$, generated by $e+d$ elements.

\emph{$\varphi$ is surjective.} Since $k(x)=k(y)$ we have $B=f^\#(A)+\fm_x$, and since $\fm_x=(f^\#a_i,t_j)$ we have $\fm_x^n=\varphi_0((\fm_y,T)^n)+\fm_x^{n+1}$ for all $n$, where $\varphi_0:A[T_1,\dots,T_d]\to B$ is the polynomial version of $\varphi$. Hence $A[T]/(\fm_y,T)^n\to B/\fm_x^n$ is surjective for all $n$. Passing to the limit and using Lemma~1.3.3 exactly as in the proof of Prop.~2.27 (note $\varprojlim_nA[T]/(\fm_y,T)^n=\widehat A[[T]]$), the map $\varphi:\widehat A[[T]]\to\widehat B$ is surjective.

\emph{$\varphi$ is injective.} By Lemma~2.26, $\widehat B$ is regular of dimension $e+d$, hence a domain (Prop.~2.11); so $q=\Ker\varphi$ is a prime ideal of $P$ and $P/q\simeq\widehat B$ has dimension $e+d$. On the other hand $\dim P\le e+d$, because the maximal ideal of $P$ is generated by $e+d$ elements (Cor.~2.5.14(b)). As $\dim P\ge\dim P/q=e+d$, $P$ is a regular local ring of dimension $e+d$, hence a domain, and the equality $\dim P/q=\dim P$ forces $q=0$ (a nonzero prime in a Noetherian local domain strictly lowers the dimension). Thus $\varphi$ is an isomorphism:
\[ \widehat{\mathcal O}_{X,x}\simeq\widehat{\mathcal O}_{Y,y}[[T_1,\dots,T_d]]. \]

\emph{$f$ is smooth at $x$.} By Def.~3.35, we need (1) $f$ flat at $x$ and (2) $X_y$ smooth at $x$ over $k(y)$. (2) is (c). For (1): $\widehat A$ is flat over $A$ (Thm.~1.3.15), and $\widehat A[[T]]$ is flat over $\widehat A$ (it is the $(T)$-adic completion of the Noetherian ring $\widehat A[T]$, which is free over $\widehat A$, and completion is flat by Thm.~1.3.15). Hence $\widehat B\simeq\widehat A[[T]]$ is flat over $A$. Let $N'\to N$ be an injection of $A$-modules, $L$ the kernel of $B\otimes_AN'\to B\otimes_AN$. Tensoring with $\widehat B$ over $B$ (flat, Thm.~1.3.15), $\widehat B\otimes_BL$ is the kernel of $\widehat B\otimes_AN'\to\widehat B\otimes_AN$, which is $0$. Since $\widehat B$ is faithfully flat over $B$ (Thm.~1.3.16(b)), $L=0$. So $B$ is flat over $A$, i.e.\ $f$ is flat at $x$, and $f$ is smooth at $x$. \hfill$\square$

\exercise{3.26}{Let $X$ be a scheme of finite type over a locally Noetherian regular scheme $S$. Let $x\in X$.
(a) Show that in a neighborhood of $x$, there exists a closed immersion $i:X\to Z$ into a scheme $Z$ that is of finite type over $S$ and regular at $x$ (locally embed $X$ into an affine space over $S$ and use Theorem 3.36).
(b) Using exact sequence (2.1) in the proof of Proposition 2.5, show that one can successively reduce the dimension of $T_{Z,i(x)}$ in such a way that the canonical map $T_{X,x}\to T_{Z,i(x)}$ is an isomorphism.}
\pf
\textbf{(a)} The statement is local on $X$. Let $s\in S$ be the image of $x$, $V=\operatorname{Spec}R$ an affine open neighborhood of $s$ in $S$ ($R$ is Noetherian and regular), and $U=\operatorname{Spec}R[T_1,\dots,T_n]/I$ an affine open neighborhood of $x$ in $X$ lying over $V$ (possible since $X$ is of finite type over $S$). Let $Z=\mathbb A^n_V=\operatorname{Spec}R[T_1,\dots,T_n]$ and $i:U\to Z$ the closed immersion defined by $I$. The scheme $Z$ is of finite type over $V$, hence over $S$. The morphism $Z\to V$ is flat (the algebra $R[T]$ is free over $R$), of finite type, and its fibers $\mathbb A^n_{k(v)}$ are smooth over $k(v)$ (Def.~3.28 and the examples after it). So $Z\to V$ is smooth (Def.~3.35), and since $V$ is regular, $Z$ is regular by Thm.~3.36. In particular $Z$ is regular at $i(x)$.

\textbf{(b)} Let $z=i(x)$, $\fm_z$ the maximal ideal of $\mathcal O_{Z,z}$ and $J\subseteq\fm_z$ the ideal of $i$ at $z$, so $\mathcal O_{X,x}=\mathcal O_{Z,z}/J$ and $\fm_x=\fm_z/J$. As in the proof of Prop.~2.5 (exact sequence (2.1), which is valid in general), we have an exact sequence
\[ 0\to(J+\fm_z^2)/\fm_z^2\to\fm_z/\fm_z^2\to\fm_x/\fm_x^2\to0. \]
Dualizing, $T_{X,x}\to T_{Z,z}$ is always injective, and it is an isomorphism if and only if $J\subseteq\fm_z^2$.

Suppose $J\not\subseteq\fm_z^2$ and choose $g\in J\setminus\fm_z^2$ (so $g\in\fm_z$, $g\neq0$). Let $\mathcal I\subseteq\mathcal O_Z$ be the (quasi-coherent) ideal sheaf of the closed immersion $i$, so that $\mathcal I_z=J$. Choose an affine open neighborhood $W$ of $z$ and a section $\tilde g\in\mathcal I(W)$ whose germ at $z$ is $g$. Let $Z_1=V(\tilde g)\subseteq W$, a closed subscheme of $W$, of finite type over $S$. Since $\tilde g$ lies in the ideal of $i$, the closed immersion $i$ restricted to $i^{-1}(W)$ factors through a closed immersion $i_1:i^{-1}(W)\to Z_1$; we replace $X$ by the open neighborhood $i^{-1}(W)$ of $x$. We have $\mathcal O_{Z_1,z}=\mathcal O_{Z,z}/g\mathcal O_{Z,z}$. As $\mathcal O_{Z,z}$ is regular and $g\in\fm_z\setminus\fm_z^2$, $g\neq0$, Cor.~2.12 shows that $\mathcal O_{Z_1,z}$ is regular, of dimension $\dim\mathcal O_{Z,z}-1$. So $Z_1$ is of finite type over $S$ and regular at $z$. Moreover, with $\fm_1=\fm_z/(g)$, the exact sequence (2.1) for $(g)$ gives $\dim\fm_1/\fm_1^2=\dim\fm_z/\fm_z^2-\dim((g)+\fm_z^2)/\fm_z^2=\dim\fm_z/\fm_z^2-1$, i.e.\ $\dim T_{Z_1,z}=\dim T_{Z,z}-1$, since $g\notin\fm_z^2$.

We repeat this with $(Z_1,z)$ in place of $(Z,z)$. Each step lowers $\dim T_{Z,z}$ by exactly $1$, and this dimension is a nonnegative integer, so the process stops after at most $\dim T_{Z,z}$ steps. It stops only when $J\subseteq\fm_z^2$, i.e.\ when $T_{X,x}\to T_{Z,z}$ is an isomorphism. At the end we have a closed immersion of a neighborhood of $x$ in $X$ into a scheme $Z$ of finite type over $S$, regular at $x$, inducing an isomorphism on tangent spaces at $x$. \hfill$\square$

\section*{4.4 Zariski Main Theorem and applications}

Results of the text are cited by their numbers in Liu. Those of earlier chapters carry the chapter in front: ``Prop.~3.3.18'' is Proposition~18 of Section~3.3, ``Prop.~2.5.19'' is Proposition~19 of Section~2.5. Exercise numbers in the text of a solution (``Ex.~4.2'') refer to exercises of Section~4.4. An algebraic variety is a separated scheme of finite type over a field (not necessarily reduced), and ``quasi-finite'' means ``of finite type with finite fibres'' (Ex.~3.2.8).

Honest notes. Exercise~4.5 is false as printed (a cuspidal curve and a non-reduced scheme are counterexamples in characteristic $0$); it is proved under the additional hypotheses that $Y$ is normal and $X$ is reduced with all components dominating $Y$, and the solution says exactly what is and is not proved. In Exercise~4.4 the algebra $B$ must be nonzero. In Exercise~4.6 the induction does not close with the global fibre dimensions of the hint, so it is run with the local fibre dimension $\dim_x X_{f(x)}$ at a point; this is a slightly stronger statement and contains the one asked for. Exercises 4.1, 4.3, 4.4 and 4.6 use the finiteness theorem for projective morphisms (Thm.~5.3.2, in the form of Rem.~3.3.20: $\mathcal O_X(X)$ is finite over a Noetherian ring $A$ when $X$ is proper over $A$), through Thm.~4.4, Cor.~4.8 or Lemma~4.13; the book itself uses this theorem ``in anticipation'' at these places. Exercise~4.2 does not use it.

\medskip
\noindent\textbf{Fact (J).} \emph{Let $S$ be a scheme of finite type over a field $\kappa$. Every nonempty locally closed subset $L\subseteq S$ contains a point which is closed in $S$.} Indeed, $L$ with its reduced structure is a scheme of finite type over $\kappa$; let $\Spec C\subseteq L$ be a nonempty affine open subset of $L$ and $\fm$ a maximal ideal of $C$. By the Nullstellensatz $C/\fm$ is finite over $\kappa$, so the corresponding point $u$ has $k(u)$ algebraic over $\kappa$, hence is closed in $S$ (Ex.~2.5.9). Two consequences: a closed subset of $S$ containing all closed points of $S$ is $S$; and if the underlying set of $S$ is finite, then every point of $S$ is closed (a point of a finite $T_0$ space is locally closed, so Fact (J) applies to it).

\exercise{4.1}{Let $f:X\to Y$ be a projective morphism of algebraic varieties over an algebraically closed field $k$. Let us suppose that $f$ is injective and $T_{f,x}$ is injective for every closed point $x\in X$. Show that $f$ is a finite morphism, and, using Nakayama's lemma, show that $f$ is a closed immersion.}
\pf
Only the injectivity of $f$ on closed points is used. For a closed point $x$ of $X$ we have $k(x)=k$ (Ex.~2.5.9 and $k=\overline k$). Put $y=f(x)$; then $k\subseteq k(y)\subseteq k(x)=k$, so $k(y)=k$ and $y$ is closed (Ex.~2.5.9). The tangent map $T_{f,x}:T_{X,x}\to T_{Y,y}\otimes_{k(y)}k(x)=T_{Y,y}$ is the $k$-dual of the map $\fm_y/\fm_y^2\to\fm_x/\fm_x^2$ induced by $f^\#$; these are finite-dimensional vector spaces, so $T_{f,x}$ is injective if and only if $\fm_y/\fm_y^2\to\fm_x/\fm_x^2$ is surjective.

\emph{Step 1: the fibre over a closed point $y$ is empty or a single point.} The fibre $X_y=f^{-1}(y)$ is closed in $X$. If it is nonempty it contains a point closed in $X$ by Fact (J). Two distinct closed points of $X$ cannot both lie over $y$, since $f$ is injective on closed points; so $X_y$ contains exactly one closed point $x_0$ of $X$. If $X_y\neq\{x_0\}$, then $X_y\setminus\{x_0\}$ is locally closed in $X$ and nonempty, hence contains a point closed in $X$ by Fact (J), which is a second closed point over $y$, a contradiction. So $X_y=\{x_0\}$.

\emph{Step 2: $f$ is finite.} Let $y\in Y$ be a closed point. By Step 1, $X_y$ is finite, so by Ex.~4.2 below (whose proof does not use the present exercise) there is an affine open neighbourhood $V_y$ of $y$ such that $f^{-1}(V_y)\to V_y$ is finite. Let $W=\bigcup_yV_y$ over all closed points $y$. Then $Y\setminus W$ is closed and contains no closed point of $Y$, hence is empty by Fact (J). Thus the affine open sets $V_y$ cover $Y$ and $f^{-1}(V_y)$ is affine and finite over $V_y$; this implies that $f$ is finite (Ex.~3.3.15).

\emph{Step 3: $f$ is a closed immersion.} Let $V=\Spec A$ be an affine open subset of $Y$. Since $f$ is finite, $f^{-1}(V)=\Spec B$ with $B$ a finite $A$-module via $\varphi:A\to B$. We show that $\varphi$ is surjective. The $A$-module $M=B/\varphi(A)$ is finite. If $M\neq0$ there is a maximal ideal $\fm$ of $A$ with $M_{\fm}\neq0$. The corresponding point $y$ is closed in $V$, hence has $k(y)=A/\fm$ finite over $k$, so $k(y)=k$ and $y$ is closed in $Y$. We show $M_\fm=0$, a contradiction.

The ring $B_\fm=B\otimes_AA_\fm$ is a finite $A_\fm$-module with $B_\fm/\fm B_\fm=B/\fm B=\mathcal O(X_y)$. If $X_y=\emptyset$ then $B_\fm=0$ by Nakayama's lemma and $M_\fm=0$. Otherwise $X_y=\{x\}$ by Step 1. The maximal ideals of the finite $A_\fm$-algebra $B_\fm$ are the primes lying over $\fm A_\fm$, i.e. the points of $X_y$; so $B_\fm$ is local with maximal ideal corresponding to $x$, and $B_\fm=\mathcal O_{X,x}$. Since $T_{f,x}$ is injective, $\fm_x=\fm_y\mathcal O_{X,x}+\fm_x^2$. Let $Q=\fm_x/\fm_y\mathcal O_{X,x}$, a finitely generated module over the Noetherian local ring $\mathcal O_{X,x}$ (varieties are Noetherian). The equality gives $Q=\fm_xQ$, hence $Q=0$ by Nakayama, i.e. $\fm_x=\fm_y\mathcal O_{X,x}$. Therefore
\[ B_\fm/\fm B_\fm=\mathcal O_{X,x}/\fm_y\mathcal O_{X,x}=\mathcal O_{X,x}/\fm_x=k(x)=k, \]
so $A_\fm\to B_\fm/\fm B_\fm$ is surjective, i.e. $B_\fm=\varphi(A_\fm)+\fm B_\fm$. Thus $M_\fm=\fm M_\fm$, and $M_\fm=0$ by Nakayama's lemma applied to the finite $A_\fm$-module $M_\fm$.

Hence $A\to B$ is surjective for every affine open $V\subseteq Y$, which says that $f$ is a closed immersion. \hfill$\square$

\exercise{4.2}{Let $f:X\to Y$ be a projective morphism. Let $y\in Y$ be such that $X_y$ is finite. We want to show that there exists a neighborhood $V$ of $y$ such that $f^{-1}(V)\to V$ is finite, without using Zariski's ``Main Theorem''. We may suppose $Y$ is affine. (a) Let us fix an embedding $X\to\mathbb P^n_Y$. Show that there exists a non-empty principal closed subset $V_+(P)$ of $\mathbb P^n_Y$ such that $V_+(P)\cap X_y=\emptyset$ (use Prop.~3.3.36). (b) Show that there exists an affine neighborhood $V$ of $y$ such that $X\cap V_+(P)\cap\mathbb P^n_V=\emptyset$. (c) Show that $f^{-1}(V)=X\cap D_+(P)\cap\mathbb P^n_V$ and that $f^{-1}(V)$ is affine. Conclude using Prop.~3.3.18.}
\pf
Let $Y=\Spec A$, with $A$ any ring, and let $X\subseteq\mathbb P^n_Y=\operatorname{Proj}A[T_0,\dots,T_n]$ be a closed subscheme, $f$ being the restriction of the projection $\pi:\mathbb P^n_Y\to Y$. Let $\fp\subset A$ be the prime of $y$ and $\kappa=k(y)=\operatorname{Frac}(A/\fp)$. The fibre of $\mathbb P^n_Y$ over $y$ is $\mathbb P^n_\kappa$, and $F:=X_y$ is a finite subset of it.

\textbf{(a)} The points of $F$ correspond to finitely many homogeneous primes $\fq_1,\dots,\fq_r$ of $\kappa[T_0,\dots,T_n]$ not containing the irrelevant ideal $I_+=(T_0,\dots,T_n)$. By the argument of Prop.~3.3.36(a) with $Z=\emptyset$ (homogeneous prime avoidance, Ex.~2.5.5) there is a homogeneous element $g\in I_+$ lying in no $\fq_i$ (if $r=0$ take $g=T_0$). Then $e:=\deg g\ge1$ and $V_+(g)\cap F=\emptyset$. Multiplying $g$ by a nonzero element of $\kappa$ we may assume that its coefficients lie in the image of $A/\fp$, and we lift it to a homogeneous polynomial $P\in A[T_0,\dots,T_n]$ of degree $e$. The image of $P$ in $\kappa[T_0,\dots,T_n]$ is a nonzero multiple of $g$, so $V_+(P)\cap\pi^{-1}(y)=V_+(g)$ and therefore $V_+(P)\cap X_y=\emptyset$. (If $n\ge1$ the set $V_+(P)$ is nonempty because $\deg P\ge1$; nonemptiness plays no role below.)

\textbf{(b)} The set $C=X\cap V_+(P)$ is closed in $\mathbb P^n_Y$, because $X$ is closed. The morphism $\pi$ is proper (Thm.~3.3.30, in particular closed), so $\pi(C)$ is closed in $Y$. By (a), $y\notin\pi(C)$. Choose $a\in A$ with $y\in D(a)\subseteq Y\setminus\pi(C)$ and put $V=D(a)$, an affine neighbourhood of $y$. Then $C\cap\pi^{-1}(V)=\emptyset$, i.e. $X\cap V_+(P)\cap\mathbb P^n_V=\emptyset$.

\textbf{(c)} We have $f^{-1}(V)=X\cap\mathbb P^n_V$, and $\mathbb P^n_V=(\mathbb P^n_V\cap D_+(P))\cup(\mathbb P^n_V\cap V_+(P))$. By (b) the second part does not meet $X$, so $f^{-1}(V)=X\cap D_+(P)\cap\mathbb P^n_V$. Since $\deg P\ge1$, $D_+(P)=\Spec\bigl(A[T_0,\dots,T_n]_P\bigr)_0$ is affine, and $D_+(P)\cap\mathbb P^n_V$ is the principal open subset of $D_+(P)$ defined by the image of $a$, hence affine. The open subscheme $f^{-1}(V)$ of $X$ is a closed subscheme of $\mathbb P^n_V$ contained in the open subscheme $D_+(P)\cap\mathbb P^n_V$, hence a closed subscheme of this affine scheme; so $f^{-1}(V)=\Spec B'$ is affine.

The morphism $f^{-1}(V)\to V$ is proper (restriction of the proper morphism $f$ over the open $V$, Thm.~3.3.30 and Prop.~3.3.16), and of finite type (Cor.~3.3.10). Since $f^{-1}(V)$ is proper over $V=\Spec A_a$, the ring $B'=\mathcal O(f^{-1}(V))$ is integral over $A_a$ by Prop.~3.3.18. An algebra which is of finite type and integral is finite. Hence $f^{-1}(V)\to V$ is finite. \hfill$\square$

\exercise{4.3}{Let $Y$ be a locally Noetherian scheme, and let $f:X\to Y$ be a quasi-finite and quasi-projective morphism. Show that $f$ factors into an open immersion $X\to Z$ followed by a finite morphism $Z\to Y$.}
\pf
By Def.~3.3.35 there is a projective morphism $p:\overline X\to Y$ and an open immersion of finite type $X\to\overline X$ over $Y$; we regard $X$ as an open subscheme of $\overline X$. The following proposition does not use quasi-finiteness and is also used in Ex.~4.4.

\textbf{Proposition R.} \emph{There exist a finite $Y$-scheme $\pi:Z\to Y$, a $Y$-morphism $g:\overline X\to Z$ and an open subset $\Omega\subseteq\overline X$ such that $g|_\Omega:\Omega\to Z$ is an open immersion and $\Omega$ contains every point $x\in\overline X$ which is isolated in its fibre $\overline X_{p(x)}$. If $Y=\Spec A$ is affine (and Noetherian), $Z=\Spec\mathcal O_{\overline X}(\overline X)$ and $g$ is the canonical morphism.}

\emph{Proof.} For an affine open $W\subseteq Y$ let $\mathcal A(W)=\mathcal O_{\overline X}(p^{-1}W)$. It is a finite $\mathcal O_Y(W)$-algebra (Rem.~3.3.20 and Thm.~5.3.2: $p^{-1}W$ is proper over the Noetherian ring $\mathcal O_Y(W)$ by Thm.~3.3.30). For $a\in\mathcal O_Y(W)$ we have $\mathcal A(D(a))=\mathcal A(W)_a$ (flat base change, Prop.~3.1.24). Hence $\mathcal A=p_*\mathcal O_{\overline X}$ is a coherent sheaf of $\mathcal O_Y$-algebras, and we put $Z=\Spec_Y\mathcal A$, glued from the affine schemes $\Spec\mathcal A(W)$; it is finite over $Y$. The canonical morphisms $p^{-1}W\to\Spec\mathcal O(p^{-1}W)$ (Prop.~2.3.25) glue to $g:\overline X\to Z$. If $Y$ is affine we simply have $Z=\Spec\mathcal O_{\overline X}(\overline X)$.

Let $W\subseteq Y$ be an affine open subset; it is a Noetherian scheme. The morphism $p^{-1}W\to W$ is separated of finite type (Cor.~3.3.10) and $\mathcal O(p^{-1}W)$ is a finite $\mathcal O_Y(W)$-module. By Thm.~4.4, the set $\Omega_W$ of points of $p^{-1}W$ which are isolated in their fibres is open in $p^{-1}W$, and $g$ induces an open immersion $\Omega_W\to\Spec\mathcal A(W)=\pi^{-1}(W)$. Let $\Omega=\bigcup_W\Omega_W$; since fibres of $p^{-1}W$ and of $\overline X$ coincide, $\Omega\cap p^{-1}W=\Omega_W$, and $\Omega$ is exactly the set of points isolated in their fibres. Being an open immersion is local on the target $Z=\bigcup\pi^{-1}(W)$, so $g|_\Omega$ is an open immersion. \hfill$\square$

\emph{Quasi-finiteness gives $X\subseteq\Omega$.} Let $x\in X$ and $y=f(x)$. The fibre $X_y$ is a scheme of finite type over $k(y)$ (as $f$ is of finite type) with finitely many points, so every point of $X_y$ is closed in $X_y$ (Fact (J)), i.e. $X_y$ is a finite discrete space and $\{x\}$ is open in $X_y$. Now $X_y=X\cap\overline X_y$ is open in $\overline X_y$ because $X$ is open in $\overline X$. So $\{x\}$ is open in $\overline X_y$, i.e. $x$ is isolated in its fibre in $\overline X$. Hence $X\subseteq\Omega$.

Consequently $g|_X:X\to Z$ is an open immersion (the restriction of the open immersion $g|_\Omega$ to the open subscheme $X\subseteq\Omega$), and $\pi:Z\to Y$ is finite. Since $\pi\circ g=p$ and $X\to\overline X$ is over $Y$, the morphism $f$ is the composition $X\to Z\to Y$. \hfill$\square$

\exercise{4.4}{Let $A$ be a discrete valuation ring, and let $B$ be a finitely generated flat $A$-algebra. Show that there exists a subalgebra $C$ of $B$ with a finite injective homomorphism $A[T_1,\dots,T_d]\to C$, and such that the morphism $\Spec B\to\Spec C$ is an open immersion. This is a sort of generalization of Noether's normalization lemma. Show that in general we cannot have $C=B$.}
\pf
We assume $B\neq0$ (for $B=0$ there is no injection $A[T_1,\dots,T_d]\to C\subseteq B$; the zero ring is flat and finitely generated, so the statement needs $B\neq0$). Let $t$ be a uniformizer, $K=\operatorname{Frac}(A)=A[1/t]$, $k=A/tA$, and $B_K=B[1/t]=B\otimes_AK$. Since $B$ is flat, multiplication by $t$ is injective on $B$, so $B$ is torsion free and $B\to B_K$ is injective; in particular $B_K\neq0$.

\textbf{Lemma A.} \emph{There exist $s\ge0$ and $z_1,\dots,z_s\in B$ such that (i) the $K$-algebra map $K[T_1,\dots,T_s]\to B_K$, $T_i\mapsto z_i$, is injective and makes $B_K$ a finite module; (ii) $B/tB$ is a finite module over $k[T_1,\dots,T_s]$ via $T_i\mapsto\overline{z_i}$.}

\emph{Proof.} Induction on the number $m$ of generators of the $A$-algebra $B=A[x_1,\dots,x_m]$ (all such $B\neq0$ are flat, hence torsion free). Let $J$ be the kernel of $A[X_1,\dots,X_m]\to B$. If $J=0$, then $B$ is a polynomial ring and $z_i=x_i$ ($s=m$) work. In particular this covers $m=0$: then $B=A/J$ is torsion free and nonzero, so $J=0$. Suppose $J\neq0$, so $m\ge1$. Take $0\neq g\in J$ and write $g=t^eg_1$ with $g_1\in A[X]$ having a coefficient which is a unit. Since $t$ is a nonzerodivisor on $B$, $g_1\in J$, and the reduction $\overline{g_1}\in k[X]$ is nonzero.

Choose an integer $r$ larger than every exponent of every variable occurring in $g_1$, and put $y_i=x_i-x_m^{r^i}$ for $1\le i\le m-1$. Let $B'=A[y_1,\dots,y_{m-1}]\subseteq B$, a subalgebra with $m-1$ generators, torsion free hence flat and nonzero. Consider
\[ G(Y_1,\dots,Y_{m-1},X)=g_1(Y_1+X^{r},\dots,Y_{m-1}+X^{r^{m-1}},X)\in A[Y][X]. \]
For a monomial $X^\alpha$ the substitution gives $X^{w(\alpha)}+{}$(terms of lower degree in $X$ with coefficients in $\mathbf Z[Y]$), where $w(\alpha)=\alpha_m+\alpha_1r+\dots+\alpha_{m-1}r^{m-1}$; the numbers $w(\alpha)$ are pairwise distinct for exponent vectors with entries $<r$ (base-$r$ digits). Hence $G=cX^N+{}$lower terms in $X$ with coefficients in $A[Y]$, where $c\in A$ is the coefficient in $g_1$ of the monomial of maximal weight, $c\neq0$. Likewise, the reduction $\overline G=\overline{g_1}(Y_i+X^{r^i},X)\in k[Y][X]$ equals $\bar c'X^{N'}+{}$lower terms, with $\bar c'\in k^\times$ the coefficient of the monomial of maximal weight in $\overline{g_1}$.

Now $G(y,x_m)=g_1(x)=0$ in $B$, so in $B_K$ the element $x_m$ satisfies an equation with leading coefficient $c\in K^\times$ over $K[y]$: it is integral over $B'_K=B'\otimes K$. As $x_i=y_i+x_m^{r^i}$, $B_K=B'_K[x_m]$ is a finite $B'_K$-module. In $B/tB$ we have $\overline G(\bar y,\bar x_m)=0$ with leading coefficient $\bar c'\neq0$, so $\bar x_m$ is integral over the subalgebra $\overline{B'}$ of $B/tB$ generated by the $\bar y_i$, and $B/tB=\overline{B'}[\bar x_m]$ is a finite $\overline{B'}$-module.

By induction applied to $B'$ there are $z_1,\dots,z_s\in B'\subseteq B$ with $K[T]\to B'_K$ injective and $B'_K$ finite over $K[T]$, and $B'/tB'$ finite over $k[T]$. Since $B'_K\subseteq B_K$, (i) holds for $B$. The ring $\overline{B'}$ is a quotient of $B'/tB'$, so it is finite over the subalgebra generated by the $\bar z_i$; as $B/tB$ is finite over $\overline{B'}$, (ii) holds. \hfill$\square$

\textbf{The map to affine space.} Take $z_1,\dots,z_s$ as in Lemma A, put $d=s$, $R=A[T_1,\dots,T_d]$ and let $\phi:R\to B$ send $T_i$ to $z_i$. It is injective: if $\phi(h)=0$, then $h=0$ in $K[T]$ by Lemma A(i) and $B\subseteq B_K$, so $h=0$. The morphism $\Spec B\to\Spec R$ is quasi-finite: for $q\in\Spec R$ with residue field $\kappa(q)$, if $t\in q$ then $B\otimes_R\kappa(q)=(B/tB)\otimes_{k[T]}\kappa(q)$ is finite-dimensional by Lemma A(ii); if $t\notin q$ then $B\otimes_R\kappa(q)=B_K\otimes_{K[T]}\kappa(q)$ is finite-dimensional by Lemma A(i). Thus all fibres are finite. Also $R$ is Noetherian, and $\Spec B\to\Spec R$ is of finite type with affine source and target, hence quasi-projective (Ex.~3.3.20(b)).

\textbf{The algebra $C$.} By Ex.~4.3 (more precisely, by Proposition~R and the argument there), there is a projective $R$-scheme $\overline X$ containing $X=\Spec B$ as an open subscheme, and, with $D=\mathcal O_{\overline X}(\overline X)$ (a finite $R$-algebra), the canonical morphism $g:X\to\Spec D$ is an open immersion. It corresponds to the restriction homomorphism $\psi:D\to\mathcal O_X(X)=B$. Let $C=\psi(D)\subseteq B$. It is an $R$-subalgebra of $B$ containing $\phi(R)$ and is finite over $R$ as a quotient of $D$; and $R\to C$ is injective because $R\to C\subseteq B$ is $\phi$. So $R=A[T_1,\dots,T_d]\to C$ is finite and injective.

It remains to see that $q:\Spec B\to\Spec C$ is an open immersion. The morphism $g$ factors as $X\xrightarrow{q}\Spec C\xrightarrow{\iota}\Spec D$ with $\iota$ a closed immersion. Since $g$ is an open immersion, $g(X)$ is open in $\Spec D$ and contained in the closed subset $\Spec C$, so $q$ is a homeomorphism onto the open subset $g(X)$ of $\Spec C$. For $x\in X$ the stalk map $\mathcal O_{\Spec D,g(x)}\to\mathcal O_{X,x}$ is the composite of the surjection $\mathcal O_{\Spec D,g(x)}\to\mathcal O_{\Spec C,q(x)}$ and of $q^\#_x$; as it is an isomorphism, the surjection is also injective, hence bijective, and so $q^\#_x$ is an isomorphism. Therefore $q$ is an open immersion.

\textbf{We cannot take $C=B$ in general.} Let $B=A[X]/(tX-1)\simeq K$. It is finitely generated over $A$, and torsion free, hence flat. Suppose $A[T_1,\dots,T_d]\to B=K$ were finite and injective. Every element of $K$ is algebraic over $A$ (it is a quotient $a/t^n$, a root of $t^nT-a$), so a homomorphism $A[T_1,\dots,T_d]\to K$ with $d\ge1$ has a nonzero kernel; hence $d=0$. But $A\to K$ is not finite: if $K=a_1A+\dots+a_rA$, write $a_i=u_i/t^{n}$ with $u_i\in A$; then $K\subseteq t^{-n}A$, which is false since $t^{-n-1}\notin t^{-n}A$. So $C=B$ is impossible. For this $B$ the proof gives $d=0$ and $C=A$, and indeed $\Spec K=D(t)\subseteq\Spec A$ is open. \hfill$\square$

\exercise{4.5}{Let $Y$ be a locally Noetherian integral scheme, and let $f:X\to Y$ be a projective morphism with geometrically connected fibres. Show that for any affine open subset $V$ of $Y$, there exists an $n\ge1$ such that $(f_*\mathcal O_X(V))^{p^n}$ is contained in the image of $\mathcal O_Y(V)\to f_*\mathcal O_X(V)$, where $p$ is the characteristic of $K(Y)$ (by convention, we set $p^n=1$ if $p=0$). In particular, if $\operatorname{char}K(Y)=0$, then $\mathcal O_Y\to f_*\mathcal O_X$ is an isomorphism.}
\pf
\textbf{The statement as printed is false}, so we first give counterexamples and then prove a corrected version. Here ``geometrically connected fibres'' means that every fibre is nonempty and stays connected after extension of the base field to an algebraic closure.

\emph{Counterexample 1 ($Y$ not normal).} Let $k$ be a field of characteristic $0$, $A=k[t^2,t^3]\subset B=k[t]$, $Y=\Spec A$ (integral, Noetherian) and $X=\Spec B$. The morphism $f:X\to Y$ is finite, hence projective (Ex.~3.3.22). Its fibre over the singular point is $\Spec k[t]/(t^2)$, a single point, and the other fibres are single reduced points; so all fibres are geometrically connected. But $f_*\mathcal O_X(Y)=B\neq A$ and $p^n=1$, so $B^{p^n}=B\not\subseteq A$, and $\mathcal O_Y\to f_*\mathcal O_X$ is not an isomorphism, although $X$ is integral.

\emph{Counterexample 2 ($X$ not reduced).} Let $k$ be a field of characteristic $0$, $Y=\mathbb A^1_k$ (normal) and $X=Y\times_k\Spec k[\varepsilon]$, $\varepsilon^2=0$. Then $X\to Y$ is finite (free of rank~$2$), hence projective, with fibres $\Spec\kappa[\varepsilon]$, geometrically connected. But $f_*\mathcal O_X(Y)=k[T][\varepsilon]\neq k[T]$.

So we prove the following.

\textbf{Theorem.} \emph{Let $Y$ be a normal locally Noetherian integral scheme with generic point $\eta$, $K=K(Y)$, $p=\operatorname{char}K$. Let $f:X\to Y$ be a projective morphism with $X$ reduced, every irreducible component of $X$ dominating $Y$ (e.g. $X$ integral and $f$ dominant), and with $X_\eta$ nonempty and geometrically connected over $K$ (this follows if all fibres are geometrically connected, since then $f(X)$, which is closed by Thm.~3.3.30, contains all points, i.e. $f$ is surjective). Then for every affine open $V$ of $Y$ there is $n\ge1$ with $(f_*\mathcal O_X(V))^{p^n}\subseteq\operatorname{Im}(\mathcal O_Y(V)\to f_*\mathcal O_X(V))$ (with $p^n=1$ if $p=0$). If $p=0$, then $\mathcal O_Y\to f_*\mathcal O_X$ is an isomorphism.}

\emph{Proof.} Let $V=\Spec A$, $A$ an integrally closed Noetherian domain with fraction field $K$, $B=\mathcal O_X(f^{-1}V)=f_*\mathcal O_X(V)$ and $L=\mathcal O(X_\eta)$.

\emph{(1) $L=B\otimes_AK$ and $B\to L$ is injective.} Since $A\to K$ is flat and $f^{-1}(V)\to V$ is quasi-compact and separated, $\mathcal O(X_\eta)=B\otimes_AK$ (Prop.~3.1.24). Let $b\in B$ with $ab=0$, $0\neq a\in A$. At the generic point $\xi$ of an irreducible component of $X$, $f(\xi)=\eta$ by hypothesis, and $\mathcal O_{X,\xi}$ is a field (as $X$ is reduced) which is a $K=\mathcal O_{Y,\eta}$-algebra; so $a$ is invertible there and $b_\xi=0$. Since $X$ is reduced and locally Noetherian, a section vanishing at all generic points vanishes (an affine open $\Spec R$ with $R$ reduced has $R\hookrightarrow\prod R_{\fp_i}$ over its minimal primes). So $b=0$, i.e. $B$ is torsion free and $B\hookrightarrow B\otimes K=L$.

\emph{(2) $L$ is a finite purely inseparable extension of $K$.} The $K$-scheme $X_\eta$ is projective, hence proper (Thm.~3.3.30), reduced (it is obtained from the reduced affine pieces of $X$ by inverting $A\setminus\{0\}$), nonempty, and geometrically connected, hence connected. By Cor.~3.3.21, $L=\mathcal O(X_\eta)$ is a field, finite and purely inseparable over $K$. If $p=0$ this means $L=K$. If $p>0$, write $L=K(\ell_1,\dots,\ell_r)$ with $\ell_i^{p^{n_i}}\in K$ and let $n\ge\max n_i$, $n\ge1$. Then $\ell^{p^n}\in K$ for every $\ell\in L$, since $\ell$ is a polynomial in the $\ell_i$ with coefficients in $K$ and $x\mapsto x^{p^n}$ is additive and multiplicative.

\emph{(3) Conclusion.} Let $b\in B\subseteq L$. By Prop.~3.3.18, $b$ is integral over $A$ ($f^{-1}(V)$ is proper over $V$). Then $b^{p^n}$ is integral over $A$ and lies in $K$ (in $K$ itself, if $p=0$, with $p^n=1$), hence lies in $A$ because $A$ is integrally closed. The element $b^{p^n}\in B\subseteq L$ is therefore the image of an element of $A$ (the map $A\to B\to L$ is the inclusion $A\subseteq K\subseteq L$), i.e. $B^{p^n}\subseteq\operatorname{Im}(A\to B)$.

If $p=0$ this gives $B\subseteq\operatorname{Im}(A)$, so $A\to B$ is surjective. It is injective as well, because $A\to B\to L$ is the composite $A\subseteq K\to L$, injective since $K$ is a field and $L\neq0$. As this holds for every affine open $V$, $\mathcal O_Y\to f_*\mathcal O_X$ is an isomorphism. \hfill$\square$

\emph{What is not proved.} If $Y$ is not normal, the same argument only gives $b^{p^n}\in K\cap B$, i.e. $b^{p^n}$ lies in the integral closure of $A$ in $K$; Counterexample~1 shows that in characteristic $0$ one cannot do better. For $Y$ not normal in characteristic $p>0$, or for $X$ non-reduced in characteristic $p>0$ (where torsion in $B$ has to be controlled using the special fibres), I do not prove the printed statement. Only the geometric connectedness of the generic fibre was used. \hfill$\square$

\exercise{4.6}{Let $Y$ be a locally Noetherian irreducible scheme, and let $f:X\to Y$ be a dominant morphism of finite type with $X$ irreducible. We are going to generalize Lemma~4.13 and show that $\dim X_\eta\le\dim X_y$ if $X_y\neq\emptyset$. By this lemma, we may suppose that $\dim X_y\ge1$. We will use induction on $\dim X_y+\dim\mathcal O_{Y,y}$. (a) Let $x\in X_y$ be such that $\dim\mathcal O_{X_y,x}=\dim X_y$. Show that there exists a $b\in\fm_x\mathcal O_{X,x}$ which belongs to none of the minimal prime ideals containing $\fm_y\mathcal O_{X,x}$. Restricting $X$ if necessary, we may suppose $X$ is affine and $b\in\mathcal O_X(X)$. (b) Let $b$ be as above. Show that there exists an irreducible component $Z$ of $V(b)$ passing through $x$ and not contained in $X_y$. Let $z'$ be the generic point of $Z$, and $y'=f(z')$ the generic point of $f(Z)$ (Ex.~2.4.11). Show that $\dim\mathcal O_{Y,y'}<\dim\mathcal O_{Y,y}$ and $\dim X_{y'}-\dim X_y\le\dim Z_{y'}-\dim Z_y$. By applying the induction hypothesis to $Z\to f(Z)$ and $X\times_Y\Spec\mathcal O_{Y,y'}\to\Spec\mathcal O_{Y,y'}$, show that $\dim X_\eta\le\dim X_y$.}
\pf
\textbf{Notation and the form in which the statement is proved.} Let $\eta$ be the generic point of $Y$ and $\xi$ that of $X$; since $f$ is dominant, $f(\xi)=\eta$. For $y\in Y$ the fibre $X_y$ is a scheme of finite type over $k(y)$ whose underlying space is $f^{-1}(y)$. For a scheme $S$ of finite type over a field and $s\in S$ put $\dim_sS=\max\{\dim T: T$ an irreducible component of $S$ containing $s\}$, and for $x\in X$ put $e(x)=\dim_xX_{f(x)}$. The inequality of (b) is a comparison of fibre dimensions near $x$ and $z'$ rather than of whole fibres (the whole fibre $X_{y'}$ may have components far from $z'$, and controlling them would need an upper semicontinuity theorem which is not available at this point). So we prove the local statement

\textbf{Theorem.} \emph{Let $Y$ be a locally Noetherian irreducible scheme, $f:X\to Y$ a dominant morphism of finite type, $X$ irreducible, $x\in X$. Then $\dim X_\eta\le e(x)$.}

This contains the exercise: if $X_y\neq\emptyset$, take $x$ the generic point of an irreducible component of $X_y$ of maximal dimension; the only component through $x$ is that one (a component through $x$ contains $\overline{\{x\}}$, hence equals it), so $e(x)=\dim X_y$.

\textbf{Facts used.} (D1) For an irreducible scheme $S$ of finite type over a field, a nonempty open subset has the same dimension (Prop.~2.5.19), and an irreducible closed subset $T\subsetneq S$ has $\dim T<\dim S$. (D2) If $S$ is an irreducible variety and $b\in\mathcal O(S)$ is non-nilpotent, every irreducible component of $V(b)$ has dimension $\dim S-1$ (Cor.~2.5.26). (D3) Krull's principal ideal theorem and prime avoidance. (D4) Lemma~4.13.

Moreover, $X_\eta$ is irreducible: it contains $\xi$, and a subset of an irreducible space containing the generic point is irreducible. By (D1), for every nonempty open $W\subseteq X$ (which contains $\xi$), $W_\eta=W\cap X_\eta$ is a nonempty open subset of $X_\eta$, so $\dim W_\eta=\dim X_\eta$. (5)

\textbf{Proof of the Theorem.} We use strong induction on $N=e(x)+\dim\mathcal O_{Y,y}$, $y=f(x)$, simultaneously for all data $(Y,X,f,x)$ as in the Theorem. Put $d=e(x)$.

\emph{Case $y=\eta$.} Then $X_y=X_\eta$ is irreducible, so $e(x)=\dim X_\eta$.

\emph{Case $y\neq\eta$ and $d=0$} (the case ``$\dim X_y=0$'' treated by Lemma~4.13). Then $\{x\}$ is a component of $X_y$ and no other component contains $x$, so $\{x\}$ is open in $X_y$. Hence there is an open $W\subseteq X$ with $W\cap X_y=\{x\}$. Then $W$ is irreducible, $W\to Y$ is of finite type, and $W_y=\{x\}$ is nonempty and finite. By Lemma~4.13, $W_\eta$ has at most one point; it contains $\xi$, so $W_\eta=\{\xi\}$ and $\dim X_\eta=\dim W_\eta=0$ by (5).

\emph{Case $y\neq\eta$ and $d\ge1$.} Note $\dim\mathcal O_{Y,y}\ge1$.

\emph{Reduction to a good point.} Let $C_1,\dots,C_r$ be the irreducible components of $X_y$ containing $x$, and $F$ the union of the others. Then $F=X_y\cap F''$ for a closed subset $F''$ of $X$, and $x\notin F''$. Replacing $X$ by $X\setminus F''$ we may assume that all irreducible components of $X_y$ pass through $x$; by (D1) their dimensions are unchanged, and (5) holds. If $x$ is the generic point of a component of $X_y$, then this component is the only one; we then replace $x$ by a closed point of $X_y$ (it still lies on the unique component, so $e(x)=d$ is unchanged). In every case we may therefore assume: (G1) every irreducible component of $X_y$ contains $x$; (G2) $x$ is not the generic point of a component of $X_y$ (if $x$ were the generic point of one component, all others would contain $x$, hence equal it); (G3) $e(x)=d=\dim X_y\ge1$. Shrinking $X$ further to an affine open neighbourhood $W=\Spec B_0$ of $x$ with $f(W)$ inside an affine open $\Spec A_0\subseteq Y$ preserves (G1)--(G3) and (5). So assume $X=\Spec B_0$ is affine.

\emph{(a) The element $b$.} Let $R=\mathcal O_{X,x}$ and $\bar R=R/\fm_y R=\mathcal O_{X_y,x}$. The minimal primes $\bar q_1,\dots,\bar q_r$ of $\bar R$ correspond to the components of $X_y$ through $x$. By (G2), $\bar q_j\neq\bar\fm$ (the maximal ideal of $\bar R$), i.e. $\bar\fm\not\subseteq\bar q_j$ for all $j$. By prime avoidance there is $\bar b\in\bar\fm\setminus\bigcup_j\bar q_j$; lift it to $b\in\fm_xR$. Then $b$ lies in no minimal prime $q_j$ of $R$ containing $\fm_yR$. Shrinking $X$ to a smaller affine neighbourhood of $x$ (this preserves (G1)--(G3) and (5)) we may assume $b\in\mathcal O_X(X)$.

For a component $C$ of $X_y$ the generic point $c$ of $C$ is a generization of $x$ and $b\notin\bar q_C$, so $b$ is a unit of $\mathcal O_{X_y,c}=\bar R_{\bar q_C}$; thus $b$ does not vanish at $c$, and $V(b)\cap C$ is a proper closed subset of $C$. By (D1) every irreducible closed subset of $V(b)\cap X_y$ has dimension $\le d-1$. (6)

\emph{(b) The component $Z$.} Irreducible components of $V(b)$ through $x$ correspond to the minimal primes $P$ of $bR$; each has $\operatorname{ht}P\le1$ by Krull. Suppose every such $P$ contained $\fm_yR$. Take one, $P$. It contains a minimal prime $q$ of $\fm_yR$, and $b\in P\setminus q$, so $q\subsetneq P$. The ring $R$ has a unique minimal prime $p_0$ (that of $\xi$: every generization of $x$ is a specialization of $\xi$), and $p_0\subseteq q\subsetneq P$ with $\operatorname{ht}P\le1$ forces $q=p_0$. Then $\fm_yR\subseteq p_0$ means $f(\xi)=y$, i.e. $y=\eta$, contrary to assumption. So some irreducible component $Z$ of $V(b)$ through $x$ has generic point $z'$ with $f(z')\neq y$, i.e. $Z\not\subseteq X_y$. Put $y'=f(z')$. We give $Z$ the reduced structure and $Y'=\overline{\{y'\}}$ the reduced structure; then $f$ induces a dominant morphism of finite type $Z\to Y'$ (as $Z$ is reduced), with $Z$ irreducible, $Y'$ locally Noetherian irreducible with generic point $y'$, and $x\in Z$, $f(x)=y\in Y'$ (as $f(x)\in f(\overline{\{z'\}})\subseteq\overline{\{y'\}}$).

\emph{(b1) $\dim\mathcal O_{Y,y'}<\dim\mathcal O_{Y,y}$.} Indeed $y'\neq y$ is a generization of $y$, so $\mathcal O_{Y,y'}$ is the localization of $\mathcal O_{Y,y}$ at a prime $\fp\subsetneq\fm_y$, and $\dim\mathcal O_{Y,y'}=\operatorname{ht}\fp<\operatorname{ht}\fm_y=\dim\mathcal O_{Y,y}$.

\emph{(b2) $e_Z(x):=\dim_xZ_y\le d-1$,} where $Z_y$ is the fibre of $Z\to Y'$ over $y$. Indeed $Z_y\subseteq V(b)\cap X_y$ is closed in $X_y$ and (6) applies to each of its components. Also $\dim\mathcal O_{Y',y}\le\dim\mathcal O_{Y,y}$, since $\mathcal O_{Y',y}$ is a quotient of $\mathcal O_{Y,y}$.

\emph{(b3) $\dim_{z'}X_{y'}\le\dim Z_{y'}+1$,} where $Z_{y'}$ is the generic fibre of $Z\to Y'$. Let $C'$ be an irreducible component of $X_{y'}$ containing $z'$ (a $k(y')$-scheme of finite type; give $C'$ the reduced structure). The set $Z_{y'}=Z\cap X_{y'}$ is the closure of $\{z'\}$ in $X_{y'}$, so $Z_{y'}\subseteq C'$. Consider two cases. If $C'\subseteq V(b)$: $C'$ is irreducible, contains $z'$, and lies in $V(b)$, so it lies in an irreducible component of $V(b)$ containing $z'$, which must be $Z$ (any such component contains $\overline{\{z'\}}=Z$ and $Z$ is a component); then $C'\subseteq Z\cap X_{y'}=Z_{y'}$ and $C'=Z_{y'}$. If $C'\not\subseteq V(b)$: $b$ is non-nilpotent on the integral variety $C'$, and $Z_{y'}$ is an irreducible component of $V(b)\cap C'$: if $T\supseteq Z_{y'}$ is an irreducible closed subset of $V(b)\cap C'$, its closure $\overline T$ in $X$ is irreducible, contained in $V(b)$ and contains $z'$, hence $\overline T\subseteq Z$, so $T\subseteq Z\cap X_{y'}=Z_{y'}$. By (D2), $\dim Z_{y'}=\dim C'-1$. In both cases $\dim C'\le\dim Z_{y'}+1$.

By (b2) and (b3), $\dim_{z'}X_{y'}-\dim_xX_y\le\dim Z_{y'}+1-d\le\dim Z_{y'}-\dim_xZ_y$; this is the local form of the inequality asked for in (b).

\emph{(c) Induction.} First apply the induction hypothesis to $Z\to Y'$ and the point $x$: its invariant is $e_Z(x)+\dim\mathcal O_{Y',y}\le d-1+\dim\mathcal O_{Y,y}<N$. It gives $\dim Z_{y'}\le e_Z(x)\le d-1$, because $y'$ is the generic point of $Y'$. By (b3), $e(z')=\dim_{z'}X_{y'}\le d$.

Next consider $X''=X\times_Y\Spec\mathcal O_{Y,y'}\to Y''=\Spec\mathcal O_{Y,y'}$. Here $Y''$ is a Noetherian irreducible scheme with generic point $\eta$ and closed point $y'$. The morphism $X''\to X$ is a homeomorphism onto the set of points of $X$ lying over generizations of $y'$ (localization on the affine scheme $X=\Spec B_0$ over $\Spec A_0$); this set contains $\xi$, so $X''$ is irreducible, $X''\to Y''$ is dominant and of finite type, $X''_{y'}=X_{y'}$ and $X''_\eta=X_\eta$. The point $z'\in X''$ lies over $y'$, and its invariant is $e(z')+\dim\mathcal O_{Y,y'}\le d+\dim\mathcal O_{Y,y}-1<N$ by (b1). The induction hypothesis gives $\dim X_\eta=\dim X''_\eta\le e(z')\le d=e(x)$. \hfill$\square$

\emph{What was proved.} The statement of the exercise follows from the Theorem as explained above. Items (a) and (b) are proved as stated, except that the inequality $\dim X_{y'}-\dim X_y\le\dim Z_{y'}-\dim Z_y$ is proved in the local form $\dim_{z'}X_{y'}-\dim_xX_y\le\dim Z_{y'}-\dim_xZ_y$ (all that the induction needs), and $x$ is chosen as in the reduction step. The proof relies on Lemma~4.13, hence on Thm.~5.3.2. \hfill$\square$

\end{document}
